\chapter[Affine VOAs for every simple Lie algebra; the exceptional Lie algebras]{Affine vertex algebras for every simple Lie algebra, and the exceptional Lie algebras}\label{ch:L23}
\begingroup\setlength{\emergencystretch}{3em}\sloppy

\providecommand{\lxxiiig}{\mathfrak g}
\providecommand{\lxxiiivac}{|0\rangle}
\providecommand{\lxxiiiV}{V^k(\mathfrak g)}
\providecommand{\lxxiiiL}{L_k(\mathfrak g)}
\providecommand{\lxxiiihv}{h^\vee}
\providecommand{\lxxiiiad}{\operatorname{ad}}
\providecommand{\lxxiiiKil}{\operatorname{Kil}}
\providecommand{\lxxiiiCas}{C}

This chapter was written against an earlier state of the LeanVOA library and updated on 28 October 2026 (the named simple
quotient of Section~\ref{l23:sec:simplename}, the quoted paths, and the renaming of the universal affine VOA). Every Lean
quotation in a \texttt{leancode} box was copied from the source files; for this edition, \texttt{scripts/textbook/sync.py} checks
that every quotation occurs verbatim in the source of the release. Checks that were run at an earlier state of the library and
not re-run for this edition are marked as such. How to read the quotations (what they omit, and the Lean
constructs that occur in them) is explained in Section~\ref{l23:sec:readlean}; what the machine confirmed for this
chapter, a map of the files it covers, and the cost of the kernel computations are collected in one place, the
verification record of Section~\ref{l23:sec:verification}.

\paragraph{What the chapter is about.}
Chapter~\ref{ch:L08} built the universal affine vertex operator algebra $V^k(\mathfrak{sl}_2)$, and the bespoke files
\leanpath{VOA/\allowbreak Examples/\allowbreak AffineSl3*.lean} (whose chapter is in preparation) construct $V^k(\mathfrak{sl}_3)$ as \lean{VOA.AffineSl3.affineSl3VOA}; both are
long chains of files written against an explicit basis ($\texttt{Fin 3}$, resp.\ $\texttt{Fin 8}$). This chapter does the construction \emph{once}, for an
arbitrary finite-dimensional complex Lie algebra $\lxxiiig$ equipped with a small package of linear algebra (an
invariant form, a basis, a dual basis and the eigenvalue of the adjoint Casimir), shows that the package exists for
every simple $\lxxiiig$, and then constructs \emph{inside Lean} concrete simple Lie algebras to which the theory can be
applied.

Why the library needs this. The later Chapters~\ref{ch:L24}--\ref{ch:L27} study the simple affine vertex algebras
$L_k(\lxxiiig)$ for the seven Lie algebras of the \emph{Deligne exceptional series}
$A_2\subset G_2\subset D_4\subset F_4\subset E_6\subset E_7\subset E_8$, at the negative levels $k=-\lxxiiihv/6-1$
($\lxxiiihv$ the dual Coxeter number). These levels are singled out by a theorem of Arakawa and Moreau: at them
the associated variety of $L_k(\lxxiiig)$ (the variety cut out by the relations among the currents at leading
order, Section~\ref{l23:sec:deligne}) is the closure of the minimal nilpotent orbit of $\lxxiiig$, and the vacuum
module acquires null vectors at conformal weight two. In physics, by the work of Beem, Lemos, Liendo, Peelaers,
Rastelli and van Rees, $L_k(\lxxiiig)$ at these levels is the chiral algebra of a rank-one four-dimensional
$\mathcal N=2$ superconformal field theory with flavour symmetry $\lxxiiig$ (for $A_2,D_4,E_6,E_7,E_8$; no
four-dimensional theory is known for $G_2,F_4$). To state anything about these algebras in Lean one needs two things
that Mathlib does not provide: an affine vertex algebra for an arbitrary simple Lie algebra, and the seven simple Lie
algebras themselves, as concrete Lean objects proved simple. In addition, the integer certificates used in
Chapter~\ref{ch:L24} were computed outside Lean, by Python, in a fixed basis; to use them as checked inputs, Lean
needs Lie data whose structure constants are \emph{literally} those certificate tables. This chapter supplies all
three.

\paragraph{The logical arc.} The mathematical backbone is a chain of five steps:
\begin{equation*}
\underbrace{\text{Lie data}}_{\S\ref{l23:sec:liedata}}
\ \xrightarrow{\ \S\ref{l23:sec:universal}\ }\
\underbrace{V^k(\lxxiiig)}_{\text{a VOA, }k+h\neq0}
\ \xrightarrow{\ \S\ref{l23:sec:simple}\ }\
\underbrace{L_k(\lxxiiig)}_{\text{simple quotient}}
\ \xrightarrow{\ \S\ref{l23:sec:deligne}\ }\
\underbrace{\text{inputs of the capstone of Ch.~\ref{ch:L24}}}_{\text{jet surjectivity; classical freeness if }hgeom,\,hchar}
\end{equation*}
\begin{enumerate}
\item \emph{Lie data} (Section~\ref{l23:sec:liedata}): the linear algebra of $\lxxiiig$ that the construction uses.
\item \emph{The universal affine VOA} $V^k(\lxxiiig)$ is built from Lie data for every non-critical level $k$
  (Section~\ref{l23:sec:universal}), with the Sugawara conformal vector.
\item \emph{Lie data exist for every finite-dimensional simple $\lxxiiig$} (Section~\ref{l23:sec:everysimple}), by Schur's
  lemma for the Casimir.
\item \emph{The simple quotient} $L_k(\lxxiiig)$ is simple, its weight-one space is spanned by the currents, and
  singular vectors of positive weight vanish in it (Section~\ref{l23:sec:simple}).
\item \emph{The Deligne instances} (Section~\ref{l23:sec:deligne}): for each of the seven types, $L_k(\lxxiiig)$ at the
  Deligne level is fed into the capstone of Chapter~\ref{ch:L24}.
\end{enumerate}
Steps 3--5 are only useful if some concrete Lean Lie algebra is known to be simple. Sections~\ref{l23:sec:sln},
\ref{l23:sec:lattice} and~\ref{l23:sec:sc} provide them: $\mathfrak{sl}_n$; the weight-one spaces of the
$A_2,D_4,E_6,E_7,E_8$ root-lattice vertex algebras (on paper $\mathfrak{sl}_3,\mathfrak{so}_8,\mathfrak e_6,\mathfrak e_7,\mathfrak e_8$); and
algebras of dimensions $14,52,78$ from kernel-checked tables of structure constants, modelling $\mathfrak g_2,\mathfrak f_4,\mathfrak e_6$
(the Cartan types are named on paper only, Remark~\ref{l23:rem:type}); with their dimensions and the
normalization of the invariant form. Section~\ref{l23:sec:cert} is engineering rather than mathematics: it realizes the
fixed ``certificate basis'' of the Python computations inside these algebras.

\paragraph{What is proved, what is conditional, and what is not claimed.}
\begin{itemize}
\item \textbf{Proved, with only the intended hypotheses} (\texttt{\#print axioms} lists only Lean's three standard axioms):
  $V^k(\lxxiiig)$ is a VOA of central charge $k\dim\lxxiiig/(k+h)$ for Lie data and $k+h\ne0$, where $2h$ is the eigenvalue of the
  adjoint Casimir \emph{in the normalization of the chosen form $\kappa$} (not, in general, the dual Coxeter number $\lxxiiihv$; see
  below); Lie data exist for every
  finite-dimensional simple $\lxxiiig$ with a nondegenerate symmetric invariant form; $L_k(\lxxiiig)$ is simple, with weight
  one spanned by the currents; $V^k(\lxxiiig)_1\cong\lxxiiig$ \emph{as vector spaces} for $k\ne0$ (a linear isomorphism
  \lean{weightOneEquiv}; compatibility with the bracket $u_{(0)}v\leftrightarrow[x,y]$ is not proved for general $\lxxiiig$);
  positive-weight singular vectors lie in the
  maximal ideal; $\mathfrak{sl}_n$ ($n\ge2$), the five lattice algebras and the three table algebras are simple, with the
  stated dimensions and (per type, via the \lean{*_standard} theorems) normalizations; the certificate tables are the structure constants of Lie data; and for every
  Deligne type the input ``current-jet surjectivity'' of the capstone holds.
\item \textbf{Naming, fixed; and a normalization caveat} (release audit 2026-10-27, findings M1--M2). (i) The audit found that the
  declaration for the \emph{universal} algebra $V^k(\lxxiiig)$ of a simple $\lxxiiig$ was called \lean{simpleAffineVOA} (``simple''
  referring to $\lxxiiig$), inviting the false reading $L_k(\lxxiiig)$, and that the simple quotient had no name. Both are fixed: the
  universal algebra is now \lean{universalAffineVOA} (the old name survives only as a deprecated alias), and the simple quotient
  $L_k(\lxxiiig)$, the term \lean{simpleVOA (simpleData 𝔤 κ hinv hnd) k hk}, is named \lean{simpleAffineQuotientVOA}
  (Section~\ref{l23:sec:simplename}). (ii) The number $h$ in
  \lean{universalAffineVOA} is half an eigenvalue picked by \lean{Classical.choose}, in the normalization of $\kappa$; it is computed only
  for the Killing form ($h=\frac12$). \lean{standardAffineVOA} takes $h$ as a free nonzero parameter \lean{hv}. That $h=\lxxiiihv$ is
  established only type by type, through the \lean{*_standard} theorems (Theorem~\ref{l23:thm:normalization} and
  Section~\ref{l23:sec:sc}), and for $\mathfrak g_2,\mathfrak f_4,\mathfrak e_6$ the longness of the chosen root is itself not checked in
  Lean. For $\mathfrak{sl}_n$ it is not established in Lean at all.
\item \textbf{Conditional.} The statement that the Deligne case is \emph{classically free}
  (\lean{classicallyFree_of_two_inputs} and its seven specializations) is proved only \emph{assuming} two named inputs,
  \lean{hgeom} (a ``geometric carrier embedding'') and \lean{hchar} (a ``carrier character identity''), explained in
  Section~\ref{l23:sec:deligne}. This chapter does not discharge them; for the presentation built here it cannot
  (warning box in Section~\ref{l23:warn:fullkernel}).
\item \textbf{Not claimed in Lean}, although true on paper or in the literature:
  (a) $L_k(\lxxiiig)_1\cong\lxxiiig$ for general $\lxxiiig$ (only spanning is proved);
  (b) that the algebras named \lean{e8Lie}, \lean{f4Lie}, \dots\ are of the named Cartan type (the classification of simple
  Lie algebras is not formalised; Remark~\ref{l23:rem:type}). Lean proves \emph{concrete simple Lie algebras of dimensions
  $14,52,78,133,248$}, built as the exceptional ones; no Lean statement identifies a Cartan type, and dimension $78$ could not do so
  in any case ($\dim B_6=\dim C_6=\dim E_6=78$, all of rank $6$ with $72$ roots);
  (c) the Frenkel--Kac isomorphism $V_{E_8}\cong L_1(\mathfrak e_8)$: the library uses the lattice VOA $V_{E_8}$ \emph{only} as a
  source of a concrete simple Lie algebra (its weight-one space with the bracket $[u,v]=u_{(0)}v$), and no Lean statement
  relates \lean{e8VOA} to \lean{standardAffineVOA e8Lie}; wherever Frenkel--Kac is mentioned below it is motivation;
  (d) the identification of $L_k(\lxxiiig)$ at the Deligne levels with four-dimensional chiral algebras (physics literature);
  (e) any isomorphism between the three Lean models of $\mathfrak e_6$;
  (f) that the form used by \lean{slAffineVOA} is the trace form.
\end{itemize}

\paragraph{Where this goes.}
Chapter~\ref{ch:L24} (presentations, carriers, divisor comparison) builds the explicit presentations
$\mathrm{Sym}\,\lxxiiig/(I_2^{\mathrm{ns}})$ of the $C_2$-algebra of $L_k(\lxxiiig)$ at the Deligne levels from
Theorem~\ref{l23:thm:singular}, Proposition~\ref{l23:prop:clsG} and the Lie data of Sections~\ref{l23:sec:sc}
and~\ref{l23:sec:cert}, and states the capstone quoted in Section~\ref{l23:sec:deligne}. Chapter~\ref{ch:L25} proves lower
bounds on $\dim L_k(\lxxiiig)_N$. Chapter~\ref{ch:L26} (fusion modules) uses the certificate bases of
Section~\ref{l23:sec:cert} in its lattice realizations. Chapter~\ref{ch:L27} states the final theorems for the seven types,
conditional on named literature inputs.

\paragraph{What we rely on.}
Chapter~\ref{ch:L01} for formal distributions, modes ($\lean{mode f n}=f_{(n)}$, the coefficient of
$z^{-n-1}$) and locality; Chapter~\ref{ch:L02} for invariant forms (\lean{IsInvariant}), the affine Kac--Moody
algebra $\widehat{\lxxiiig}$, the vacuum module \lean{vacuumSpace} and its nonvanishing by relative PBW;
Chapter~\ref{ch:L03} for the structures \lean{PreVertexAlgebra}, \lean{VertexOperatorAlgebra} and the
reconstruction theorem; Chapter~\ref{ch:L04} for the Lie algebra structure $[u,v]=u_{(0)}v$ on weight one, ideals,
quotients, the maximal graded ideal and \lean{IsSimple}; Chapter~\ref{ch:L05} for lattice vertex algebras
\lean{latticeVOAGen} and the FLM cocycle; the files \leanpath{VOA/\allowbreak Examples/\allowbreak AffineSl3*.lean} (chapter in preparation) for the $\mathfrak{sl}_3$ construction, of which the
general one is a port; Chapter~\ref{ch:L11} for Li's filtration, jets and generator words; Chapter~\ref{ch:L14} for
modular forms and Eisenstein series; Chapter~\ref{ch:02} for integral lattices and theta series.

%======================================================================
\section{Reading the Lean, and the notation of this chapter}\label{l23:sec:readlean}
%======================================================================

\paragraph{A primer is in Chapter~\ref{ch:01}.} Section~\ref{ch01:sec:lean} (``The Lean toolkit: reading Lean for
mathematicians'') and the glossary of Section~\ref{ch01:sec:glossary} teach everything needed to read the boxes below. For
convenience, here are the constructs that occur in this chapter.
\begin{itemize}
\item \textbf{Omitted bodies.} A quotation stops where the proof begins. A box ending in \lean{:= by} (or \lean{:=}) after a
  \lean{theorem} has had its proof removed; it is not truncated by accident. For a \lean{def} or \lean{structure} ending in
  \lean{where} or \lean{:=}, the body that follows defines the object; where that body is data rather than a proof
  (\lean{ofNondegenerate}, \lean{smulForm}, \lean{reindex}, \lean{ofSimple}, \lean{standardData}, the truncated \lean{pairOK},
  the \lean{casimir} field of \lean{sl3Data}) the dictionary box says what it contains. Docstrings and comments are also
  omitted. A one-line proof term is kept when it is informative (e.g.\ \lean{simpleVOA_isSimple}).
\item \textbf{Arguments.} \lean{(x : X)} is an explicit argument, \lean{{x : X}} an implicit one (inferred by Lean), \lean{[C X]}
  an instance argument (a structure, such as \lean{[Fintype ι]} ``$\iota$ is finite'' or \lean{[DecidableEq ι]} ``equality on $\iota$
  can be decided'', which Lean finds automatically). For reading, all three are universally quantified hypotheses.
  \lean{Type*} means ``some type'' (in some universe).
\item \textbf{\lean{variable} and \lean{include}.} A \lean{variable} line declares arguments once for a whole section of a
  file; they become arguments of every declaration below that uses them. They are therefore hypotheses of a theorem even
  though they are not printed in its signature. Every dictionary box lists the relevant section variables. \lean{include h}
  forces a section variable \lean{h} to be a hypothesis even of declarations that do not mention it.
\item \textbf{\lean{abbrev}, \lean{irreducible_def}, \lean{noncomputable}.} \lean{abbrev} is a definition that Lean unfolds
  eagerly (for reading, the same as \lean{def}); \lean{irreducible_def} is a definition that is opaque except through its
  defining equation (used for speed). \lean{noncomputable} means Lean cannot run it as a program (it uses choice).
\item \textbf{Anonymous constructors and functions.} \lean{⟨a, b⟩} builds a structure from its fields in order, e.g.\ an element
  of a subspace from a vector \lean{a} and a proof \lean{b} that it lies in the subspace. \lean{fun i => e} is the function
  $i\mapsto e$; \lean{![a, b, c]} is a vector literal.
\item \textbf{Finite index sets.} \lean{Fin n} is $\{0,\dots,n-1\}$; \lean{∑ i, f i} sums over a finite type; \lean{Fintype.card ι} is $|\iota|$.
\item \textbf{Kernel checks.} \lean{decide +kernel} proves a decidable statement (typically ``this Boolean function of an integer
  table evaluates to \lean{true}'') by letting Lean's kernel evaluate it. The kernel is the trusted core of Lean, so such a
  proof is as reliable as any other; the Python scripts that \emph{produced} the tables are not trusted and need not be.
\end{itemize}

\paragraph{Notation.} The letter $h$ is overloaded in the literature; this chapter keeps the meanings apart as follows.
\begin{center}\small
\begin{tabular}{lp{0.72\textwidth}}
\toprule
symbol & meaning\\
\midrule
$\kappa$ & a symmetric invariant bilinear form on $\lxxiiig$ (Lean \lean{κ})\\
$b_a$, $b^a$ & a basis of $\lxxiiig$ and its $\kappa$-dual family, $\kappa(b_a,b^c)=\delta_{ac}$\\
$\kappa_{ab}$, $\kappa^{ab}$ & the Gram matrix $\kappa(b_a,b_b)$ (\lean{formTable}); the $b$-th coordinate of $b^a$ (\lean{kappaInv}), the inverse Gram matrix\\
$f^{ab}{}_c$ & the $c$-th coordinate of $[b_a,b_b]$ (\lean{structureConstant a b c}). No summation convention: every sum is written\\
$h$ & the \emph{Casimir number} of $(\lxxiiig,\kappa)$: $\sum_a[b_a,[b^a,x]]=2hx$ (field \lean{h} of Lie data). It depends on $\kappa$\\
$\lxxiiihv$ & the dual Coxeter number of $\lxxiiig$, a fixed integer ($2,3,4,6,9,12,18,30$ for $A_1,A_2,G_2,D_4,F_4,E_6,E_7,E_8$). $h=\lxxiiihv$ exactly when $\kappa$ is standard (long roots of squared length $2$)\\
$\nu$ & a free nonzero complex parameter (Lean \lean{hv} in \lean{standardData}); \emph{not} assumed to be $\lxxiiihv$\\
$h_\alpha$, $e_\alpha$ & coroot (Cartan) vectors and root vectors of a concrete Lie algebra (Sections~\ref{l23:sec:sln}--\ref{l23:sec:sc})\\
$H$ & the diagonal matrix $\mathrm{diag}(1,-1)\in\mathfrak{sl}_2$\\
$J^a_m$ & the mode $b_a\otimes t^m$ acting on the vacuum module; $J^a(z)=\sum_mJ^a_mz^{-m-1}$\\
$x_{(n)}$ & the $n$-th product (mode) of the field of $x$: $Y(x,z)=\sum_nx_{(n)}z^{-n-1}$; for a current, $x_{(n)}=J^x_n$\\
$k$ & the level\\
\bottomrule
\end{tabular}
\end{center}

%======================================================================
\section{The physical picture}\label{l23:sec:physics}
%======================================================================

A two-dimensional conformal field theory with a continuous internal symmetry group $G$ has conserved currents
$J^a(z)$, one for each element $t^a$ of a basis of the Lie algebra $\lxxiiig$ of $G$. Holomorphy and
conformal weight one force their operator product expansion to be
\begin{equation}\label{l23:eq:JJ}
J^a(z)\,J^b(w)\sim \frac{k\,\kappa(t^a,t^b)}{(z-w)^2}+\frac{\sum_c f^{ab}{}_c\,J^c(w)}{z-w},
\end{equation}
where $f^{ab}{}_c$ are the structure constants, $\kappa$ is an invariant symmetric form, and the number $k$ (the
\emph{level}) measures the strength of the Schwinger term, i.e.\ the two-point function of the currents. In
modes $J^a(z)=\sum_m J^a_m z^{-m-1}$ this is the affine Kac--Moody algebra
\begin{equation}\label{l23:eq:affine}
[J^a_m,J^b_n]=\sum_c f^{ab}{}_c\,J^c_{m+n}+m\,k\,\kappa(t^a,t^b)\,\delta_{m+n,0}.
\end{equation}
The \emph{vacuum module} is the space of states generated from a vacuum $\lxxiiivac$ annihilated by all
$J^a_m$ with $m\ge0$; the state--field correspondence turns it into a vertex algebra, the \emph{universal affine
vertex algebra} $\lxxiiiV$. Its stress tensor is not an extra datum: the \emph{Sugawara construction}
\begin{equation}\label{l23:eq:sugawara}
T(z)=\frac{1}{2(k+h)}\sum_{a}:\!J^{a}(z)J_{a}(z)\!:\,,\qquad c=\frac{k\dim\lxxiiig}{k+h},
\end{equation}
($J_a$ the current of the $\kappa$-dual basis element; $h$ the Casimir number of $(\lxxiiig,\kappa)$, defined in
Definition~\ref{l23:def:liedata} below by $\sum_a[b_a,[b^a,x]]=2hx$, which equals the dual Coxeter number $\lxxiiihv$ when $\kappa$ is
normalized so that long roots have squared length $2$) builds it from the currents. It fails exactly at the \emph{critical level} $k=-h$, where $T$ becomes central.

Two classes of examples will matter. At positive integer level, $L_k(\lxxiiig)$ is the chiral algebra of the
Wess--Zumino--Witten model of $G$; at level $1$ for simply-laced $\lxxiiig$ it is isomorphic to the lattice vertex algebra
of the root lattice (the Frenkel--Kac theorem). The library uses the lattice side of this picture \emph{only} to obtain
a concrete $\mathfrak e_8$: the conserved currents of eight chiral bosons on the $E_8$ torus (the theory underlying the
$E_8\times E_8$ heterotic string) form a Lie algebra, and that Lie algebra is the library's $\mathfrak e_8$. The
isomorphism $V_{E_8}\cong L_1(\mathfrak e_8)$ itself is \emph{not} formalised, and nothing below depends on it. At
the \emph{negative} levels $k=-\lxxiiihv/6-1$ of the Deligne series, $L_k(\lxxiiig)$ is (by the work of Beem,
Lemos, Liendo, Peelaers, Rastelli and van Rees, and Arakawa--Moreau) the chiral algebra of a four-dimensional
$\mathcal N=2$ superconformal field theory of rank one, whose Higgs branch is the minimal nilpotent orbit of
$\lxxiiig$; the two-dimensional central charge is $c_{2d}=-12c_{4d}=-2\lxxiiihv-2$. These are the theories studied in
Chapters~\ref{ch:L24}--\ref{ch:L27}. This identification is physics literature; it motivates the Lean statements but is
not one of them.

%======================================================================
\section{Lie data}\label{l23:sec:liedata}
%======================================================================

The construction of $\lxxiiiV$ uses a finite-dimensional Lie algebra only through a small amount of linear algebra:
a basis $b_1,\dots,b_N$, the dual basis $b^1,\dots,b^N$ with respect to $\kappa$, and the value of the adjoint
Casimir. The library packages these as a structure.

\begin{definition}[Lie data]\label{l23:def:liedata}
Let $\lxxiiig$ be a complex Lie algebra and $\kappa$ a bilinear form on $\lxxiiig$. \emph{Lie data} for
$(\lxxiiig,\kappa)$ indexed by a finite set $\iota$ consist of
\begin{enumerate}
\item a proof that $\kappa$ is symmetric and invariant: $\kappa(x,y)=\kappa(y,x)$ and $\kappa([x,y],z)=\kappa(x,[y,z])$;
\item a basis $(b_i)_{i\in\iota}$ of $\lxxiiig$;
\item a family $(b^i)_{i\in\iota}$ (the \emph{dual family}) with $\kappa(b_i,b^j)=\delta_{ij}$;
\item a complex number $h$ (the \emph{Casimir number}; it equals the dual Coxeter number $\lxxiiihv$ when $\kappa$ is the
  standard normalization, and scales inversely with $\kappa$ otherwise);
\item a proof that the adjoint Casimir acts by $2h$:
\begin{equation}\label{l23:eq:cas}
\lxxiiiCas(x):=\sum_{i\in\iota}[b_i,[b^i,x]]=2h\,x\qquad\text{for all }x\in\lxxiiig .
\end{equation}
\end{enumerate}
\end{definition}
\noindent File \leanpath{VOA/\allowbreak AffineGeneral/\allowbreak LieData.lean}.
\begin{leancode}
structure LieData {𝔤 : Type*} [LieRing 𝔤] [LieAlgebra ℂ 𝔤]
    (κ : LinearMap.BilinForm ℂ 𝔤) (ι : Type*) [Fintype ι] [DecidableEq ι] where
  invariant : κ.IsInvariant
  basis : Module.Basis ι ℂ 𝔤
  dual : ι → 𝔤
  pair : ∀ i j, κ (basis i) (dual j) = if i = j then 1 else 0
  h : ℂ
  casimir : ∀ x : 𝔤, ∑ i, ⁅basis i, ⁅dual i, x⁆⁆ = (2 * h) • x
\end{leancode}

\begin{dictionary}
\lean{LinearMap.BilinForm ℂ 𝔤} is a $\C$-bilinear form $\lxxiiig\times\lxxiiig\to\C$, applied as \lean{κ x y}.
\lean{κ.IsInvariant} is the structure of Chapter~\ref{ch:L02} (Definition~\ref{l02:def:invariant},
\leanpath{VOA/\allowbreak Foundations/\allowbreak AffineLie.lean}) with two fields, \lean{symm : ∀ x y, κ x y = κ y x} and
\lean{invariant : ∀ x y z, κ ⁅x, y⁆ z = κ x ⁅y, z⁆}; it is plain symmetry, not Mathlib's \lean{IsSymm} (which over
$\C$ is the Hermitian condition). \lean{Module.Basis ι ℂ 𝔤} is a basis indexed by $\iota$; \lean{[Fintype ι]} says
$\iota$ is finite, so $\lxxiiig$ is finite-dimensional, and \lean{[DecidableEq ι]} is needed for the Kronecker
\lean{if i = j then 1 else 0}. \lean{⁅x, y⁆} is the Lie bracket and \lean{•} scalar multiplication.

Every field is a hypothesis about $\lxxiiig$, not an axiom. Nondegeneracy of $\kappa$ is not a field but is a
consequence: the pairing field forces $\kappa(y,x)=0\ \forall y\Rightarrow x=0$
(\lean{LieData.form_right_nondegenerate} in the same file). The field \lean{casimir} is the only genuinely
Lie-theoretic input; it is \emph{not} automatic for a general Lie algebra (for a reductive non-simple algebra
the Casimir has several eigenvalues), and it is proved for every simple $\lxxiiig$ in
Section~\ref{l23:sec:everysimple}.

Could the structure be vacuous? No: Section~\ref{l23:sec:everysimple} constructs it for every simple $\lxxiiig$, and
\lean{sl3Data} (Section~\ref{l23:sec:sl3}) is an explicit instance. Is the name ``$h$'' honest? $h$ is
\emph{defined} by \eqref{l23:eq:cas}; it equals the dual Coxeter number $\lxxiiihv$ exactly when $\kappa$ is the
standard normalization (long roots of squared length $2$). In general $h$ scales inversely with $\kappa$
(Definition~\ref{l23:def:smulform}). Same as the textbook.
\end{dictionary}

\begin{example}[$\mathfrak{sl}_2$ with the trace form]\label{l23:ex:sl2}
Take $\lxxiiig=\mathfrak{sl}_2$, $\kappa(x,y)=\tr(xy)$, basis $b_0=e$, $b_1=H$, $b_2=f$ ($H=\mathrm{diag}(1,-1)$), with
$[H,e]=2e$, $[H,f]=-2f$, $[e,f]=H$. Since $\tr(ef)=1$, $\tr(HH)=2$ and all other traces of products vanish, the dual family is
$b^0=f$, $b^1=\tfrac12H$, $b^2=e$. On $x=e$:
$\lxxiiiCas(e)=[e,[f,e]]+[H,[\tfrac12H,e]]+[f,[e,e]]=[e,-H]+\tfrac12[H,2e]+0=2e+2e$. So $\lxxiiiCas(e)=4e$ and
$h=2=\lxxiiihv(\mathfrak{sl}_2)$, as it must be, since the trace form is the standard normalization. This example is
continued in Examples~\ref{l23:ex:sl2sug} and~\ref{l23:ex:sl2sing}.
\end{example}

\textbf{Why this is the right definition.} The construction of $V^k(\lxxiiig)$ uses the Lie algebra only through
brackets of basis vectors, the Gram matrix of $\kappa$ (for the central term) and the double contraction
$\sum_a[b_a,[b^a,\cdot]]$ (for the Sugawara normalization). Packaging exactly these, rather than assuming
$\lxxiiig$ simple, lets the same construction serve every case, including reductive algebras with a suitable form, and
moves the one genuinely Lie-theoretic fact (the Casimir is a scalar) to a separate theorem.

\begin{physics}
Lie data are exactly what the current algebra \eqref{l23:eq:JJ} and the Sugawara tensor \eqref{l23:eq:sugawara}
need: a choice of currents $J^a$ (the basis), the two-point normalization $\kappa$ (the ``metric'' raising and
lowering colour indices, here encoded by the dual family), and the Casimir eigenvalue in the adjoint, which is the
one-loop shift $k\mapsto k+h$ in the Sugawara normalization (the famous ``$k+h^\vee$'' of WZW models).
\end{physics}

\begin{status}
A definition (no proof obligation). Instances are constructed unconditionally (Sections~\ref{l23:sec:everysimple},
\ref{l23:sec:sl3}).
\end{status}

From the fields one derives, in coordinates, the identities the Sugawara computation needs. Write
$f^{ab}{}_c$ for the structure constants (\lean{structureConstant a b c}, the $c$-th coordinate of $[b_a,b_b]$, so
$[b_a,b_b]=\sum_cf^{ab}{}_cb_c$), $\kappa_{ab}=\kappa(b_a,b_b)$ for the Gram matrix (\lean{formTable}) and $\kappa^{ac}$ for
the $c$-th coordinate of $b^a$ (\lean{kappaInv}; so $b^a=\sum_c\kappa^{ac}b_c$, and $(\kappa^{ac})$ is the inverse of
$(\kappa_{ab})$). The position of an index (up or down) carries no meaning beyond this; there is no summation convention, and
every sum is written out.

\begin{proposition}[Finite identities]\label{l23:prop:finite}
For Lie data as above:
\begin{enumerate}
\item (Casimir in coordinates) $\sum_{a,d}\big(\sum_b\kappa^{ab}f^{bc}{}_d\big)f^{ad}{}_e=2h\,\delta_{ce}$;
\item (trace of $\lxxiiiad$ vanishes) $\sum_a f^{ac}{}_a=0$ for every $c$ (note $\tr\lxxiiiad(b_c)=\sum_af^{ca}{}_a=-\sum_af^{ac}{}_a$);
\item (central charge numerator) $\sum_{a,b}\kappa^{ab}\kappa_{ab}=|\iota|=\dim\lxxiiig$.
\end{enumerate}
\end{proposition}
\noindent File \leanpath{VOA/\allowbreak AffineGeneral/\allowbreak LieData.lean}.
\begin{leancode}
theorem struct_casimir_contract (c e : ι) :
    (∑ a, ∑ d, (∑ b, D.kappaInv a b * D.structureConstant b c d) * D.structureConstant a d e) =
      if c = e then 2 * D.h else 0 := by

theorem structureConstant_diag_sum (c : ι) : (∑ a, D.structureConstant a c a) = 0 := by

theorem trace_kappaInv_formTable :
    (∑ a, ∑ b, D.kappaInv a b * D.formTable a b) = Fintype.card ι := by
\end{leancode}

\begin{dictionary}
\lean{D : LieData κ ι} is a variable of the section; \lean{D.kappaInv a b} is \lean{D.basis.repr (D.dual a) b},
the $b$-th coordinate of $b^a$, which is the inverse Gram matrix by \lean{kappaInv_kappa} in the same file.
\lean{Fintype.card ι} is $|\iota|=\dim\lxxiiig$. Same as the textbook.
\end{dictionary}

\textbf{Proof on paper.} (1) is the $e$-th coordinate of \eqref{l23:eq:cas} at $x=b_c$, after expanding
$b^a=\sum_b\kappa^{ab}b_b$. (2): the trace of $\lxxiiiad(x)$ is $\sum_a\kappa([x,b_a],b^a)$; invariance gives
$\kappa([x,b_a],b^a)=-\kappa(b_a,[x,b^a])$, and re-expanding in the dual basis shows the sum equals its own negative,
hence vanishes (no simplicity is needed: any Lie algebra with a nondegenerate invariant form is unimodular).
(3): row by row, $\sum_b\kappa^{ab}\kappa_{ab}=1$ because $\kappa^{-1}\kappa=1$; summing over $a$ gives $|\iota|$. \qed

\begin{physics}
(1) is the statement that the Casimir appears in the double contraction $f f$ of the current OPE; it produces the
$h$ in $k+h$. (2) makes each current a Virasoro primary of weight one (no ``anomalous'' $\partial J$ term in
$T(z)J(w)$). (3) is why the central charge numerator is $k\dim\lxxiiig$. \emph{Heuristically}, each of the $\dim\lxxiiig$
currents contributes one free-boson-like unit; this is literally true only as $k\to\infty$, where $c=k\dim\lxxiiig/(k+h)\to\dim\lxxiiig$
and the currents become $\dim\lxxiiig$ free bosons. At finite $k$ the interaction (the $f$ terms) reduces it by the factor $k/(k+h)$.
\end{physics}

\begin{status}
Proved; \texttt{\#print axioms} standard for all three (project record, Section~\ref{l23:sec:verification}).
\end{status}

Lie data can be built from a nondegenerate form, rescaled, and reindexed.

\begin{definition}[Constructions of Lie data]\label{l23:def:smulform}
(i) Given a symmetric invariant nondegenerate $\kappa$, a basis $b$ and a number $h$ with
$\sum_i[b_i,[b^i,x]]=2hx$ for the $\kappa$-dual basis $b^i$, there is Lie data with these components.
(ii) If $D$ is Lie data for $\kappa$ and $c\ne0$, then $(b_i,\,c^{-1}b^i,\,h/c)$ is Lie data for $c\kappa$.
(iii) Lie data can be transported along a bijection $\iota\simeq\iota'$.
\end{definition}
\noindent File \leanpath{VOA/\allowbreak AffineGeneral/\allowbreak LieData.lean}.
\begin{leancode}
def ofNondegenerate (hinv : κ.IsInvariant) (hnd : LinearMap.BilinForm.Nondegenerate κ) (b : Module.Basis ι ℂ 𝔤)
    (h : ℂ) (hcas : ∀ x : 𝔤, ∑ i, ⁅b i, ⁅LinearMap.BilinForm.dualBasis κ hnd b i, x⁆⁆ = (2 * h) • x) :
    LieData κ ι where

def smulForm (D : LieData κ ι) (c : ℂ) (hc : c ≠ 0) : LieData (c • κ) ι where

theorem smulForm_h (D : LieData κ ι) (c : ℂ) (hc : c ≠ 0) : (D.smulForm c hc).h = D.h / c := rfl

noncomputable def reindex {ι' : Type*} [Fintype ι'] [DecidableEq ι'] (D : LieData κ ι) (e : ι ≃ ι') :
    LieData κ ι' where
\end{leancode}

\begin{dictionary}
\lean{LinearMap.BilinForm.Nondegenerate κ} is Mathlib's nondegeneracy; \lean{dualBasis κ hnd b} is Mathlib's
$\kappa$-dual basis. \lean{c • κ} is the rescaled form; \lean{ι ≃ ι'} is a bijection. The rescaling law $h\mapsto h/c$
is forced: the dual family scales by $c^{-1}$, so the Casimir does too. Same as the textbook.
\end{dictionary}

\textbf{Proof on paper.} (ii): $c\kappa(b_i,c^{-1}b^j)=\delta_{ij}$ and $\sum_i[b_i,[c^{-1}b^i,x]]=c^{-1}2hx=2(h/c)x$.
(iii) is a reindexing of finite sums. \qed

\begin{physics}
Rescaling $\kappa$ by $c$ while keeping the currents fixed rescales the level $k\mapsto k/c$ in \eqref{l23:eq:JJ} and
the dual Coxeter number $h\mapsto h/c$; only the ratio $k/h$ (and hence $c=\dim\lxxiiig\,k/(k+h)$) is physical. This is
why a level is meaningless until a normalization of $\kappa$ is fixed.
\end{physics}

\begin{status}
Definitions; the identity \lean{smulForm_h} holds by \lean{rfl}.
\end{status}

%======================================================================
\section{The universal affine vertex algebra \texorpdfstring{$V^k(\mathfrak g)$}{V\^{}k(g)}}\label{l23:sec:universal}
%======================================================================

\subsection{The vacuum module and the currents}

Recall from Chapter~\ref{ch:L02} (Definitions~\ref{l02:def:affine} and~\ref{l02:def:vacuum}) the affine Lie algebra
$\widehat{\lxxiiig}=\lxxiiig\otimes\C[t,t^{-1}]\oplus\C K\,(\oplus\C d)$ and its vacuum module, realized in Lean as a
quotient of the universal enveloping algebra of $\lxxiiig\otimes\C[t,t^{-1}]\oplus\C K\oplus\C\ell_0$ by the left ideal
(\lean{vacuumIdeal}, \leanpath{VOA/\allowbreak Foundations/\allowbreak AffineLie.lean}) generated by the non-negative loop modes
$\lxxiiig\otimes t^{n}$ ($n\ge0$), the derivation $\ell_0$, and $K-k$. The general construction works on that carrier.

\begin{definition}[Vacuum module and current modes]\label{l23:def:vacuum}
For Lie data $D$ and $k\in\C$, $\lxxiiiV$ is the vacuum module of $\widehat{\lxxiiig}$ at level $k$, and $J^a_m$ is the
action of $b_a\otimes t^m$ on it.
\end{definition}
\noindent Files \leanpath{VOA/\allowbreak AffineGeneral/\allowbreak Currents.lean},
\leanpath{VOA/\allowbreak AffineGeneral/\allowbreak Grading.lean}.
\begin{leancode}
abbrev Vacuum (k : ℂ) := vacuumSpace κ D.invariant k

noncomputable def J (k : ℂ) (a : ι) (m : ℤ) : Module.End ℂ (Vacuum D k) :=
  vacuumRho κ D.invariant k (ofLoop κ D.invariant (D.generator a) m)

theorem current_comm (k : ℂ) (a b : ι) (m n : ℤ) :
    J D k a m * J D k b n - J D k b n * J D k a m
      = (∑ c, structureConstant D a b c • J D k c (m + n))
        + (if m + n = 0 then ((m : ℂ) * k * formTable D a b) • LinearMap.id else 0) := by

noncomputable def vacuumA (k : ℂ) : Vacuum D k :=
  vacuumVector κ D.invariant k

theorem vacuumA_ne_zero (k : ℂ) : vacuumA D k ≠ 0 :=
  vacuumVector_ne_zero κ D.invariant k
\end{leancode}

\begin{dictionary}
\lean{vacuumSpace κ hκ k}, \lean{vacuumRho}, \lean{ofLoop} and \lean{vacuumVector} are the objects of
Chapter~\ref{ch:L02} (\leanpath{VOA/\allowbreak Foundations/\allowbreak AffineLie.lean}): the vacuum module, the action
of $\widehat{\lxxiiig}$ on it, the loop element $x\otimes t^m$, and the class of $1$. \lean{Module.End ℂ (Vacuum D k)} is
the algebra of linear endomorphisms; \lean{*} is composition, so the left side of \lean{current_comm} is the
commutator. \lean{LinearMap.id} is the identity. \lean{D.generator a} is an \lean{abbrev} for \lean{D.basis a}, i.e.\ $b_a$, and
\lean{D.dualGenerator a} for \lean{D.dual a}, i.e.\ $b^a$ (\leanpath{VOA/\allowbreak AffineGeneral/\allowbreak LieData.lean}).
\lean{current_comm} is \eqref{l23:eq:affine} verbatim, with
$\kappa(t^a,t^b)$ the Gram entry \lean{formTable D a b}.

The nonvanishing \lean{vacuumA_ne_zero} is the delicate point of the whole construction: if the vacuum were $0$,
every subsequent statement would be about the zero space. It is proved by \lean{AffineLie.vacuumVector_ne_zero},
which rests on the \emph{theorem} \lean{VOA.pbw} (relative PBW, Chapter~\ref{ch:L02},
Theorem~\ref{l02:thm:vacind}); there is no project axiom. The file's docstrings claim otherwise (see the warning box in Section~\ref{l23:warn:freefield}).
\end{dictionary}

\textbf{Proof on paper.} The commutator of two operators of a representation is the representation of the bracket;
the bracket of $\widehat{\lxxiiig}$ is the loop bracket plus the cocycle $m\delta_{m+n,0}\kappa(x,y)K$, and $K$ acts by
$k$ on the vacuum module. Nonvanishing: by the relative PBW theorem, the induced module
$U(\widehat{\lxxiiig})\otimes_{U(\lxxiiig[t]\oplus\C K)}\C_k$ restricts to $\lxxiiig[t]\oplus\C K$ with $\C_k$ as a direct summand, so
$1\otimes1\ne0$. \qed

\begin{physics}
$J^a_m$ are the Laurent modes of the conserved currents; \lean{current_comm} is the current algebra with Schwinger
term $mk\kappa$; the vacuum is the $SL_2$-invariant state annihilated by all non-negative modes. Nonvanishing of the
vacuum says the theory has a state at all, which is not a formality: the vacuum module is defined as a quotient.
\end{physics}

\begin{status}
Proved; \texttt{\#print axioms} standard for \lean{current_comm} and \lean{vacuumA_ne_zero} (project record,
Section~\ref{l23:sec:verification}).
\end{status}

\begin{warning}[title={Trap: Docstring defect: ``free-field evaluation, without PBW''}]\label{l23:warn:freefield}
In \leanpath{VOA/\allowbreak AffineGeneral/\allowbreak Grading.lean} the docstring of \lean{vacuumA_ne_zero} reads ``The
vacuum is nonzero by its actual polynomial evaluation, without PBW'', and the module docstring says the vacuum line
is ``proved nonzero using the explicit free-field evaluation''; \leanpath{Currents.lean} says ``Its previously
constructed free-field evaluation proves that this space is nonzero''. These sentences were inherited from the
$\mathfrak{sl}_3$ files (\leanpath{VOA/\allowbreak Examples/\allowbreak AffineSl3*.lean}, where a Wakimoto realization was used). In the general files they are
\textbf{false}: the proof term printed above is \lean{vacuumVector_ne_zero}, i.e.\ relative PBW. (The second paragraph of
the \leanpath{Grading.lean} module docstring states the correct fact, so the file contradicts itself.) Listed in
Section~\ref{l23:sec:defects}.
\end{warning}

\subsection{Grading, Sugawara vector and the vertex algebra}

The operator $\ell_0$ of Chapter~\ref{ch:L02} is the action of the derivation $t\,\tfrac{d}{dt}$, i.e.\
$[\ell_0,x\otimes t^n]=n\,(x\otimes t^n)$ (\lean{ell0_lie_ofLoop}); it kills the vacuum, and a creation mode
$x\otimes t^{-m}$ ($m\ge1$) \emph{lowers} its eigenvalue by $m$. So $\ell_0$ measures \emph{minus} the energy: it
diagonalizes $\lxxiiiV$ with eigenvalues $0,-1,-2,\dots$ and finite-dimensional eigenspaces, and a state of
weight $n\ge0$ is an $\ell_0$-eigenvector of eigenvalue $-n$ (on $\lxxiiiV$, $\ell_0=-L_0$). The library takes as grading
$\lxxiiiV_n=\ker(\ell_0+n)$, the $\ell_0$-eigenspace of eigenvalue $-n$ (\lean{gradingZA D k n}), proves that it is an internal direct sum, that
$\lxxiiiV_0=\C\lxxiiivac$ and $\lxxiiiV_n=0$ for $n<0$, and builds the vertex operators $Y(u,z)$ by the reconstruction
theorem of Chapter~\ref{ch:L03} from the generating fields $J^a(z)$ (files \leanpath{WordFields.lean},
\leanpath{Reconstruction.lean}, \leanpath{Assembly.lean}, \leanpath{GradedFields.lean}). The resulting vertex algebra is
\lean{affineVertexAlgebra D k hk}. The conformal vector is the Sugawara state.

\begin{definition}[Sugawara field and conformal vector]\label{l23:def:sugawara}
For $k+h\neq0$,
\begin{equation*}
T(z)=\frac{1}{2(k+h)}\sum_{a,b\in\iota}\kappa^{ab}:\!J^a(z)J^b(z)\!:,\qquad L_n=T_{(n+1)},\qquad \omega=L_{-2}\lxxiiivac,
\end{equation*}
and the central charge is $c=|\iota|\,k/(k+h)$.
\end{definition}
\noindent Files \leanpath{VOA/\allowbreak AffineGeneral/\allowbreak SugawaraField.lean},
\leanpath{VOA/\allowbreak AffineGeneral/\allowbreak Virasoro.lean}.
\begin{leancode}
noncomputable def sugawaraField (k : ℂ) : FormalDist (End ℂ (Vacuum D k)) :=
  (2 * (k + D.h))⁻¹ •
    ∑ a : ι, ∑ b : ι,
      kappaInv D a b • nprod (Jfield_isField D k a) (Jfield_isField D k b)

noncomputable def Lsug (k : ℂ) (n : ℤ) : End ℂ (Vacuum D k) :=
  mode (sugawaraField D k) (n + 1)

noncomputable def omegaA (k : ℂ) : Vacuum D k := Lsug D k (-2) (vacuumA D k)

noncomputable def centralChargeA (k : ℂ) : ℂ := (Fintype.card ι : ℂ) * k / (k + D.h)
\end{leancode}

\begin{dictionary}
\lean{FormalDist (End ℂ V)} is the space of formal distributions $\sum_n A_nz^{-n-1}$ with operator coefficients
(Chapter~\ref{ch:L01}); \lean{Jfield D k a} is $J^a(z)$ and \lean{nprod} is the normally ordered product $:\!ab\!:$
(the $(-1)$-product of fields, Chapter~\ref{ch:L03}); \lean{Jfield_isField} certifies lower truncation, which
\lean{nprod} requires. \lean{mode f n} is $f_{(n)}$, so \lean{Lsug D k n} is $T_{(n+1)}=L_n$ (weight-two shift).
\lean{(2 * (k + D.h))⁻¹} is Lean's inverse, with $0^{-1}=0$: at critical level the definition still typechecks and
gives $T=0$; every theorem about $T$ assumes \lean{k + D.h ≠ 0}. $\sum_{a,b}\kappa^{ab}:\!J^aJ^b\!:=\sum_b:\!J^{b^b}\,J^b\!:$
after expanding the dual vectors, i.e.\ \eqref{l23:eq:sugawara}. Same as the textbook.
\end{dictionary}

\textbf{Why this definition.} The quadratic expression $\sum_{a,b}\kappa^{ab}:\!J^aJ^b\!:$ is the only weight-two field built
from two currents that is invariant under the zero modes (it is the image of the Casimir element
$\sum_ab_ab^a\in U(\lxxiiig)$). Its normalization is fixed by requiring $[L_m,J^c_n]=-nJ^c_{m+n}$, i.e.\ that $L_0$ measure
the weight; the proof of Theorem~\ref{l23:thm:affineVOA} shows that this forces the prefactor $1/(2(k+h))$, not the classical $1/(2k)$.

\begin{physics}
$T(z)$ is the stress tensor of the WZW-type theory, built entirely from the currents; the factor $1/(2(k+h))$
rather than the naive $1/(2k)$ is the quantum correction from normal ordering (the double contraction produces
$h$ via Proposition~\ref{l23:prop:finite}(1)).
\end{physics}

\begin{status}
Definitions (no proof obligation); every theorem about them assumes $k+h\ne0$.
\end{status}

\begin{theorem}[The universal affine VOA]\label{l23:thm:affineVOA}
For Lie data $D$ and $k+h\ne0$, the vacuum module $\lxxiiiV$ with the vertex operators reconstructed from the currents,
vacuum $\lxxiiivac$ and conformal vector $\omega$ is a vertex operator algebra of central charge
$c=\dim\lxxiiig\cdot k/(k+h)$: the modes $L_n$ satisfy
$[L_m,L_n]=(m-n)L_{m+n}+\delta_{m+n,0}\frac{m^3-m}{12}c$; $L_0$ acts on $\lxxiiiV_n=\ker(\ell_0+n)$ as multiplication by $n$
(so $L_0=-\ell_0$ on $\lxxiiiV$: $L_0$ is the energy, $\ell_0$ minus the energy); $L_{-1}$ is translation; the
grading is non-negative and $\lxxiiiV_0=\C\lxxiiivac$.
\end{theorem}
\noindent File \leanpath{VOA/\allowbreak AffineGeneral/\allowbreak Virasoro.lean}.
\begin{leancode}
theorem hYA_virasoro (k : ℂ) (hk : k + D.h ≠ 0) (m n : ℤ) :
    ⁅mode (YA D k hk (omegaA D k)) (m + 1), mode (YA D k hk (omegaA D k)) (n + 1)⁆
      = ((m : ℂ) - n) • mode (YA D k hk (omegaA D k)) (m + n + 1)
        + (if m + n = 0 then
            (((m : ℂ) ^ 3 - m) / 12 * centralChargeA D k) • (1 : End ℂ (Vacuum D k))
           else 0) := by

noncomputable def affineVOA (k : ℂ) (hk : k+D.h ≠ 0) :
    VertexOperatorAlgebra (Vacuum D k) where
  toPreVertexAlgebra := affineVertexAlgebra D k hk
  conformalVector := omegaA D k
  conformal_mem := omegaA_mem D k
  centralCharge := centralChargeA D k
  virasoro := hYA_virasoro D k hk
  L0_grading := hYA_L0 D k hk
  translation := hYA_translation D k hk
  positivity := gradingZA_positivity D k
  vacuum_grade := gradingZA_zero_eq_smul_vacuum D k
\end{leancode}

\begin{dictionary}
The fields of \lean{VertexOperatorAlgebra} are those of Chapter~\ref{ch:L03}: a vertex algebra (grading, vacuum,
$Y$, creation, locality, regularity), a conformal vector in weight $2$, its central charge, the Virasoro relations,
$L_0$-grading, $L_{-1}$-translation, positivity ($V_n=0$ for $n<0$) and $V_0=\C\lxxiiivac$ (CFT type). \lean{YA D k hk u}
is the reconstructed vertex operator $Y(u,z)$; since $\omega$ has weight $2$, $L_n$ is \lean{mode (YA ω) (n+1)}, which
explains the shifts \lean{m + 1}, \lean{n + 1}, \lean{m + n + 1}. \lean{(1 : End ℂ _)} is the identity. The hypothesis
\lean{hk : k + D.h ≠ 0} is non-criticality, the only hypothesis besides the Lie data and the instance argument
\lean{[FiniteDimensional ℂ 𝔤]}, which every file from \leanpath{Currents.lean} to \leanpath{DeligneInstance.lean} declares as a
section variable (it is implied by the finite basis of $D$, so it loses nothing); non-criticality is genuinely needed (at $k=-h$ there is no
Virasoro element of this form). Every VOA axiom is a \emph{field} filled by a theorem; nothing is assumed. The central
charge is \lean{centralChargeA D k}, i.e.\ $|\iota|k/(k+h)$, checked by \lean{rfl} in the instances. Same as the
textbook.

How could it be wrong? A VOA structure on the zero space would satisfy all axioms vacuously; this is excluded by
\lean{vacuumA_ne_zero}. A wrong central charge would make \lean{hYA_virasoro} false, and it is proved.
\end{dictionary}

\textbf{Proof on paper.} This is the classical Frenkel--Zhu construction. 

\emph{(a) Vertex algebra.} The currents $J^a(z)$ are mutually local: their commutator \eqref{l23:eq:affine} is a sum of
$\delta(z-w)$ and $\partial_w\delta(z-w)$ terms, so $(z-w)^2[J^a(z),J^b(w)]=0$. They are translation-covariant, and their modes
generate $\lxxiiiV$ from $\lxxiiivac$. The reconstruction theorem (Chapter~\ref{ch:L03}) gives a unique vertex algebra structure
with $Y(J^a_{-1}\lxxiiivac,z)=J^a(z)$.

\emph{(b) The key commutator $[L_m,J^c_n]=-n\,J^c_{m+n}$.} In modes,
$L_m=\frac{1}{2(k+h)}\sum_{a,b}\kappa^{ab}\sum_{p\in\Z}:\!J^a_pJ^b_{m-p}\!:$, where normal ordering puts the mode with the larger
index on the right (so that on any vector only finitely many terms are nonzero). Commute $J^c_n$ through this sum using
$[J^c_n,J^a_p]=\sum_df^{ca}{}_dJ^d_{n+p}+nk\kappa_{ca}\delta_{n+p,0}$. Two kinds of terms appear.
\begin{itemize}
\item \emph{Central (Schwinger) terms}, one from each of the two factors:
  $nk\sum_{a,b}\kappa^{ab}\big(\kappa_{ca}J^b_{m+n}+\kappa_{cb}J^a_{m+n}\big)=2nk\,J^c_{m+n}$, because
  $\sum_a\kappa_{ca}\kappa^{ab}=\delta_{cb}$ ($\kappa^{-1}\kappa=1$).
\item \emph{Structure-constant terms}, $\sum_{a,b,d}\kappa^{ab}\big(f^{ca}{}_d:\!J^dJ^b\!:+f^{cb}{}_d:\!J^aJ^d\!:\big)$. These would cancel
  exactly, because the tensor $\sum_{a,b}\kappa^{ab}b_a\otimes b_b$ is invariant under $\lxxiiiad(b_c)$ (this is invariance of
  $\kappa$). But after commuting, $|n|$ of the products are no longer normally ordered; restoring the order costs one commutator
  $[J^d_\cdot,J^b_\cdot]$ for each of them. The structure-constant part of these commutators is
  $n\sum_{a,b,d}\kappa^{ab}f^{ca}{}_df^{db}{}_e\,J^e_{m+n}$, which is $n\cdot2h\,J^c_{m+n}$ by Proposition~\ref{l23:prop:finite}(1)
  (the Casimir in coordinates); their central part is proportional to $\sum_af^{ca}{}_a$, a would-be multiple of the identity, which
  vanishes by Proposition~\ref{l23:prop:finite}(2).
\end{itemize}
Altogether $[J^c_n,L_m]=\frac{2nk+2nh}{2(k+h)}J^c_{m+n}=n\,J^c_{m+n}$; this is where the prefactor $1/(2(k+h))$ comes from, and
where $k+h\ne0$ is used. Consequences: since also $L_m\lxxiiivac=0$ for $m\ge-1$, $L_0$ applied to a monomial
$J^{a_1}_{-n_1}\cdots J^{a_r}_{-n_r}\lxxiiivac$ gives $(n_1+\dots+n_r)$ times it, so $L_0$ is the energy and equals $-\ell_0$; and
$L_{-1}$ agrees with the translation operator on the generating fields, hence everywhere.

\emph{(c) Virasoro.} By the Borcherds commutator formula $[L_m,L_n]$ is determined by the products $\omega_{(j)}\omega$, $j=0,1,2,3$:
$\omega_{(0)}\omega=L_{-1}\omega$, $\omega_{(1)}\omega=2\omega$ (by (b), $\omega$ has weight $2$), $\omega_{(2)}\omega=L_1\omega=0$ (by (b),
$L_1J^a_{-1}J^b_{-1}\lxxiiivac=J^a_0J^b_{-1}\lxxiiivac=[b_a,b_b]_{(-1)}\lxxiiivac$, and $\sum_{a,b}\kappa^{ab}[b_a,b_b]=0$ because
$\kappa^{ab}$ is symmetric and the bracket antisymmetric), and $\omega_{(3)}\omega=L_2\omega=\frac c2\lxxiiivac$. The last one is a two-line computation. On the vacuum only
$p=-1$ survives in $L_{-2}$, so $\omega=\frac1{2(k+h)}\sum_{a,b}\kappa^{ab}J^a_{-1}J^b_{-1}\lxxiiivac$. By (b),
$[L_2,J^a_{-1}]=J^a_1$, and $L_2\lxxiiivac=0=J^b_1\lxxiiivac$, so
\begin{equation*}
L_2\,J^a_{-1}J^b_{-1}\lxxiiivac=J^a_1J^b_{-1}\lxxiiivac=[J^a_1,J^b_{-1}]\lxxiiivac
=\Big(\sum_df^{ab}{}_dJ^d_0+k\kappa_{ab}\Big)\lxxiiivac=k\kappa_{ab}\lxxiiivac .
\end{equation*}
Hence $L_2\omega=\frac{k}{2(k+h)}\sum_{a,b}\kappa^{ab}\kappa_{ab}\,\lxxiiivac=\frac{k\dim\lxxiiig}{2(k+h)}\lxxiiivac$ by
Proposition~\ref{l23:prop:finite}(3), i.e.\ $c=k\dim\lxxiiig/(k+h)$. \qed

\begin{example}[$\mathfrak{sl}_2$: the Sugawara construction by hand]\label{l23:ex:sl2sug}
Continue Example~\ref{l23:ex:sl2}: $\kappa=\tr$, $h=2$, basis $e,H,f$, dual family $f,\tfrac12H,e$. Then
\begin{equation*}
L_m=\frac{1}{2(k+2)}\sum_{p\in\Z}:\!\big(J^e_pJ^f_{m-p}+J^f_pJ^e_{m-p}+\tfrac12J^H_pJ^H_{m-p}\big)\!: .
\end{equation*}
\emph{Check of (b) on $v=J^e_{-1}\lxxiiivac$}: we show $L_0v=v$. In $L_0$, the terms with $p\neq0$ that can act nontrivially on the
weight-one vector $v$ are those whose right-hand factor has mode $+1$ (modes $\ge2$ kill weight-one vectors); by the symmetry of the sum
the $p=-1$ and $p=+1$ terms contribute equally, giving
$\frac{2}{2(k+2)}\big(J^e_{-1}J^f_1+J^f_{-1}J^e_1+\tfrac12J^H_{-1}J^H_1\big)v$. Using $[J^x_1,J^e_{-1}]=J^{[x,e]}_0+k\tr(xe)$:
$J^f_1v=k\lxxiiivac$ (since $\tr(fe)=1$ and $J^{-H}_0\lxxiiivac=0$), $J^e_1v=0$ ($\tr(ee)=0$), $J^H_1v=0$ ($\tr(He)=0$). So this part is
$\frac{2k}{2(k+2)}v$: the \emph{Schwinger} contribution $2k$. The $p=0$ term is
$\frac1{2(k+2)}\big(J^e_0J^f_0+J^f_0J^e_0+\tfrac12J^H_0J^H_0\big)v$; zero modes act on weight one by the adjoint action,
$J^x_0J^y_{-1}\lxxiiivac=J^{[x,y]}_{-1}\lxxiiivac$, so this is $\frac1{2(k+2)}J^{\lxxiiiCas(e)}_{-1}\lxxiiivac=\frac{4}{2(k+2)}v$: the
\emph{Casimir} contribution $2h=4$. Total: $L_0v=\frac{2k+4}{2(k+2)}v=v$, as claimed; with the naive prefactor $1/(2k)$ one would get
$\frac{k+2}{k}v$ instead.

\emph{Central charge.} $\omega=\frac1{2(k+2)}\big(J^e_{-1}J^f_{-1}+J^f_{-1}J^e_{-1}+\tfrac12J^H_{-1}J^H_{-1}\big)\lxxiiivac$, and by the
computation in (c), $L_2\omega=\frac{k}{2(k+2)}\big(\tr(ef)+\tr(fe)+\tfrac12\tr(HH)\big)\lxxiiivac=\frac{3k}{2(k+2)}\lxxiiivac$. Hence
$c=\frac{3k}{k+2}$: at $k=1$, $c=1$, the central charge of one free boson, in accordance with the Frenkel--Kac picture
$L_1(\mathfrak{sl}_2)\cong V_{\sqrt2\Z}$ (motivation only, as for $\mathfrak e_8$).
\end{example}

\begin{physics}
This is the chiral algebra of the WZW model of $G$ at level $k$, before any null vectors are removed: the ``universal''
theory whose states are all polynomials in the current modes. The theorem says its stress tensor is the Sugawara tensor
and its central charge is $k\dim G/(k+h)$. For instance, for $\mathfrak e_8$ in the standard normalization ($h=\lxxiiihv=30$)
at level $1$, $c=248\cdot1/31=8$, the central charge of eight free bosons on the $E_8$ torus, as the Frenkel--Kac
isomorphism $L_1(\mathfrak e_8)\cong V_{E_8}$ requires. (This match is arithmetic on paper; the library does not relate the two VOAs.)
\end{physics}

\begin{status}
Proved; \texttt{\#print axioms affineVOA} = \lean{[propext, Classical.choice, Quot.sound]} (re-run for this chapter,
Section~\ref{l23:sec:verification}). \textbf{\texttt{\#print axioms}-clean, conditional on the Lie data} $D$ (which is a
hypothesis about $\lxxiiig$, discharged for every simple $\lxxiiig$ in Section~\ref{l23:sec:everysimple}).
\end{status}

\begin{warning}[title={Trap: The general files keep $\mathfrak{sl}_3$ docstrings}]\label{l23:warn:sl3docs}
The files of \leanpath{VOA/\allowbreak AffineGeneral/} up to \leanpath{Virasoro.lean} were generated from the
$\mathfrak{sl}_3$ files \leanpath{VOA/\allowbreak Examples/\allowbreak AffineSl3*.lean} by a mechanical substitution
(\leanpath{Notes/\allowbreak deligne/\allowbreak campaign/\allowbreak aff\_port.py}: \texttt{Fin 8}$\to\iota$,
$6\to2h$, $3\to h$). The code was ported; the prose was not. Each such file now opens with a disclaimer (``Where the text
below says $\mathfrak{sl}_3$, \texttt{Fin 8}, 6 or 3, read $\lxxiiig$, $\iota$, $2h$ and $h$''), but many docstrings remain
literally false for general $\lxxiiig$: the module title of \leanpath{Virasoro.lean} is ``The actual sl3 Sugawara Virasoro
algebra'' and says ``central charge $8k/(k+3)$''; the docstring of \lean{centralChargeA} says ``using dimension eight and
dual Coxeter number three'' although the code is \lean{(Fintype.card ι : ℂ) * k / (k + D.h)}; the docstring of
\lean{sugawaraField} gives the prefactor ``$1/(2(k+3))$''; \lean{hYA_virasoro}'s says ``$c = 8k/(k+3)$''; \leanpath{WeightOne.lean}
is titled ``The actual weight-one space of the affine sl3 vacuum''. A reader must trust the \emph{code}, which this book
quotes. Listed in Section~\ref{l23:sec:defects}.
\end{warning}

%======================================================================
\section{Every simple Lie algebra}\label{l23:sec:everysimple}
%======================================================================

\subsection{The Casimir is a scalar}

The one non-automatic field of Lie data is the Casimir eigenvalue. For simple $\lxxiiig$ it is supplied by Schur's lemma.

\begin{theorem}[Adjoint Casimir of a simple Lie algebra]\label{l23:thm:casimir}
Let $\lxxiiig$ be a finite-dimensional simple complex Lie algebra, $\kappa$ a symmetric invariant form, $(b_i)$ a basis and
$(b^i)$ a family with $\kappa(b_i,b^j)=\delta_{ij}$. Then there is $c\in\C$ with $\sum_i[b_i,[b^i,x]]=cx$ for all $x$.
\end{theorem}
\noindent File \leanpath{VOA/\allowbreak AffineGeneral/\allowbreak Casimir.lean}.
\begin{leancode}
def casimirOp : Module.End ℂ 𝔤 :=
  ∑ i, (LieModule.toEnd ℂ 𝔤 𝔤 (b i)) * (LieModule.toEnd ℂ 𝔤 𝔤 (d i))

theorem casimir_commute (y x : 𝔤) : casimirOp b d ⁅y, x⁆ = ⁅y, casimirOp b d x⁆ := by

theorem casimir_scalar [FiniteDimensional ℂ 𝔤] [LieAlgebra.IsSimple ℂ 𝔤] :
    ∃ c : ℂ, ∀ x : 𝔤, ∑ i, ⁅b i, ⁅d i, x⁆⁆ = c • x := by
\end{leancode}

\begin{dictionary}
The section variables are \lean{hinv : κ.IsInvariant}, \lean{b : Module.Basis ι ℂ 𝔤}, \lean{d : ι → 𝔤} and
\lean{pair : ∀ i j, κ (b i) (d j) = if i = j then 1 else 0}; they are hypotheses of both theorems (added by
\lean{include hinv pair}). \lean{LieModule.toEnd ℂ 𝔤 𝔤 x} is $\lxxiiiad(x)$. \lean{LieAlgebra.IsSimple ℂ 𝔤} is Mathlib's
simplicity: every ideal is $0$ or $\lxxiiig$, and $\lxxiiig$ is not abelian (so $\lxxiiig\ne0$). Same as the textbook.
\end{dictionary}

\textbf{Proof on paper.} \emph{Commutation}: by the Jacobi identity,
$[y,[b_i,[b^i,x]]]=[[y,b_i],[b^i,x]]+[b_i,[[y,b^i],x]]+[b_i,[b^i,[y,x]]]$; expanding $[y,b_i]=\sum_j\kappa([y,b_i],b^j)b_j$ and
$[y,b^i]=\sum_j\kappa([y,b^i],b_j)b^j$, the first two sums cancel because invariance gives
$\kappa([y,b_i],b^j)+\kappa([y,b^j],b_i)=0$. \emph{Scalar}: $\lxxiiig\ne0$ is finite-dimensional over the algebraically
closed $\C$, so $\lxxiiiCas$ has an eigenvalue $\mu$; by commutation the eigenspace is an ideal; it is nonzero, hence all of
$\lxxiiig$ by simplicity. \qed

\begin{physics}
The adjoint Casimir is the quadratic Casimir of $G$ in the representation carried by the currents themselves; Schur's
lemma (irreducibility of the adjoint of a simple group) makes it a number. Half of it, $h$, is the shift that turns the
classical level $k$ into $k+h$ in the quantum theory. Its value depends on the normalization of $\kappa$: it equals the dual
Coxeter number $\lxxiiihv$ only when $\kappa$ gives long roots squared length $2$, and the theorem here does not compute it.
\end{physics}

\begin{status}
Proved; \texttt{\#print axioms casimir\_scalar} standard (re-run).
\end{status}

\subsection{\texorpdfstring{$V^k(\mathfrak g)$}{V\^{}k(g)} for every simple \texorpdfstring{$\mathfrak g$}{g}}

\begin{theorem}[Universal affine VOA of a simple Lie algebra]\label{l23:thm:simpleaffine}
Let $\lxxiiig$ be a finite-dimensional simple complex Lie algebra.
\begin{enumerate}
\item For every nondegenerate symmetric invariant $\kappa$ there is Lie data indexed by $\{0,\dots,\dim\lxxiiig-1\}$, with
  $h=\frac12\mu$, $\mu$ the Casimir eigenvalue; for every $k$ with $k+h\ne0$ the universal affine VOA exists, with central
  charge $\dim\lxxiiig\cdot k/(k+h)$.
\item The Killing form $\lxxiiiKil(x,y)=\tr(\lxxiiiad x\,\lxxiiiad y)$ is symmetric, invariant and nondegenerate, and in its
  normalization $h=\tfrac12$.
\item For every complex number $\nu\neq0$ the form $\kappa=\lxxiiiKil/(2\nu)$ gives Lie data with Casimir number $h=\nu$, and for
  $k+\nu\neq0$ the VOA $V^k(\lxxiiig)$ in this normalization has central charge $\dim\lxxiiig\cdot k/(k+\nu)$. (Here $\nu$ is a free
  parameter, Lean's \lean{hv}; only for $\nu=\lxxiiihv(\lxxiiig)$ is $\lxxiiiKil/(2\nu)$ the standard normalization, which must be proved
  separately for each type.)
\end{enumerate}
\end{theorem}
\noindent File \leanpath{VOA/\allowbreak AffineGeneral/\allowbreak Simple.lean} (and
\leanpath{VOA/\allowbreak AffineGeneral/\allowbreak Casimir.lean} for \lean{LieData.ofSimple}).
\begin{leancode}
def ofSimple [FiniteDimensional ℂ 𝔤] [LieAlgebra.IsSimple ℂ 𝔤] (hinv : κ.IsInvariant)
    (hnd : LinearMap.BilinForm.Nondegenerate κ) (b : Module.Basis ι ℂ 𝔤) : LieData κ ι :=

def simpleData (κ : LinearMap.BilinForm ℂ 𝔤) (hinv : κ.IsInvariant)
    (hnd : LinearMap.BilinForm.Nondegenerate κ) : LieData κ (Fin (Module.finrank ℂ 𝔤)) :=
  LieData.ofSimple hinv hnd (Module.finBasis ℂ 𝔤)

def universalAffineVOA (κ : LinearMap.BilinForm ℂ 𝔤) (hinv : κ.IsInvariant)
    (hnd : LinearMap.BilinForm.Nondegenerate κ) (k : ℂ) (hk : k + (simpleData 𝔤 κ hinv hnd).h ≠ 0) :
    VertexOperatorAlgebra (AffineLie.vacuumSpace κ hinv k) :=
  affineVOA (simpleData 𝔤 κ hinv hnd) k hk

theorem universalAffineVOA_centralCharge (κ : LinearMap.BilinForm ℂ 𝔤) (hinv : κ.IsInvariant)
    (hnd : LinearMap.BilinForm.Nondegenerate κ) (k : ℂ) (hk : k + (simpleData 𝔤 κ hinv hnd).h ≠ 0) :
    (universalAffineVOA 𝔤 κ hinv hnd k hk).centralCharge =
      (Module.finrank ℂ 𝔤 : ℂ) * k / (k + (simpleData 𝔤 κ hinv hnd).h) := by

theorem killingInvariant : (killingForm ℂ 𝔤).IsInvariant where

def killingAffineVOA (k : ℂ)
    (hk : k + (simpleData 𝔤 (killingForm ℂ 𝔤) (killingInvariant 𝔤)
      (LieAlgebra.IsKilling.killingForm_nondegenerate ℂ 𝔤)).h ≠ 0) :
    VertexOperatorAlgebra (AffineLie.vacuumSpace (killingForm ℂ 𝔤) (killingInvariant 𝔤) k) :=

theorem killing_h :
    (simpleData 𝔤 (killingForm ℂ 𝔤) (killingInvariant 𝔤)
      (LieAlgebra.IsKilling.killingForm_nondegenerate ℂ 𝔤)).h = 1 / 2 := by

def standardData (hv : ℂ) (hhv : hv ≠ 0) :
    LieData ((2 * hv)⁻¹ • killingForm ℂ 𝔤) (Fin (Module.finrank ℂ 𝔤)) :=

theorem standardData_h (hv : ℂ) (hhv : hv ≠ 0) :
    (standardData 𝔤 hv hhv).h = hv := by

def standardAffineVOA (hv : ℂ) (hhv : hv ≠ 0) (k : ℂ) (hk : k + hv ≠ 0) :
    VertexOperatorAlgebra (Vacuum (standardData 𝔤 hv hhv) k) :=
  affineVOA (standardData 𝔤 hv hhv) k (by rw [standardData_h]; exact hk)

theorem standardAffineVOA_centralCharge (hv : ℂ) (hhv : hv ≠ 0) (k : ℂ)
    (hk : k + hv ≠ 0) :
    (standardAffineVOA 𝔤 hv hhv k hk).centralCharge = (Module.finrank ℂ 𝔤 : ℂ) * k / (k + hv) := by
\end{leancode}

\begin{dictionary}
The section variables are \lean{(𝔤 : Type*) [LieRing 𝔤] [LieAlgebra ℂ 𝔤] [FiniteDimensional ℂ 𝔤] [LieAlgebra.IsSimple ℂ 𝔤]}:
every statement assumes $\lxxiiig$ finite-dimensional and simple, and nothing else. \lean{Module.finBasis ℂ 𝔤} is a
basis indexed by \lean{Fin (finrank ℂ 𝔤)} chosen by Mathlib (noncomputably); \lean{Module.finrank} is the dimension.
\lean{killingForm ℂ 𝔤} is Mathlib's Killing form; its nondegeneracy \lean{LieAlgebra.IsKilling.killingForm_nondegenerate}
is available because Mathlib derives \lean{IsKilling} from simplicity via Cartan's criterion
(\lean{HasTrivialRadical.instIsKilling}), so no separate hypothesis is needed.

\emph{The $h$ of \lean{simpleData} is not computed}: it is half of an eigenvalue chosen by \lean{Classical.choose} from
\lean{casimir_scalar}. Hence the non-criticality hypothesis of \lean{universalAffineVOA} is stated as
\lean{k + (simpleData 𝔤 κ hinv hnd).h ≠ 0}, which a user can only check by computing that $h$; \lean{killing_h} does
this for the Killing form ($h=\frac12$), and \lean{standardData_h} for its rescalings.

\emph{The word ``standard''.} Write $\nu$ for Lean's \lean{hv}. \lean{standardData 𝔤 hv hhv} is Lie data for $\lxxiiiKil/(2\nu)$ and
\lean{standardData_h} says its Casimir number equals $\nu$ --- for \emph{every} nonzero $\nu$. This is a tautology of rescaling (the
Killing normalization has $h=\frac12$, and multiplying the form by $(2\nu)^{-1}$ divides $h$ by $(2\nu)^{-1}$, giving $\nu$), not a
computation of the dual Coxeter number. The form $\lxxiiiKil/(2\nu)$ is the standard normalization (long roots of squared length $2$)
if and only if $\nu=\lxxiiihv(\lxxiiig)$; for the types used later this is proved separately (Theorem~\ref{l23:thm:normalization} and the
\lean{*_standard} theorems of Section~\ref{l23:sec:sc}). Thus \lean{standardAffineVOA 𝔤 hv _ k _} is \emph{the} $V^k(\lxxiiig)$ of
the literature only once $\nu=\lxxiiihv$ is established. For $\mathfrak{sl}_n$ this identification is not in Lean (see the warning box in Section~\ref{l23:warn:slnnorm}).
Otherwise same as the textbook.
\end{dictionary}

\textbf{Proof on paper.} (1) Theorem~\ref{l23:thm:casimir} provides the Casimir eigenvalue; Mathlib's \lean{dualBasis} provides
the dual family; Theorem~\ref{l23:thm:affineVOA} does the rest. (2) Symmetry and invariance of the trace form are Mathlib
facts. For $h$: $\tr\lxxiiiCas=\sum_i\tr(\lxxiiiad b_i\,\lxxiiiad b^i)=\sum_i\lxxiiiKil(b_i,b^i)=\dim\lxxiiig$, while
$\lxxiiiCas=2h\cdot\mathrm{id}$ has trace $2h\dim\lxxiiig$, so $h=\frac12$. (3) Rescaling (Definition~\ref{l23:def:smulform}):
$h=\tfrac12\cdot 2\nu=\nu$. Central charges: $|\mathrm{Fin}\,n|=n$. \qed

\begin{physics}
Every simple compact gauge group $G$ has a WZW current algebra, and this theorem says the library has its chiral algebra
for every level except the critical one. In the Killing normalization the critical level is $k=-\frac12$; in the
normalization physicists use (long roots of length $2$, $\lxxiiihv$ the dual Coxeter number) it is $k=-\lxxiiihv$, and
integrable levels are positive integers.
\end{physics}

\begin{status}
Proved; \texttt{\#print axioms} standard for \lean{universalAffineVOA}, \lean{universalAffineVOA_centralCharge},
\lean{killing_h}, \lean{standardAffineVOA_centralCharge} (re-run, Section~\ref{l23:sec:verification}).
\textbf{Unconditional} in the sense of this book: the only hypotheses are simplicity, finite dimension, nondegeneracy of the
chosen form and non-criticality, all part of the intended statement. \emph{What is proved}: the \emph{universal} VOA $V^k(\lxxiiig)$
(not $L_k(\lxxiiig)$), with central charge $\dim\lxxiiig\cdot k/(k+h)$ where $2h$ is the adjoint Casimir eigenvalue of the chosen form,
chosen by \lean{Classical.choose}. \emph{Not proved} for general $\lxxiiig$: that $h=\lxxiiihv$; in \lean{standardAffineVOA} the number
$h$ is the free parameter \lean{hv}, and $h=\lxxiiihv$ holds only per type through the \lean{*_standard} theorems.
\end{status}

\begin{warning}[title={Naming: \texttt{universalAffineVOA} is $V^k(\mathfrak g)$, \texttt{simpleAffineQuotientVOA} is $L_k(\mathfrak g)$}]\label{l23:warn:simplename}
The type of \lean{universalAffineVOA} is \lean{VertexOperatorAlgebra (AffineLie.vacuumSpace κ hinv k)}: the carrier is the full vacuum
module, so this is $\lxxiiiV$; the hypothesis \lean{LieAlgebra.IsSimple ℂ 𝔤} concerns the Lie algebra. Until 28 October 2026 the
declaration was called \lean{simpleAffineVOA}, with ``simple'' referring to $\lxxiiig$, and the release audit flagged the name as
misleading; it was renamed, and \lean{simpleAffineVOA} survives only as a deprecated alias of \lean{universalAffineVOA}. Its docstring
says what it is (``universal affine VOA \dots\ of a simple Lie algebra''). The simple quotient $\lxxiiiL$ of an
arbitrary simple $\lxxiiig$ is obtained by applying Definition~\ref{l23:def:simpleVOA} to \lean{simpleData}, i.e.\ the term
\lean{simpleVOA (simpleData 𝔤 κ hinv hnd) k hk}. At the time of the release audit the library had no named declaration for it (the
audit's search for \lean{simpleVOA (simpleData} found nothing); since 2026-10-28 it is named \lean{simpleAffineQuotientVOA}
(Definition~\ref{l23:def:simpleaffinequotient}), whose module docstring records the former name. The naming now matches the
mathematics: \lean{universalAffineVOA} is $\lxxiiiV$ and \lean{simpleAffineQuotientVOA} is $\lxxiiiL$. Anyone reading older code or
notes should know that ``\lean{simpleAffineVOA}'' there means $\lxxiiiV$, not $L_k(\lxxiiig)$. Two project
notes made the slip (a layer-one summary, and a progress note speaking of ``universal and simple affine VOAs''); the
correct reading is ``$\lxxiiiV$ (\lean{universalAffineVOA}) and its simple quotient $\lxxiiiL$ (\lean{simpleAffineQuotientVOA}, by
definition \lean{simpleVOA} applied to \lean{simpleData})''. Source: release audit 2026-10-27, finding M1; the missing name is
resolved by \leanpath{VOA/\allowbreak AffineGeneral/\allowbreak SimpleAffineQuotient.lean}, the misleading old name by the renaming.
\end{warning}

\subsection{The \texorpdfstring{$\mathfrak{sl}_3$}{sl(3)} instance}\label{l23:sec:sl3}

As a consistency check the bespoke $\mathfrak{sl}_3$ data of \leanpath{VOA/\allowbreak Examples/\allowbreak AffineSl3*.lean} are packaged as Lie data.

\begin{proposition}[$\mathfrak{sl}_3$ Lie data]\label{l23:prop:sl3}
The trace form, the Chevalley basis, its trace-dual family and $h=3$ are Lie data for $\mathfrak{sl}_3$, and the resulting
$V^k(\mathfrak{sl}_3)$ has central charge $8k/(k+3)$.
\end{proposition}
\noindent File \leanpath{VOA/\allowbreak AffineGeneral/\allowbreak Sl3Instance.lean}.
\begin{leancode}
def sl3Data : LieData AffineSl3.form (Fin 8) where
  invariant := form_invariant
  basis := chevalleyBasis
  dual := dualGenerator
  pair i j := by rw [chevalleyBasis_apply]; exact generator_dual i j
  h := 3

theorem sl3VOA_centralCharge (k : ℂ) (hk : k + 3 ≠ 0) :
    (sl3VOA k hk).centralCharge = 8 * k / (k + 3) := by
\end{leancode}

\begin{dictionary}
\lean{AffineSl3.form}, \lean{chevalleyBasis}, \lean{dualGenerator} and \lean{adjoint_casimir} (used in the omitted
\lean{casimir} field) are the Lie-theoretic objects of the bespoke $\mathfrak{sl}_3$ files. \lean{sl3VOA k hk} abbreviates
\lean{affineVOA sl3Data k hk}. The carrier is the same \lean{vacuumSpace} as the bespoke \lean{affineSl3VOA} of
\leanpath{VOA/\allowbreak Examples/\allowbreak AffineSl3Virasoro.lean}, but the two VOA structures are \emph{not} proved equal in Lean. Same as the textbook.
\end{dictionary}

\textbf{Proof on paper.} The fields are the $\mathfrak{sl}_3$ facts proved in the bespoke files; $\lxxiiiCas=6=2\cdot3$. \qed

\begin{physics}
$V^k(\mathfrak{sl}_3)$, the current algebra of $SU(3)$; at $k=1$, $c=2$, the two bosons of the $A_2$ torus.
\end{physics}

\begin{status}
Proved; axioms standard (project record, Section~\ref{l23:sec:verification}).
\end{status}

%======================================================================
\section{The simple quotient \texorpdfstring{$L_k(\mathfrak g)$}{L\_k(g)} and its weight-one space}\label{l23:sec:simple}
%======================================================================

\subsection{The simple quotient}

Chapter~\ref{ch:L04} (Section~\ref{l04:sec:quot}) defines, for any VOA $W$, the maximal graded ideal
\lean{maxGradedIdeal W} (the supremum of all proper graded ideals), the quotient \lean{maxQuotientVOA W}, and proves that
the quotient is simple among \emph{all} ideals as soon as the vacuum of $W$ is nonzero
(\lean{isSimple_maxQuotientVOA'}). The general construction applies this to $\lxxiiiV$.

\begin{definition}[The simple affine VOA]\label{l23:def:simpleVOA}
For Lie data $D$ and $k+h\ne0$, $\lxxiiiL=\lxxiiiV/I_{\max}$, where $I_{\max}$ is the maximal graded ideal; the
\emph{current states} of $\lxxiiiL$ are the images $[x_{(-1)}\lxxiiivac]$, $x\in\lxxiiig$.
\end{definition}
\noindent File \leanpath{VOA/\allowbreak AffineGeneral/\allowbreak SimpleQuotient.lean}.
\begin{leancode}
def maxIdeal (k : ℂ) (hk : k + D.h ≠ 0) : Submodule ℂ (Vacuum D k) :=
  QuotVOA.maxGradedIdeal (affineVOA D k hk)

def simpleVOA (k : ℂ) (hk : k + D.h ≠ 0) : VertexOperatorAlgebra (Vacuum D k ⧸ maxIdeal D k hk) :=
  QuotVOA.maxQuotientVOA (affineVOA D k hk)

def simpleProjection (k : ℂ) (hk : k + D.h ≠ 0) : Vacuum D k →ₗ[ℂ] Vacuum D k ⧸ maxIdeal D k hk :=
  (maxIdeal D k hk).mkQ

def simpleCurrentState (k : ℂ) (hk : k + D.h ≠ 0) : 𝔤 →ₗ[ℂ] Vacuum D k ⧸ maxIdeal D k hk :=
  (simpleProjection D k hk).comp (currentState D k)
\end{leancode}

\begin{dictionary}
\lean{⧸} is the quotient vector space, \lean{mkQ} the quotient map. \lean{currentState D k x} is
$x_{(-1)}\lxxiiivac$, the action of $x\otimes t^{-1}$ on the vacuum (Definition~\ref{l23:def:currentstate} below).
\lean{→ₗ[ℂ]} is a $\C$-linear map. Same as the textbook, with one remark: the literature's $L_k(\lxxiiig)$ is the quotient
by the unique maximal proper ideal; since $\lxxiiiV$ is of CFT type every ideal not containing $\lxxiiivac$ is graded (by
$L_0$-stability), so the maximal graded ideal is that ideal (Chapter~\ref{ch:L04}).
\end{dictionary}

\textbf{Why this definition.} Quotienting by the \emph{maximal} graded ideal is the only choice that yields a simple VOA
without knowing the ideal explicitly; this is what makes $L_k(\lxxiiig)$ available for every $k$ even though its null vectors are
known only at special levels. The price is that the definition displays nothing about $I_{\max}$; facts about it (such as
Theorem~\ref{l23:thm:singular}) must be proved separately.

\begin{physics}
$\lxxiiiL$ is the chiral algebra in which all null states (states orthogonal to everything, equivalently generating proper
ideals) are set to zero: the physical theory. At positive integer level it is the WZW chiral algebra; at the Deligne levels it is
(by the physics literature quoted in Section~\ref{l23:sec:physics}) the four-dimensional $\mathcal N=2$ chiral algebra of
Chapters~\ref{ch:L24}--\ref{ch:L27}.
\end{physics}

\begin{status}
Definitions (no proof obligation); \lean{simpleVOA} is a VOA because \lean{maxQuotientVOA} is (Chapter~\ref{ch:L04}). For an
arbitrary finite-dimensional simple $\lxxiiig$ with form $\kappa$, $\lxxiiiL$ is the term
\lean{simpleVOA (simpleData 𝔤 κ hinv hnd) k hk}, named \lean{simpleAffineQuotientVOA} (Section~\ref{l23:sec:simplename}); it is \emph{not} \lean{universalAffineVOA}, which is $\lxxiiiV$ (warning box in
Section~\ref{l23:warn:simplename}).
\end{status}

\begin{theorem}[Simplicity, and weight one is spanned by the currents]\label{l23:thm:simple}
For Lie data $D$ and $k+h\ne0$: $\lxxiiiL$ is a simple VOA, and its weight-one space is the image of
$x\mapsto[x_{(-1)}\lxxiiivac]$; in particular each basis current lies in weight one.
\end{theorem}
\noindent File \leanpath{VOA/\allowbreak AffineGeneral/\allowbreak SimpleQuotient.lean}.
\begin{leancode}
theorem simpleVOA_isSimple (k : ℂ) (hk : k + D.h ≠ 0) : (simpleVOA D k hk).IsSimple :=
  QuotVOA.isSimple_maxQuotientVOA' (affineVOA D k hk) (vacuumA_ne_zero D k)

theorem simple_grading_one_eq_range (k : ℂ) (hk : k + D.h ≠ 0) :
    (simpleVOA D k hk).grading 1 = LinearMap.range (simpleCurrentState D k hk) := by

theorem simpleCurrentState_mem (k : ℂ) (hk : k + D.h ≠ 0) (i : ι) :
    simpleCurrentState D k hk (D.generator i) ∈ (simpleVOA D k hk).grading 1 := by
\end{leancode}

\begin{dictionary}
\lean{IsSimple} is the definition of Chapter~\ref{ch:L04}: the carrier is nonzero and every ideal is $0$ or everything.
\lean{LinearMap.range} is the image. The proof term of \lean{simpleVOA_isSimple} is displayed because it is informative: the
only input beyond the general theory of Chapter~\ref{ch:L04} is \lean{vacuumA_ne_zero} (relative PBW). Without it, the
``simple quotient'' of a zero VOA would be zero and \lean{IsSimple} would fail; so the theorem is not vacuous.
Same as the textbook.
\end{dictionary}

\textbf{Proof on paper.} Simplicity is the general fact of Chapter~\ref{ch:L04}. Weight one: the vacuum module is induced,
$\lxxiiiV=U(\widehat{\lxxiiig}_-)\lxxiiivac$ with $\widehat{\lxxiiig}_-=\lxxiiig\otimes t^{-1}\C[t^{-1}]$, so by the PBW theorem it is
spanned by the monomials $J^{a_1}_{-n_1}\cdots J^{a_r}_{-n_r}\lxxiiivac$ ($n_i\ge1$), of weight $n_1+\dots+n_r$. A monomial of weight one
has exactly one letter, of mode $-1$; so $\lxxiiiV_1$ is spanned by the $x_{(-1)}\lxxiiivac$. (In Lean this ``creation-word spanning''
is \lean{grading_one_eq_range}, Theorem~\ref{l23:thm:weightone}.) The quotient map preserves the grading and is surjective in each
degree. \qed

\begin{physics}
The weight-one states of a chiral algebra are its conserved currents, i.e.\ its continuous symmetries. The theorem says
$\lxxiiiL$ has no currents other than those of $\lxxiiig$: quotienting cannot create new symmetries.
\end{physics}

\begin{status}
Proved; axioms standard (re-run).
\end{status}

\subsection{The named simple affine VOA \texorpdfstring{$L_k(\mathfrak g)$}{L\_k(g)} of a simple Lie algebra}\label{l23:sec:simplename}

For Lie data built from a simple $\lxxiiig$ by \lean{simpleData} (Theorem~\ref{l23:thm:simpleaffine}), the simple quotient of
Definition~\ref{l23:def:simpleVOA} is the term \lean{simpleVOA (simpleData 𝔤 κ hinv hnd) k hk}. Since 2026-10-28 the library gives it a name of its own and transports the facts just proved to it. Nothing new is proved here
mathematically; the point is that a reader (or a later chapter) can now write ``$\lxxiiiL$'' as a single Lean name, and the name
says what it is. Together with the renaming recorded in Section~\ref{l23:warn:simplename}, the two names now say what they are:
\lean{universalAffineVOA} is the \emph{universal} algebra $\lxxiiiV$, \lean{simpleAffineQuotientVOA} its simple quotient.

\begin{definition}[The simple affine VOA $L_k(\mathfrak g)$, named]\label{l23:def:simpleaffinequotient}
Let $\lxxiiig$ be a finite-dimensional simple complex Lie algebra, $\kappa$ a nondegenerate invariant bilinear form on it, $h$ the
Casimir number of $\kappa$ (half the eigenvalue of the adjoint Casimir built from $\kappa$), and $k\in\C$ with $k+h\ne0$. Then
\begin{equation*}
L_k(\lxxiiig)\;=\;V^k(\lxxiiig)/I_{\max},
\end{equation*}
the quotient of the universal affine VOA by its maximal graded ideal, with the induced VOA structure, and
$\pi:V^k(\lxxiiig)\to L_k(\lxxiiig)$ denotes the quotient map.
\end{definition}
\noindent File \leanpath{VOA/\allowbreak AffineGeneral/\allowbreak SimpleAffineQuotient.lean}.
\begin{leancode}
def simpleAffineQuotientVOA (κ : LinearMap.BilinForm ℂ 𝔤) (hinv : κ.IsInvariant)
    (hnd : LinearMap.BilinForm.Nondegenerate κ) (k : ℂ) (hk : k + (simpleData 𝔤 κ hinv hnd).h ≠ 0) :
    VertexOperatorAlgebra
      (Vacuum (simpleData 𝔤 κ hinv hnd) k ⧸ maxIdeal (simpleData 𝔤 κ hinv hnd) k hk) :=
  simpleVOA (simpleData 𝔤 κ hinv hnd) k hk

def simpleAffineQuotientProjection (κ : LinearMap.BilinForm ℂ 𝔤) (hinv : κ.IsInvariant)
    (hnd : LinearMap.BilinForm.Nondegenerate κ) (k : ℂ) (hk : k + (simpleData 𝔤 κ hinv hnd).h ≠ 0) :
    Vacuum (simpleData 𝔤 κ hinv hnd) k →ₗ[ℂ]
      Vacuum (simpleData 𝔤 κ hinv hnd) k ⧸ maxIdeal (simpleData 𝔤 κ hinv hnd) k hk :=
  simpleProjection (simpleData 𝔤 κ hinv hnd) k hk
\end{leancode}

\begin{dictionary}
The section variables are those of Theorem~\ref{l23:thm:simpleaffine}: \lean{(𝔤 : Type*) [LieRing 𝔤] [LieAlgebra ℂ 𝔤]
[FiniteDimensional ℂ 𝔤] [LieAlgebra.IsSimple ℂ 𝔤]}. The explicit arguments are the form $\kappa$ (\lean{LinearMap.BilinForm ℂ 𝔤}, a
bilinear form), its invariance \lean{hinv} ($\kappa([x,y],z)=\kappa(x,[y,z])$), its nondegeneracy \lean{hnd}, the level $k$, and the
non-criticality hypothesis \lean{hk}: $k+h\ne0$ with $h$ the Casimir number of \lean{simpleData}, exactly the hypothesis of
\lean{universalAffineVOA}. (Symmetry of $\kappa$ is not among the explicit hypotheses.) The carrier
\lean{Vacuum (simpleData 𝔤 κ hinv hnd) k ⧸ maxIdeal (simpleData 𝔤 κ hinv hnd) k hk} is the quotient vector space of the vacuum
module by the maximal graded ideal of Definition~\ref{l23:def:simpleVOA}; \lean{Vacuum D k} is an \lean{abbrev} for
\lean{vacuumSpace κ D.invariant k} (file \leanpath{VOA/\allowbreak AffineGeneral/\allowbreak Currents.lean}), the carrier of
\lean{universalAffineVOA}.

Both definitions are one-line abbreviations, and the one-line bodies are shown because they \emph{are} the definition:
\lean{simpleAffineQuotientVOA} is literally \lean{simpleVOA} (Definition~\ref{l23:def:simpleVOA}) applied to \lean{simpleData}, and
\lean{simpleAffineQuotientProjection} is literally \lean{simpleProjection}, i.e.\ \lean{mkQ}, the quotient map. So the named object is
the term the earlier status boxes described; no new choice is made. The docstring of the module says that \lean{universalAffineVOA} is the universal algebra and records its former name. As in \lean{universalAffineVOA}, the Casimir number $h$ is the one chosen
by \lean{Classical.choose}; it equals $\lxxiiihv$ only after the per-type normalization theorems (Section~\ref{l23:sec:sc}).
Same as the textbook, with the normalization caveat of Theorem~\ref{l23:thm:simpleaffine}.
\end{dictionary}

\textbf{Why it could have been wrong, and is not.} A wrapper with this name could have been a definition of something else (for
instance the universal algebra again, which is what the former name \lean{simpleAffineVOA} suggested to readers). The body \lean{simpleVOA (simpleData …) k hk} and
the carrier type, a quotient by \lean{maxIdeal}, exclude this. A quotient by the maximal graded ideal could also be degenerate
(zero) if the vacuum were zero; Theorem~\ref{l23:thm:simpleaffinequotient} below records simplicity, which includes nonvanishing.

\begin{physics}
$L_k(\lxxiiig)$ is the physical current algebra (the chiral algebra of the WZW model at positive integer level, of the
four-dimensional $\mathcal N=2$ theories at the Deligne levels), in which every null state has been set to zero. The projection
$\pi$ is the statement ``compute in the universal algebra, then impose the null relations'': physicists do exactly this when they
write correlators of currents first and then remove null descendants.
\end{physics}

\begin{status}
Definitions (no proof obligation). \texttt{\#print axioms} for both reports \lean{[propext, Classical.choice, Quot.sound]} (re-run for
this chapter, Section~\ref{l23:sec:verification}). The file is imported by the umbrella through
\leanpath{VOA/\allowbreak AffineGeneral/\allowbreak All.lean}; the targeted and umbrella builds are recorded green with standard axioms
(\texttt{Notes/RELEASE\_PLAN.md}, ``Caveats removed in Lean'', 2026-10-28); the aggregate build was green in pre-release builds on 2026-10-28
(18{,}901 jobs; 18{,}907 after a merge); for the release build see Chapter~\ref{ch:15}, Section~\ref{ch15:sec:record}. Not re-run for this chapter.
\end{status}

\begin{theorem}[Properties of the named $L_k(\mathfrak g)$]\label{l23:thm:simpleaffinequotient}
In the setting of Definition~\ref{l23:def:simpleaffinequotient}:
\begin{enumerate}
\item $L_k(\lxxiiig)$ is a simple VOA (nonzero, and every ideal is $0$ or everything);
\item its central charge is $\dim\lxxiiig\cdot k/(k+h)$, the same as that of $V^k(\lxxiiig)$;
\item $\pi$ is a homomorphism for all products: $\pi(u_{(n)}v)=\pi(u)_{(n)}\pi(v)$ for all $u,v\in V^k(\lxxiiig)$ and $n\in\Z$;
\item $\pi$ is surjective;
\item $\pi$ preserves conformal weight: if $x\in V^k(\lxxiiig)_n$ then $\pi(x)\in L_k(\lxxiiig)_n$, for every $n\in\Z$.
\end{enumerate}
\end{theorem}
\noindent File \leanpath{VOA/\allowbreak AffineGeneral/\allowbreak SimpleAffineQuotient.lean}.
\begin{leancode}
theorem simpleAffineQuotientVOA_isSimple (κ : LinearMap.BilinForm ℂ 𝔤) (hinv : κ.IsInvariant)
    (hnd : LinearMap.BilinForm.Nondegenerate κ) (k : ℂ) (hk : k + (simpleData 𝔤 κ hinv hnd).h ≠ 0) :
    (simpleAffineQuotientVOA 𝔤 κ hinv hnd k hk).IsSimple :=
  simpleVOA_isSimple (simpleData 𝔤 κ hinv hnd) k hk

theorem simpleAffineQuotientVOA_centralCharge (κ : LinearMap.BilinForm ℂ 𝔤) (hinv : κ.IsInvariant)
    (hnd : LinearMap.BilinForm.Nondegenerate κ) (k : ℂ) (hk : k + (simpleData 𝔤 κ hinv hnd).h ≠ 0) :
    (simpleAffineQuotientVOA 𝔤 κ hinv hnd k hk).centralCharge =
      (Module.finrank ℂ 𝔤 : ℂ) * k / (k + (simpleData 𝔤 κ hinv hnd).h) :=
  universalAffineVOA_centralCharge 𝔤 κ hinv hnd k hk

theorem simpleAffineQuotientProjection_smul (κ : LinearMap.BilinForm ℂ 𝔤) (hinv : κ.IsInvariant)
    (hnd : LinearMap.BilinForm.Nondegenerate κ) (k : ℂ) (hk : k + (simpleData 𝔤 κ hinv hnd).h ≠ 0)
    (u : Vacuum (simpleData 𝔤 κ hinv hnd) k) (n : ℤ) (v : Vacuum (simpleData 𝔤 κ hinv hnd) k) :
    (simpleAffineQuotientVOA 𝔤 κ hinv hnd k hk).smul
        (simpleAffineQuotientProjection 𝔤 κ hinv hnd k hk u) n
        (simpleAffineQuotientProjection 𝔤 κ hinv hnd k hk v) =
      simpleAffineQuotientProjection 𝔤 κ hinv hnd k hk
        ((universalAffineVOA 𝔤 κ hinv hnd k hk).smul u n v) :=

theorem simpleAffineQuotientProjection_surjective (κ : LinearMap.BilinForm ℂ 𝔤)
    (hinv : κ.IsInvariant) (hnd : LinearMap.BilinForm.Nondegenerate κ) (k : ℂ)
    (hk : k + (simpleData 𝔤 κ hinv hnd).h ≠ 0) :
    Function.Surjective (simpleAffineQuotientProjection 𝔤 κ hinv hnd k hk) :=

theorem simpleAffineQuotientProjection_grading (κ : LinearMap.BilinForm ℂ 𝔤)
    (hinv : κ.IsInvariant) (hnd : LinearMap.BilinForm.Nondegenerate κ) (k : ℂ)
    (hk : k + (simpleData 𝔤 κ hinv hnd).h ≠ 0) {n : ℤ} {x : Vacuum (simpleData 𝔤 κ hinv hnd) k}
    (hx : x ∈ (universalAffineVOA 𝔤 κ hinv hnd k hk).grading n) :
    simpleAffineQuotientProjection 𝔤 κ hinv hnd k hk x ∈
      (simpleAffineQuotientVOA 𝔤 κ hinv hnd k hk).grading n :=
\end{leancode}

\begin{dictionary}
\emph{Products.} \lean{W.smul u n v} is \lean{mode (W.Y u) n v}, the $n$-th product $u_{(n)}v$ of the VOA \lean{W} (the coefficient of
$z^{-n-1}$ in $Y(u,z)v$, Chapter~\ref{ch:L01}), for every integer $n$, negative ones included. In (3) the left side is the $n$-th product of $L_k(\lxxiiig)$
applied to the images $\pi(u),\pi(v)$; the right side is $\pi$ of the $n$-th product computed in \lean{universalAffineVOA}, i.e.\ in
$V^k(\lxxiiig)$. Since all $n$ are covered, (3) says that
$\pi$ intertwines the full vertex operators: $\pi(Y(u,z)v)=Y(\pi u,z)\pi v$.

\emph{Grading.} \lean{W.grading n} is the weight-$n$ subspace of \lean{W}, on which $L_0$ acts by $n$ (field \lean{L0_grading}, Chapter~\ref{ch:L01}); the weight
index is an integer, and \lean{{n : ℤ}} and \lean{{x : …}} are implicit arguments inferred from \lean{hx}. (5) is the inclusion
$\pi(V^k_n)\subseteq (L_k)_n$; together with (4) and the directness of both gradings this gives $\pi(V^k_n)=(L_k)_n$ on paper, but
the equality is not stated in Lean.

\emph{Central charge.} \lean{centralCharge} is the field $c$ of the VOA structure appearing in its Virasoro relation
$[L_m,L_n]=(m-n)L_{m+n}+\frac{m^3-m}{12}c\,\delta_{m+n,0}$ (Chapter~\ref{ch:L01}).
\lean{Module.finrank ℂ 𝔤} is $\dim\lxxiiig$, cast to $\C$. The right side is literally the right side of
\lean{universalAffineVOA_centralCharge}, and the displayed proof term shows the equality is \emph{definitional transport}: the quotient
VOA carries the image of the conformal vector, and its \lean{centralCharge} field is copied from the parent
(\lean{centralCharge := W.centralCharge} in \lean{QuotVOA.quotientVOA}, Chapter~\ref{ch:L04}); the content is that the Virasoro relation
holds in the quotient with the same $c$.

\emph{Simplicity.} \lean{IsSimple} is the definition of Chapter~\ref{ch:L04} (nonzero carrier, only trivial ideals among \emph{all}
ideals, not only graded ones); the proof term is \lean{simpleVOA_isSimple}, Theorem~\ref{l23:thm:simple}.

\emph{Hypotheses.} Every statement carries exactly $\kappa$, \lean{hinv}, \lean{hnd}, $k$, \lean{hk} and the section variables;
there is no hidden assumption. Same as the textbook, in the normalization of $\kappa$ (so $h=\lxxiiihv$ only per type).
\end{dictionary}

\textbf{Proof on paper.} (1) is Theorem~\ref{l23:thm:simple} for $D=$\lean{simpleData}. (2) The conformal vector of the quotient is
the class $[\omega]$ of the Sugawara vector, and by (3) its modes are the modes $L_n$ of $V^k$ acting on classes; the Virasoro relation
of $V^k$ therefore descends to the quotient with the same $c$ (and the scalar is detected, since the quotient is nonzero by (1)). The
central charge is unchanged. (3) $I_{\max}$ is a two-sided
ideal, so $u_{(n)}v$ modulo $I_{\max}$ depends only on the classes of $u$ and $v$; the quotient VOA is \emph{defined} by this formula
(Chapter~\ref{ch:L04}). (4) Every quotient map is surjective. (5) $I_{\max}$ is graded, so the quotient grading is the image of the
grading of $V^k$. \qed

\begin{physics}
(1) says $L_k(\lxxiiig)$ has no proper consistent truncation: it is the chiral algebra of a theory, not of a reducible collection of
sectors. (2) Imposing null relations does not change the Virasoro central charge --- the Sugawara stress tensor of the currents is the
same field before and after; the familiar $c=\dim\lxxiiig\,k/(k+\lxxiiihv)$ of the WZW model is the value. (3)--(5) say that operator
product expansions and conformal weights computed in the universal algebra descend to the physical one, and that every state of the
physical theory comes from a polynomial in current modes acting on the vacuum.
\end{physics}

\begin{status}
Proved; \texttt{\#print axioms} reports \lean{[propext, Classical.choice, Quot.sound]} for all five theorems (re-run for this
chapter, Section~\ref{l23:sec:verification}). \textbf{Unconditional} in the sense of this book: the only hypotheses are those of the
intended statement. Not stated in Lean: equality of weight spaces in (5); that $h=\lxxiiihv$ for general $\lxxiiig$.
\end{status}

\subsection{Weight one of the universal algebra}

\begin{definition}[Current state]\label{l23:def:currentstate}
$x\mapsto x_{(-1)}\lxxiiivac$, the linear map $\lxxiiig\to\lxxiiiV$.
\end{definition}
\noindent File \leanpath{VOA/\allowbreak AffineGeneral/\allowbreak WeightOne.lean}.
\begin{leancode}
def currentState (k : ℂ) : 𝔤 →ₗ[ℂ] Vacuum D k where
  toFun x := vacuumRho κ D.invariant k (ofLoop κ D.invariant x (-1)) (vacuumA D k)
\end{leancode}

\begin{dictionary}
The action of the loop element $x\otimes t^{-1}$ on the vacuum. Linearity is proved in the omitted fields. Same as the textbook.
\end{dictionary}

\begin{physics}
The state--field correspondence sends $x_{(-1)}\lxxiiivac$ to the current $J^x(z)$.
\end{physics}

\begin{status}
A definition (no proof obligation).
\end{status}

\begin{theorem}[Weight one of $V^k(\mathfrak g)$]\label{l23:thm:weightone}
$\lxxiiiV_1$ is the image of the current-state map; $y_{(1)}x_{(-1)}\lxxiiivac=k\kappa(y,x)\lxxiiivac$; and for $k\ne0$ the
current-state map is injective, so $\lxxiiig\cong\lxxiiiV_1$.
\end{theorem}
\noindent File \leanpath{VOA/\allowbreak AffineGeneral/\allowbreak WeightOne.lean}.
\begin{leancode}
theorem grading_one_eq_range (k : ℂ) : gradingZA D k 1=LinearMap.range (currentState D k) := by

theorem positive_current_currentState (k : ℂ) (x y : 𝔤) :
    vacuumRho κ D.invariant k (ofLoop κ D.invariant y 1) (currentState D k x)=
      (k*κ y x) • vacuumA D k := by

theorem currentState_injective_of_ne_zero (k : ℂ) (hk : k≠0) :
    Function.Injective (currentState D k) := by

def weightOneEquiv (k : ℂ) (hk : k ≠ 0) : 𝔤 ≃ₗ[ℂ] gradingZA D k 1 := by
\end{leancode}

\begin{dictionary}
\lean{≃ₗ[ℂ]} is a linear isomorphism, and only that: no Lean statement in \leanpath{AffineGeneral/} says that it carries the
bracket $[x,y]$ to the weight-one product $u_{(0)}v$, so ``$\lxxiiig\cong\lxxiiiV_1$'' is proved as vector spaces, not as Lie
algebras (release audit 2026-10-27, finding m1). The hypothesis $k\ne0$ is used by the proof (the form argument below). It is not
necessary mathematically: by PBW the map is injective at every level, and for $\mathfrak{sl}_3$ this is proved
(\lean{VOA.AffineSl3.currentState_injective}, \leanpath{VOA/\allowbreak Examples/\allowbreak AffineSl3VacuumBasis.lean}); the general all-level statement is not in the library. The docstring of
\lean{currentState_injective_of_ne_zero} points to the $\mathfrak{sl}_3$ file for it, which is correct only for $\mathfrak{sl}_3$
(Section~\ref{l23:sec:defects}). Same as the textbook for $k\ne0$.
\end{dictionary}

\textbf{Proof on paper.} Spanning as above. The pairing: $[y\otimes t,x\otimes t^{-1}]=[y,x]\otimes1+k\kappa(y,x)$, and
$[y,x]\otimes1$ and $y\otimes t$ kill the vacuum. Injectivity: if $x_{(-1)}\lxxiiivac=0$ then $k\kappa(y,x)=0$ for all $y$;
with $k\ne0$ and $\kappa$ nondegenerate (a consequence of Lie data), $x=0$. \qed

\begin{physics}
The two-point function of currents is $\langle J^y J^x\rangle\propto k\kappa(y,x)$; at nonzero level it is nondegenerate,
so no current is null in the universal theory.
\end{physics}

\begin{status}
Proved; axioms standard (project record, Section~\ref{l23:sec:verification}).
\end{status}

\begin{warning}[title={Trap: $\lxxiiiL_1\cong\lxxiiig$ is not stated for general $\lxxiiig$}]\label{l23:warn:simpleweightone}
Theorem~\ref{l23:thm:simple} gives only that the currents \emph{span} $\lxxiiiL_1$. Injectivity of $\lxxiiig\to\lxxiiiL_1$ for $k\ne0$
is easy on paper: if $x_{(-1)}\lxxiiivac\in I_{\max}$ then $y_{(1)}x_{(-1)}\lxxiiivac=k\kappa(y,x)\lxxiiivac\in I_{\max}$, and a proper
ideal does not contain the vacuum, so $x=0$. It is proved for $\mathfrak{sl}_3$ (\lean{AffineSl3.simpleCurrentState_injective},
\leanpath{VOA/\allowbreak Examples/\allowbreak AffineSl3Simple.lean}) but not in \leanpath{AffineGeneral/}. Nothing in the Deligne chain needs it: the presentations of
Chapter~\ref{ch:L24} only need spanning. This is a gap, not a defect.
\end{warning}

\subsection{Singular vectors lie in the maximal ideal}

The explicit presentations of Chapter~\ref{ch:L24} need to know that specific weight-two vectors vanish in $\lxxiiiL$. The general tool:

\begin{definition}[Positive-mode singular vector]\label{l23:def:singP}
$t\in\lxxiiiV$ is \emph{positive-mode singular} if $J^a_i t=0$ for all $a$ and all $i\ge1$.
\end{definition}
\noindent File \leanpath{VOA/\allowbreak AffineGeneral/\allowbreak Singular.lean}.
\begin{leancode}
def SingP (t : Vacuum D k) : Prop := ∀ (a : ι) (i : ℤ), 1 ≤ i → J D k a i t = 0
\end{leancode}

\begin{dictionary}
Note that the zero modes $J^a_0$ are \emph{not} required to kill $t$: $t$ need not be a highest-weight vector for $\lxxiiig$.
This is the right notion because the set of such $t$ is $\lxxiiig$-stable (\lean{SingP.zeroMode}), so a $\lxxiiig$-module of singular
vectors can be handled at once. Same as the textbook notion of ``singular vector for $\widehat{\lxxiiig}_+$''.
\end{dictionary}

\begin{physics}
A singular vector is a null state: its two-point function with every state vanishes; it decouples from all correlators of the
physical theory.
\end{physics}

\begin{status}
A definition (no proof obligation).
\end{status}

\begin{theorem}[Singular vectors die in $L_k(\mathfrak g)$]\label{l23:thm:singular}
A homogeneous positive-mode singular vector of positive weight lies in the maximal ideal. In particular a quadratic state
$t=\sum_i c_i J^{a_i}_{-1}J^{b_i}_{-1}\lxxiiivac$ lies in it as soon as, for all $x,g$,
\begin{equation*}
\sum_i c_i\Big(\delta_{a_i g}\,k\kappa_{x b_i}+\delta_{b_i g}\,k\kappa_{x a_i}+\sum_c f^{x a_i}{}_c f^{c b_i}{}_g\Big)=0,
\qquad
\sum_i c_i\,k\sum_c f^{x a_i}{}_c\kappa_{c b_i}=0 .
\end{equation*}
\end{theorem}
\noindent Files \leanpath{VOA/\allowbreak AffineGeneral/\allowbreak Singular.lean},
\leanpath{VOA/\allowbreak AffineGeneral/\allowbreak SingularQuad.lean}.
\begin{leancode}
theorem mem_maxIdeal_of_SingP {t : Vacuum D k} (ht : SingP D k t)
    {w : ℤ} (hw : 0 < w) (hmem : t ∈ gradingZA D k w) : t ∈ maxIdeal D k hk := by

def quadState (terms : List (ℂ × ι × ι)) : Vacuum D k :=
  (terms.map fun t => t.1 • st2 D k t.2.1 t.2.2).sum

def quadCoef1 (terms : List (ℂ × ι × ι)) (x g : ι) : ℂ :=
  (terms.map fun t => t.1 * ((if t.2.1 = g then k * formTable D x t.2.2 else 0)
    + (if t.2.2 = g then k * formTable D x t.2.1 else 0)
    + ∑ c, structureConstant D x t.2.1 c * structureConstant D c t.2.2 g)).sum

def quadCoef2 (terms : List (ℂ × ι × ι)) (x : ι) : ℂ :=
  (terms.map fun t => t.1 * (k * ∑ c, structureConstant D x t.2.1 c * formTable D c t.2.2)).sum

theorem quadState_mem_maxIdeal (hk : k + D.h ≠ 0) (terms : List (ℂ × ι × ι))
    (h1 : ∀ x g, quadCoef1 D k terms x g = 0) (h2 : ∀ x, quadCoef2 D k terms x = 0) :
    quadState D k terms ∈ maxIdeal D k hk :=
  mem_maxIdeal_of_SingP D k hk (singP_quad D k terms h1 h2) (by norm_num) (quadState_mem D k terms)
\end{leancode}

\begin{dictionary}
\lean{st2 D k a b} is $J^a_{-1}J^b_{-1}\lxxiiivac$; a list of triples \lean{(c, a, b)} encodes $\sum c\,J^a_{-1}J^b_{-1}\lxxiiivac$.
\lean{quadCoef1 x g} is the coefficient of $J^g_{-1}\lxxiiivac$ in $J^x_1t$ and \lean{quadCoef2 x} the coefficient of $\lxxiiivac$ in
$J^x_2t$ (modes $\ge3$ kill weight-two states automatically). In \lean{mem_maxIdeal_of_SingP} the weight $w>0$ is essential: the
vacuum itself is positive-mode singular of weight $0$ and is not in the ideal. Same as the textbook.
\end{dictionary}

\textbf{Proof on paper.} Let $t\ne0$ be positive-mode singular and homogeneous of weight $w>0$, and let $M=U(\widehat{\lxxiiig})t$ be
the $\widehat{\lxxiiig}$-submodule it generates. Three facts.
\begin{itemize}
\item \emph{$M=U(\widehat{\lxxiiig}_-)U(\lxxiiig)t$, so $M$ lives in weights $\ge w$.} Write
  $\widehat{\lxxiiig}=\widehat{\lxxiiig}_-\oplus\lxxiiig\oplus\widehat{\lxxiiig}_+$ (negative, zero and positive modes, plus the central
  element, which acts by $k$). By PBW, $U(\widehat{\lxxiiig})=U(\widehat{\lxxiiig}_-)U(\lxxiiig)U(\widehat{\lxxiiig}_+)$. The positive modes kill
  $t$, and they also kill every $J^a_0t$ (the set of positive-mode singular vectors is $\lxxiiig$-stable, \lean{SingP.zeroMode}: $J^b_iJ^a_0t
  =J^a_0J^b_it+\sum_cf^{ba}{}_cJ^c_it=0$ for $i\ge1$). Zero modes preserve weight and negative modes raise it.
\item \emph{$M$ is an ideal of the vertex algebra.} An ideal of $\lxxiiiV$ is a subspace stable under all $u_{(n)}$, $u\in\lxxiiiV$.
  Stability under the current modes $J^a_n=(b_a)_{(-1)}\lxxiiivac{}_{(n)}$ holds by construction; since $\lxxiiiV$ is generated by the
  currents, every $u_{(n)}$ is, by the Borcherds identity (iterate formula), an infinite but locally finite combination of products of
  current modes, so $M$ is stable under all of them. (This is Dong's lemma in the form used in Chapter~\ref{ch:L04}; in Lean it is the
  induction on current words that turns ``positive-mode singular'' into ``strongly singular''.) It is graded because $t$ is homogeneous.
\item \emph{$M$ is proper.} The vacuum has weight $0<w$, so $\lxxiiivac\notin M$.
\end{itemize}
So $M$ is a proper graded ideal and $t\in M\subseteq I_{\max}$. For the quadratic test:
$J^x_1J^a_{-1}J^b_{-1}\lxxiiivac=\big(\sum_c f^{xa}{}_cJ^c_0+k\kappa_{xa}\big)J^b_{-1}\lxxiiivac+J^a_{-1}k\kappa_{xb}\lxxiiivac$ and
$J^c_0J^b_{-1}\lxxiiivac=\sum_g f^{cb}{}_gJ^g_{-1}\lxxiiivac$, which gives \lean{quadCoef1}; similarly
$J^x_2J^a_{-1}J^b_{-1}\lxxiiivac=\sum_c f^{xa}{}_c J^c_1J^b_{-1}\lxxiiivac=k\sum_cf^{xa}{}_c\kappa_{cb}\lxxiiivac$. \qed

\begin{example}[$\mathfrak{sl}_2$: the null vector $(J^e_{-1})^2\lxxiiivac$ at level one]\label{l23:ex:sl2sing}
Continue Example~\ref{l23:ex:sl2} ($\kappa=\tr$, $h=2$) and take $t=J^e_{-1}J^e_{-1}\lxxiiivac$, i.e.\ the list \lean{[(1, e, e)]}. The
commutators needed are $[J^x_1,J^e_{-1}]=J^{[x,e]}_0+k\tr(xe)$ and $[J^x_2,J^e_{-1}]=J^{[x,e]}_1$:
\begin{align*}
J^f_1t&=(-J^H_0+k)J^e_{-1}\lxxiiivac+J^e_{-1}(-J^H_0+k)\lxxiiivac=(k-2)J^e_{-1}\lxxiiivac+kJ^e_{-1}\lxxiiivac=(2k-2)\,J^e_{-1}\lxxiiivac,\\
J^H_1t&=2J^e_0J^e_{-1}\lxxiiivac+J^e_{-1}\,2J^e_0\lxxiiivac=2J^{[e,e]}_{-1}\lxxiiivac=0,\qquad J^e_1t=0\quad(\tr(ee)=0,\ [e,e]=0),\\
J^f_2t&=-J^H_1J^e_{-1}\lxxiiivac-J^e_{-1}J^H_1\lxxiiivac=-\big(2J^e_0+k\tr(He)\big)\lxxiiivac=0,
\end{align*}
and similarly $J^H_2t=J^e_2t=0$; modes $\ge3$ kill weight two. So $t$ is positive-mode singular \emph{if and only if $k=1$}. In the
language of the theorem, \lean{quadCoef1 f e} $=k\tr(fe)+k\tr(fe)+\sum_cf^{fe}{}_cf^{ce}{}_e=2k+(-1)\cdot2=2k-2$ (only $c=H$
contributes, with $[f,e]=-H$ and $[H,e]=2e$), and every other coefficient vanishes identically. Hence at $k=1$,
$(J^e_{-1})^2\lxxiiivac\in I_{\max}$ and $(J^e_{-1})^2\lxxiiivac=0$ in $L_1(\mathfrak{sl}_2)$, the familiar relation $e(z)^2=0$ of the
level-one WZW model (in the Frenkel--Kac picture $e(z)=e^{i\sqrt2X(z)}$, and $e(z)e(w)\sim(z-w)^2$ vanishes at $z=w$). The general
pattern is $(J^e_{-1})^{k+1}\lxxiiivac=0$ at positive integer level $k$; the weight-two case is $k+1=2$.
\end{example}

\begin{physics}
At the Deligne levels the vacuum module has null vectors at weight two, spanning the $\lxxiiig$-representation complementary to
$\C\oplus V(2\theta)$ in $\mathrm{Sym}^2\lxxiiig$ ($\theta$ the highest root). Their images in $\mathrm{Sym}\,\lxxiiig$ are the quadratic
generators of the \emph{Joseph ideal} (for $\lxxiiig$ not of type $A$, the unique completely prime primitive ideal of $U(\lxxiiig)$
whose associated variety is the minimal nilpotent orbit closure; its associated graded is the ideal of functions vanishing on
$\overline{\mathbb O_{\min}}$, Chapter~\ref{ch:L24}). Example~\ref{l23:ex:sl2sing} is the simplest analogue. This theorem is the mechanism by which those
null states are shown to vanish in the physical chiral algebra; in four-dimensional language they are the Higgs-branch chiral
ring relations of the minimal nilpotent orbit.
\end{physics}

\begin{status}
Proved; axioms standard for \lean{mem_maxIdeal_of_SingP} (re-run) and \lean{quadState_mem_maxIdeal} (project record).
The structure-constant version \lean{SCLie.ld_quadState_mem_maxIdeal} (\leanpath{SCSingular.lean}) replaces the two families of
conditions by Boolean kernel checks \lean{q1OK}, \lean{q2OK} (defined in \leanpath{SCLieData.lean}) on integer tables. Its
statement takes an \emph{integer} level and integer coefficients, \lean{(k : ℤ)} and \lean{tz : List (ℤ × Fin N × Fin N)}, for Lie
data built by \lean{lieDataOf} (Section~\ref{l23:sec:sc}); so it applies as stated only at the integer Deligne levels
($-2,-3,-4,-6$ for $D_4,E_6,E_7,E_8$), not at $-3/2$, $-5/3$, $-5/2$ ($A_2,G_2,F_4$), where the general
\lean{quadState_mem_maxIdeal} above (complex $k$ and coefficients) must be used.
\end{status}

\subsection{Current coordinates in Li degree zero}

For the presentations of Chapter~\ref{ch:L24} the library needs the map from polynomials in the currents to Li's associated graded
algebra (Chapter~\ref{ch:L11}, Section~\ref{l11:sec:li}), and two facts about its kernel. Recall that Li's filtration
$F^0V\supseteq F^1V\supseteq\cdots$ is decreasing, that $\operatorname{gr}^FV=\bigoplus_pF^pV/F^{p+1}V$ is a commutative algebra
(product induced by $u_{(-1)}v$), and that its degree-zero part is Zhu's $C_2$-algebra $R_V=V/C_2(V)$, $C_2(V)=\operatorname{span}\{u_{(-2)}v\}$
(Theorem~\ref{l11:thm:f1}). The ``Li-degree-zero class'' of a vector is its image in $R_V$. The facts below are proved in full in the
cited file; the argument for (i) is the one of Chapter~\ref{ch:L24}, Theorem~\ref{l24:thm:adprop}, summarized after the statement.

\begin{proposition}[Kernel facts for current coordinates]\label{l23:prop:clsG}
Let $\mathrm{cls}:\C[X_1,\dots,X_n]\to\operatorname{gr}^F L_k(\lxxiiig)$ send $X_i$ to the Li-degree-zero class of the current
$[J^i_{-1}\lxxiiivac]$. Then (i) the kernel is stable under the coadjoint derivations $\mathrm{ad}^*(b)$, $X_i\mapsto\sum f^{bi}{}_cX_c$,
and (ii) if a quadratic state $\sum c_iJ^{a_i}_{-1}J^{b_i}_{-1}\lxxiiivac$ lies in $I_{\max}$, then
$\sum c_iX_{a_i}X_{b_i}\in\ker\mathrm{cls}$.
\end{proposition}
\noindent File \leanpath{VOA/\allowbreak AffineGeneral/\allowbreak LiPresentation.lean}.
\begin{leancode}
abbrev clsG : MvPolynomial (Fin n) ℂ →ₐ[ℂ] liAssociatedGraded W' :=
  aeval fun i => liDegreeZeroClass W' 1 ⟨g i, hg i⟩

theorem clsG_adPoly_eq_zero (b : Fin n) {p : MvPolynomial (Fin n) ℂ}
    (hp : clsG (simpleVOA D k hk) (genS D k hk) (genS_weight_one D k hk) p = 0) :
    clsG (simpleVOA D k hk) (genS D k hk) (genS_weight_one D k hk)
      (adPolyT (structureConstant D) b p) = 0 :=

theorem clsG_quad_eq_zero (terms : List (ℂ × Fin n × Fin n)) (h : quadState D k terms ∈ maxIdeal D k hk) :
    clsG (simpleVOA D k hk) (genS D k hk) (genS_weight_one D k hk) (quadPoly terms) = 0 := by
\end{leancode}

\begin{dictionary}
\lean{MvPolynomial (Fin n) ℂ} is $\C[X_1,\dots,X_n]$; the section variable is \lean{D : LieData κ (Fin n)}, so $n$ is the size of the
index set of $D$, i.e.\ $n=\dim\lxxiiig$ (variables indexed by the basis); \lean{aeval} is the algebra map determined by the images of the variables;
\lean{liAssociatedGraded W'} and \lean{liDegreeZeroClass} are Li's associated graded algebra and the class of a weight-$N$ vector in
Li degree $0$ (Chapter~\ref{ch:L11}; in Li degree zero this is the $C_2$-algebra $R_V$). \lean{genS i} is $[J^i_{-1}\lxxiiivac]\in\lxxiiiL$.
\lean{adPolyT T b} is the derivation of the polynomial ring extending $X_i\mapsto\sum_cT(b,i,c)X_c$. Same as the textbook.
\end{dictionary}

\textbf{Proof on paper.} (i) The zero mode $J^b_0$ preserves $L_k(\lxxiiig)$, Li's filtration and $C_2$, and acts on products by the
Leibniz rule ($J^b_0$ is a derivation of every $n$-th product, by the commutator formula, since $[J^b_0,Y(u,z)]=Y(J^b_0u,z)$). So it
induces a derivation $\delta_b$ of $\operatorname{gr}^F$, and in degree zero this is the Poisson bracket $\{[J^b],\cdot\}$ of $R_V$. On
generators, $\delta_b[J^i_{-1}\lxxiiivac]=[J^{[b_b,b_i]}_{-1}\lxxiiivac]=\sum_cf^{bi}{}_c[J^c_{-1}\lxxiiivac]$, which is the derivation
$\mathrm{ad}^*(b)$ on variables. Thus $\mathrm{cls}\circ\mathrm{ad}^*(b)=\delta_b\circ\mathrm{cls}$ on generators, hence on all
polynomials (both sides are derivations along $\mathrm{cls}$); if $\mathrm{cls}(p)=0$ then $\mathrm{cls}(\mathrm{ad}^*(b)p)=\delta_b(0)=0$.
(ii) The product of $\operatorname{gr}^F$ is induced by the $(-1)$-product, so the class of $J^a_{-1}J^b_{-1}\lxxiiivac=(J^a_{-1}\lxxiiivac)_{(-1)}(J^b_{-1}\lxxiiivac)$
is the product of the classes; and the image of $I_{\max}$ in $L_k$ is zero. \qed

\begin{physics}
$\operatorname{gr}^F L_k$ in degree zero is the ring of functions on the associated variety (the Higgs branch, in the 4d
dictionary); (i) says the relations are $G$-invariant as a set, (ii) that a null state at weight two becomes a quadratic
chiral-ring relation.
\end{physics}

\begin{status}
Proved; axioms standard (project record, Section~\ref{l23:sec:verification}).
\end{status}

%======================================================================
\section{Concrete simple Lie algebras I: \texorpdfstring{$\mathfrak{sl}_n$}{sl(n)} and a root-space criterion}\label{l23:sec:sln}
%======================================================================

Theorem~\ref{l23:thm:simpleaffine} is only useful if some concrete Lie algebra in Lean carries an instance of
\lean{LieAlgebra.IsSimple}. Mathlib v4.30 has none: it defines $\mathfrak{sl}_n$
(\lean{LieAlgebra.SpecialLinear.sl}) and, through \lean{Matrix.ToLieAlgebra}, a Serre-relation construction of the exceptional
algebras, but proves neither simplicity nor finite-dimensionality for any of them.

\begin{warning}[title={Dead end: Mathlib's Serre construction}]\label{l23:warn:serre}
The obvious route to $\mathfrak e_8$ is Mathlib's \lean{LieAlgebra.e₈ := Matrix.ToLieAlgebra …} (file
\texttt{Mathlib/Algebra/Lie/SerreConstruction.lean}): the Lie algebra presented by Chevalley generators and Serre relations of
the $E_8$ Cartan matrix. Mathlib does not prove that it is finite-dimensional, nor that it is simple, nor its dimension; each
of these is a substantial theorem (Serre's theorem). The project recorded this route as a dead end and
built the exceptional algebras as lattice weight-one spaces and from checked structure constants instead.
\end{warning}

\begin{theorem}[$\mathfrak{sl}_n$ is simple]\label{l23:thm:sln}
For every finite type $n$ with at least two elements, $\mathfrak{sl}_n(\C)$ is a simple Lie algebra. Consequently
$V^k(\mathfrak{sl}_m)$ exists in the normalization $\lxxiiiKil/(2m)$ for $k\ne-m$.
\end{theorem}
\noindent File \leanpath{VOA/\allowbreak AffineGeneral/\allowbreak SlnSimple.lean}.
\begin{leancode}
instance instIsSimple [Nontrivial n] : LieAlgebra.IsSimple ℂ (sl n ℂ) where

def slAffineVOA (m : ℕ) [Nontrivial (Fin m)] (k : ℂ) (hk : k + (m : ℂ) ≠ 0) (hm : (m : ℂ) ≠ 0) :
    VertexOperatorAlgebra (Vacuum (standardData (sl (Fin m) ℂ) (m : ℂ) hm) k) :=
  standardAffineVOA (sl (Fin m) ℂ) (m : ℂ) hm k hk
\end{leancode}

\begin{dictionary}
\lean{sl n ℂ} is Mathlib's Lie subalgebra of trace-zero matrices indexed by the finite type \lean{n};
\lean{[Nontrivial n]} means $|n|\ge2$ (for $|n|\le1$, $\mathfrak{sl}_n=0$ is not simple). The hypothesis \lean{hm : (m : ℂ) ≠ 0}
of \lean{slAffineVOA} is \textbf{redundant}: it follows from \lean{[Nontrivial (Fin m)]} ($m\ge2$). Harmless, listed in
Section~\ref{l23:sec:defects}.
\end{dictionary}

\textbf{Proof on paper.} Let $I\ne0$ be an ideal and $v\in I$, $v\ne0$. If $v$ has an off-diagonal entry $v_{ji}\ne0$, then
$[E_{ij},[E_{ij},v]]=-2v_{ji}E_{ij}$, so $E_{ij}\in I$. If $v$ is diagonal and nonzero, it has two distinct diagonal entries (trace
zero), and $[E_{ij},v]=(v_{jj}-v_{ii})E_{ij}$ gives some $E_{ij}\in I$. From one $E_{ij}$, brackets
$[E_{ij},E_{jl}]=E_{il}$ and $[E_{ki},E_{ij}]=E_{kj}$ produce all $E_{kl}$, $k\ne l$, and then $[E_{kl},E_{lk}]=E_{kk}-E_{ll}$ span the
diagonal; so $I=\mathfrak{sl}_n$. Non-abelian: $[E_{ij},E_{ji}]\ne0$. \qed

\begin{physics}
$SU(n)$ current algebras, the $\mathfrak{sl}_n$ WZW models; $\mathfrak{sl}_2$ is the subject of
Chapters~\ref{ch:L08} and~\ref{ch:L09}, and a chapter on $\mathfrak{sl}_3$ is in preparation.
\end{physics}

\begin{status}
Proved; \texttt{\#print axioms SlSimple.instIsSimple} standard (re-run).
\end{status}

\begin{warning}[title={Trap: The normalization of \texttt{slAffineVOA} is not proved standard}]\label{l23:warn:slnnorm}
The docstring of \lean{slAffineVOA} says ``in the standard normalization ($\kappa=\lxxiiiKil/(2n)$, $h=n$)''. On paper this is true:
$\lxxiiiKil(x,y)=2n\tr(xy)$ on $\mathfrak{sl}_n$, so $\lxxiiiKil/(2n)$ is the trace form, whose long roots have length $2$, and
$\lxxiiihv(\mathfrak{sl}_n)=n$. But in Lean, by Theorem~\ref{l23:thm:simpleaffine}(3), $h=n$ holds for \emph{any} choice of the
parameter; no theorem identifies $\lxxiiiKil/(2n)$ with the trace form. Contrast the lattice and structure-constant types, where
\lean{*_standard} theorems do this (Theorem~\ref{l23:thm:normalization}).
\end{warning}

\subsection{A root-space criterion for simplicity}

For the exceptional algebras the $\mathfrak{sl}_n$ argument is generalized to an abstract criterion, stated for any Lie algebra with
a family of ``root vectors'' and ``coroots''.

\begin{theorem}[Root-space criterion]\label{l23:thm:rootspace}
Let $L$ be a complex Lie algebra, $R$ a finite nonempty set, $e,h:R\to L$, $p:R\times R\to\C$, and $-:R\to R$, such that
\begin{enumerate}
\item the $e_\beta$ and $h_\alpha$ together span $L$, and $e_\beta\ne0$;
\item $[h_\alpha,e_\beta]=p(\alpha,\beta)e_\beta$ and $[h_\alpha,h_\beta]=0$;
\item $p(\alpha,\alpha)\ne0$, and the weights separate: for $\alpha\ne\beta$ some $\gamma$ has $p(\gamma,\alpha)\ne p(\gamma,\beta)$;
\item $[e_\alpha,e_{-\alpha}]=c_\alpha h_\alpha$ with $c_\alpha\neq0$;
\item no nonzero element of $\operatorname{span}(h)$ commutes with all $e_\beta$;
\item the graph on $R$ with an edge $\alpha\to\beta$ when $p(\alpha,\beta)\ne0$ is connected (every $\beta$ is reachable from every
  $\alpha$).
\end{enumerate}
Then $L$ is simple.
\end{theorem}
\noindent File \leanpath{VOA/\allowbreak AffineGeneral/\allowbreak RootSpaceSimple.lean}.
\begin{leancode}
theorem isSimple_of_roots [Fintype R] [DecidableEq R] [Nonempty R]
    (hspan : span ℂ (Set.range e ∪ Set.range hb) = ⊤)
    (he0 : ∀ β, e β ≠ 0)
    (hbe : ∀ α β, ⁅hb α, e β⁆ = pair α β • e β)
    (hhb : ∀ α β, ⁅hb α, hb β⁆ = 0)
    (hdiag : ∀ α, pair α α ≠ 0)
    (sep : ∀ α β, α ≠ β → ∃ γ, pair γ α ≠ pair γ β)
    (neg : R → R) (hneg : ∀ α, ∃ c : ℂ, c ≠ 0 ∧ ⁅e α, e (neg α)⁆ = c • hb α)
    (center : ∀ y ∈ span ℂ (Set.range hb), (∀ β, ⁅y, e β⁆ = 0) → y = 0)
    (conn : ∀ α β, Relation.ReflTransGen (fun a b => pair a b ≠ 0) α β) :
    LieAlgebra.IsSimple ℂ L where
\end{leancode}

\begin{dictionary}
\lean{span ℂ S = ⊤} says $S$ spans $L$. \lean{Relation.ReflTransGen r α β} is the reflexive-transitive closure: a finite chain
$\alpha=\alpha_0,\alpha_1,\dots,\alpha_m=\beta$ with $r(\alpha_i,\alpha_{i+1})$. The map \lean{neg} is arbitrary data, constrained
only by \lean{hneg}. Nothing requires $R$ to be a root system: the criterion is purely about weight vectors, which is what makes it
checkable by finite computation. Every hypothesis is used. Same as the statement above, which is a slight generalization of the
textbook argument (Humphreys \S14) for simplicity of a Lie algebra with an irreducible root system.
\end{dictionary}

\textbf{Proof on paper.} Let $I\ne0$ be an ideal, $v\in I$ nonzero; write $v=\sum_\beta c_\beta e_\beta+y$, $y\in\operatorname{span}h$.
If all $c_\beta=0$ then $v=y\neq0$, and by (5) some $[y,e_\beta]=\lambda e_\beta\ne0$ lies in $I$, so $e_\beta\in I$. Otherwise pick
$\beta$ with $c_\beta\ne0$. The operators $\lxxiiiad(h_\gamma)-p(\gamma,\delta)$ kill $e_\delta$, act on $e_\beta$ by the scalar
$p(\gamma,\beta)-p(\gamma,\delta)$, and by (3) for each $\delta\ne\beta$ some $\gamma$ makes this scalar nonzero; applying the product
of these operators (one per $\delta\neq\beta$), and finally $\lxxiiiad(h_\beta)$ (to kill the Cartan part, using $p(\beta,\beta)\ne0$),
projects $v$ onto a nonzero multiple of $e_\beta$ inside $I$. Once $e_\beta\in I$: $[e_\beta,e_{-\beta}]=c_\beta h_\beta\in I$; then
$[h_\beta,e_\gamma]=p(\beta,\gamma)e_\gamma$ puts $e_\gamma\in I$ whenever $p(\beta,\gamma)\ne0$; by connectivity (6) all
$e_\gamma\in I$, then all $h_\gamma$, so $I=L$ by (1). Non-abelian: $[h_\alpha,e_\alpha]=p(\alpha,\alpha)e_\alpha\ne0$. \qed

\begin{physics}
The criterion says: a symmetry algebra is simple if its charged currents ($e_\beta$) carry distinguishable charges under the Cartan
currents ($h_\alpha$), no Cartan current is neutral under everything, and the charges link all the currents together. It is the
algebraic content of ``the gauge group does not factorize''.
\end{physics}

\begin{status}
Proved; \texttt{\#print axioms RootSpace.isSimple\_of\_roots} standard (re-run).
\end{status}

%======================================================================
\section{Concrete simple Lie algebras II: simply-laced types as lattice currents}\label{l23:sec:lattice}
%======================================================================

\subsection{The root lattices}

\begin{definition}[Cartan lattices]\label{l23:def:cartanlattice}
For a symmetric integer matrix $G$, $\mathrm{ofCartan}(G)$ is the lattice $\Z^n$ with bilinear form $\langle v,w\rangle=v^{\mathsf T}Gw$.
The lattices $a_2,d_4,e_6,e_7,e_8$ are $\mathrm{ofCartan}$ of Mathlib's Cartan matrices of types $A_2,D_4,E_6,E_7,E_8$; each is even and
positive definite.
\end{definition}
\noindent File \leanpath{VOA/\allowbreak Foundations/\allowbreak CartanLattices.lean}.
\begin{leancode}
noncomputable def ofCartan {n : ℕ} (G : Matrix (Fin n) (Fin n) ℤ) (hG : G.IsSymm) :
    IntegralLattice where
  carrier := Fin n → ℤ
  form := Matrix.toBilin' G

theorem sos_A2 (v : Fin 2 → ℤ) :
    ((Matrix.toBilin' (CartanMatrix.A 2) v v : ℤ) : ℚ) =
      (2 : ℚ) * ((v 0 : ℚ) + (-1/2 : ℚ) * (v 1 : ℚ)) ^ 2 + (3/2 : ℚ) * ((v 1 : ℚ)) ^ 2 := by

noncomputable def e8 : IntegralLattice := ofCartan (CartanMatrix.E₈) CartanMatrix.E₈_isSymm

theorem e8_even : e8.Even :=

theorem e8_posDef : e8.PosDef :=
\end{leancode}

\begin{dictionary}
\lean{IntegralLattice} is the library's structure of Chapter~\ref{ch:L02} (Definition~\ref{l02:def:lattice}): a free $\Z$-module
\lean{carrier} with a symmetric $\Z$-valued form. \lean{Matrix.toBilin' G} is $(v,w)\mapsto v^{\mathsf T}Gw$. \lean{CartanMatrix.E₈} etc.\
are Mathlib's Cartan matrices (Bourbaki numbering). \lean{Even}: $\langle v,v\rangle\in2\Z$; \lean{PosDef}: $\langle v,v\rangle>0$ for
$v\ne0$. These are exactly the hypotheses of \lean{latticeVOAGen} (Chapter~\ref{ch:L05}). Same as the textbook.
\end{dictionary}

\textbf{Proof on paper.} Evenness: the diagonal of $G$ is $2$, so $v^{\mathsf T}Gv=2\big(\sum v_i^2+\sum_{i<j}G_{ij}v_iv_j\big)$.
Positive definiteness: an $LDL^{\mathsf T}$ decomposition $v^{\mathsf T}Gv=\sum_kd_k\big(v_k+\sum_{i>k}L_{ki}v_i\big)^2$ with all
pivots $d_k>0$ (for $A_2$: $d=(2,\tfrac32)$, as \lean{sos_A2} shows); if the sum is $\le0$ every square vanishes and back-substitution gives
$v=0$. The identities were produced by an untrusted script and are re-checked by \lean{ring}. The products of the pivots are the
determinants $3,4,3,2,1$. \qed

\begin{physics}
These are the momentum lattices of $n$ chiral bosons on the maximal torus of the simply-laced group $G$, at the special
(``enhanced-symmetry'') point where the vectors of norm $2$ --- the roots --- give extra massless currents that enhance $U(1)^n$ to
$G$. Only $E_8$ is unimodular (determinant $1$), so only $V_{E_8}$ is a self-dual, modular-invariant theory by itself; for $A_2$,
$D_4$, $E_6$, $E_7$ (determinants $3,4,3,2$) $V_L$ is a chiral algebra whose modular-invariant completion needs the other cosets
$L^*/L$.
\end{physics}

\begin{status}
Proved; axioms standard (project record, Section~\ref{l23:sec:verification}).
\end{status}

\begin{definition}[Lattice VOAs and their weight-one Lie algebras]\label{l23:def:latticeLie}
For each type $T\in\{A_2,D_4,E_6,E_7,E_8\}$, $V_T$ is the lattice vertex algebra of the root lattice, of central charge
$\operatorname{rank}$, and $\mathfrak t=(V_T)_1$ with bracket $[u,v]=u_{(0)}v$.
\end{definition}
\noindent File \leanpath{VOA/\allowbreak Examples/\allowbreak DeligneLatticeLie.lean}.
\begin{leancode}
abbrev e8VOA : VertexOperatorAlgebra (LatModuleGen e8) :=
  latticeVOAGen e8 e8_even e8_posDef

abbrev e8Lie := e8VOA.grading 1

theorem e8VOA_centralCharge : e8VOA.centralCharge = 8 := by
\end{leancode}

\begin{dictionary}
\lean{latticeVOAGen L hEven hPos} is the lattice VOA of Chapter~\ref{ch:L05} (Theorem~\ref{l05:thm:latticeVOAGen});
\lean{e8VOA.grading 1} is its weight-one subspace (this \lean{e8Lie} is \lean{DeligneLatticeLie.e8Lie}; the \lean{CartanSimple.e8Lie}
of Theorem~\ref{l23:thm:cartansimple} is the same type, written as \lean{V1 e8 e8_even e8_posDef}), which carries the \lean{LieAlgebra ℂ} instance of Chapter~\ref{ch:L04}
(\leanpath{Core/\allowbreak LieAlgebraV1.lean}), with Jacobi and antisymmetry derived from the VOA axioms. The same pattern defines
\lean{a2VOA}, \lean{d4VOA}, \lean{e6VOA}, \lean{e7VOA} and their \lean{*Lie}, with central charges $2,4,6,7$. Same as the textbook.
\end{dictionary}

\textbf{Proof on paper.} The central charge of a lattice VOA is its rank (Chapter~\ref{ch:L05}), and $\mathrm{rank}\,\Z^n=n$. \qed

\begin{physics}
$V_{E_8}$ is the chiral algebra of eight bosons on the $E_8$ torus, central charge $8$, and its currents form $\mathfrak e_8$. By the
Frenkel--Kac theorem it is the level-one WZW chiral algebra $L_1(\mathfrak e_8)$; this isomorphism is \emph{not} in Lean, and the
library uses $V_{E_8}$ only through its weight-one Lie algebra.
\end{physics}

\begin{status}
Proved; axioms standard (\texttt{\#print axioms} in the file itself; \lean{e8VOA_centralCharge} re-run).
\end{status}

\begin{warning}[title={Trap: stale docstring in \texttt{DeligneLatticeLie.lean}}]\label{l23:warn:stalelattice}
The module docstring of \leanpath{VOA/\allowbreak Examples/\allowbreak DeligneLatticeLie.lean} says, under ``Not proved here'', that the
identification with the simple Lie algebra of the named type and the dimension $\mathrm{rank}+\#\text{roots}$ ``are the next steps''. This is
literally true of that file, but reads as if they were open. Simplicity and the dimensions are now proved, in \leanpath{CartanLatticeSimple.lean} (Theorem~\ref{l23:thm:cartansimple}),
\leanpath{RootCounts.lean} and \leanpath{E8Dimension.lean} (Theorems~\ref{l23:thm:e8dim}, \ref{l23:thm:rootcounts}). The pointer is stale;
nothing proved is affected (the identification of the \emph{type} remains unformalised, Remark~\ref{l23:rem:type}). Listed in
Section~\ref{l23:sec:defects}.
\end{warning}

\subsection{Simplicity of the lattice current algebras}

Inside $(V_L)_1$ there are the vectors $e^\beta$ for each root $\beta$ (norm $2$) and the Heisenberg states $\beta(-1)\lxxiiivac$.

\begin{definition}[Root vectors and coroots of $(V_L)_1$]\label{l23:def:eVhV}
For $L$ even positive definite, \lean{Root L hPos} is the finite set of $\beta\in L$ with $\langle\beta,\beta\rangle=2$;
$e_\beta=e^\beta$ and $h_\beta=\beta(-1)\lxxiiivac$ in $(V_L)_1$.
\end{definition}
\noindent File \leanpath{VOA/\allowbreak AffineGeneral/\allowbreak LatticeV1Simple.lean}.
\begin{leancode}
abbrev Root : Type _ := ↥(normEq L hPos 2)

abbrev V1 : Type _ := ↥((latticeVOAGen L hEven hPos).grading 1)

def eV (β : Root L hPos) : V1 L hEven hPos :=
  ⟨latChargeGen L β.1, latChargeGen_mem_grading_one L hEven hPos (root_form β)⟩

def hV (β : Root L hPos) : V1 L hEven hPos :=
  ⟨hSt L β.1, hSt_mem_gradingLatGen L β.1⟩

def RootsGenerate : Prop := AddSubgroup.closure {v : L.carrier | L.form v v = 2} = ⊤
\end{leancode}

\begin{dictionary}
\lean{normEq L hPos 2} is the finite set (\lean{Finset}) of lattice vectors of norm $2$ (finite by positive definiteness);
\lean{↥} turns it into a type. \lean{latChargeGen L β} is the state $e^\beta$ (vacuum of the Fock space of charge $\beta$),
\lean{hSt L β} is $\beta(-1)\lxxiiivac$ (Chapter~\ref{ch:L05}). \lean{RootsGenerate}: the roots generate $L$ as a group. Same as the textbook.
\end{dictionary}

\textbf{Why these vectors lie in weight one.} $e^\beta$ has weight $\langle\beta,\beta\rangle/2=1$ and $\beta(-1)\lxxiiivac$ has weight $1$;
these definitions bundle each with that proof, as elements of the subtype \lean{V1}.

\begin{physics}
$h_\beta$ are the $U(1)^n$ currents $i\partial X\cdot\beta$; $e_\beta=e^{i\beta\cdot X}$ are the vertex operators of charge $\beta$,
which have conformal weight one, hence are extra massless currents, exactly when $\langle\beta,\beta\rangle=2$.
\end{physics}

\begin{status}
Definitions (no proof obligation).
\end{status}

\begin{theorem}[Frenkel--Kac brackets, spanning, dimension and simplicity]\label{l23:thm:latticesimple}
Let $L$ be even and positive definite.
\begin{enumerate}
\item $[h_\alpha,e_\beta]=\langle\alpha,\beta\rangle e_\beta$, $[h_\alpha,h_\beta]=0$, $[e_\alpha,e_{-\alpha}]=\varepsilon(\alpha,-\alpha)h_\alpha$, $e_\beta\ne0$.
\item If the roots generate $L$, the $e_\beta$ and $h_\beta$ span $(V_L)_1$.
\item $\dim(V_L)_1=\operatorname{rank}L+\#\{\beta:\langle\beta,\beta\rangle=2\}$.
\item If a family $s$ of roots generates $L$ and is connected (the graph $\langle s_i,s_j\rangle\ne0$ is connected), $(V_L)_1$ is simple.
\end{enumerate}
\end{theorem}
\noindent File \leanpath{VOA/\allowbreak AffineGeneral/\allowbreak LatticeV1Simple.lean}.
\begin{leancode}
theorem hV_eV (α β : Root L hPos) :
    ⁅hV L hEven hPos α, eV L hEven hPos β⁆ = ((L.form α.1 β.1 : ℤ) : ℂ) • eV L hEven hPos β := by

theorem hV_hV (α β : Root L hPos) : ⁅hV L hEven hPos α, hV L hEven hPos β⁆ = 0 := by

theorem eV_eV_neg (α : Root L hPos) :
    ⁅eV L hEven hPos α, eV L hEven hPos (Root.neg L hPos α)⁆
      = latticeCocycle L α.1 (-α.1) • hV L hEven hPos α := by

theorem eV_ne_zero (β : Root L hPos) : eV L hEven hPos β ≠ 0 := by

theorem span_eV_hV (hgen : RootsGenerate L) :
    span ℂ (Set.range (eV L hEven hPos) ∪ Set.range (hV L hEven hPos)) = ⊤ := by

theorem finrank_V1 : Module.finrank ℂ (V1 L hEven hPos) = rank L + (normEq L hPos 2).card := by

theorem isSimple_of_simpleRoots {ι : Type*} [Nonempty ι] (s : ι → Root L hPos)
    (hsgen : AddSubgroup.closure (Set.range fun i => (s i).1) = ⊤)
    (hsconn : ∀ i j, Relation.ReflTransGen (fun i j => L.form (s i).1 (s j).1 ≠ 0) i j) :
    LieAlgebra.IsSimple ℂ (V1 L hEven hPos) := by
\end{leancode}

\begin{dictionary}
\lean{Root.neg L hPos α} is the root $-\alpha$. \lean{eV_ne_zero} is part (1)'s $e_\beta\ne0$: the Fock vacuum of charge $\beta$ is
nonzero. \lean{latticeCocycle L α β} is the FLM sign $\varepsilon(\alpha,\beta)=\pm1$ of Chapter~\ref{ch:L02} (Theorem~\ref{l02:thm:cocycle});
it is nonzero, which is what the criterion needs. \lean{rank L} is the $\Z$-rank. In (4) the family $s$ is \emph{any} connected
generating family of roots, not necessarily a base of the root system; for the Cartan lattices the standard basis vectors serve.
\lean{finrank_V1} needs no generation hypothesis. Same as the textbook (Frenkel--Kac; Kac, \emph{Vertex algebras for beginners},
\S5.6).
\end{dictionary}

\textbf{Proof on paper.} Recall (Chapters~\ref{ch:L02}, \ref{ch:L05}) that $\varepsilon:L\times L\to\{\pm1\}$ is a bimultiplicative
cocycle with $\varepsilon(\alpha,\beta)\varepsilon(\beta,\alpha)^{-1}=(-1)^{\langle\alpha,\beta\rangle}$, and that
$Y(e^\alpha,z)e^\beta=\varepsilon(\alpha,\beta)z^{\langle\alpha,\beta\rangle}\exp\big(\sum_{n\ge1}\alpha(-n)z^n/n\big)e^{\alpha+\beta}$, while the
bracket on weight one is $[u,v]=u_{(0)}v$, the coefficient of $z^{-1}$ in $Y(u,z)v$.
(1) $\alpha(-1)_{(0)}=\alpha(0)$ acts on $e^\beta$ by $\langle\alpha,\beta\rangle$; $\alpha(0)\beta(-1)\lxxiiivac=0$. For $\beta=-\alpha$,
$Y(e^\alpha,z)e^{-\alpha}=\varepsilon(\alpha,-\alpha)z^{-2}\big(\lxxiiivac+z\,\alpha(-1)\lxxiiivac+O(z^2)\big)$; the coefficient of $z^{-1}$ is
$\varepsilon(\alpha,-\alpha)\alpha(-1)\lxxiiivac$, which is $e^\alpha_{(0)}e^{-\alpha}$. (2) Weight one of $V_L=\bigoplus_\gamma
M(1)\otimes e^\gamma$ is $M(1)_1\otimes e^0\oplus\bigoplus_{\langle\gamma,\gamma\rangle=2}\C e^\gamma$, since $M(1)\otimes e^\gamma$ starts in
weight $\langle\gamma,\gamma\rangle/2$; $M(1)_1=\mathfrak h$ is spanned by the $h_\beta$ because the roots generate $L$. (3) The graded
dimension is $\sum_n\theta$-coefficients times partition counts; in weight one only $\gamma=0$ (contributing $\operatorname{rank}L$) and
$\langle\gamma,\gamma\rangle=2$ (contributing $1$ each) occur. (4) Apply Theorem~\ref{l23:thm:rootspace}: separation of roots and the centre
condition follow from nondegeneracy of the form on $L\otimes\C$ (an element of $\mathfrak h$ commuting with all $e_\beta$ pairs to zero
with all roots, hence with $L$); connectivity of all roots follows because every root has nonzero pairing with some $s_i$ (else it
would be orthogonal to all of $L$, contradicting $\langle\beta,\beta\rangle=2$). \qed

\begin{physics}
This is the statement that at the enhanced-symmetry point the $U(1)^n$ symmetry of $n$ bosons on the torus of a simply-laced group
is enhanced to the full non-abelian group: the Frenkel--Kac mechanism, and the origin of the $E_8\times E_8$ and
$\mathrm{Spin}(32)/\Z_2$ gauge symmetries of the heterotic string. (In the Higgs mechanism a non-abelian symmetry is broken to
$U(1)^n$ and the charged vector bosons become massive; here, moving the radius \emph{to} the special point, the charged currents
$e^{i\beta\cdot X}$ become massless and the symmetry is restored. This is an analogy, not part of the proof.)
\end{physics}

\begin{status}
Proved; \texttt{\#print axioms} standard for \lean{isSimple_of_simpleRoots}, \lean{finrank_V1} and \lean{eV_ne_zero} (re-run).
\end{status}

\begin{theorem}[Simple Lie algebras of types $A_2,D_4,E_6,E_7,E_8$]\label{l23:thm:cartansimple}
Let $G$ be a symmetric integer matrix with diagonal $2$ and connected Dynkin diagram, such that the lattice $\mathrm{ofCartan}\,G$ is
even and \emph{positive definite} (for a Cartan matrix in the usual sense, i.e.\ of finite type, positive definiteness holds; it does
\emph{not} follow from ``diagonal $2$ and connected'', e.g.\ affine Cartan matrices are only semidefinite). Then
$(V_{\mathrm{ofCartan}\,G})_1$ is simple and finite-dimensional. In particular $\mathfrak{a}_2,\mathfrak{d}_4,\mathfrak e_6,\mathfrak e_7,\mathfrak e_8:=(V_T)_1$ are simple, and $V^k$
of each exists in the normalization $\lxxiiiKil/(2h)$ with $h=3,6,12,18,30$.
\end{theorem}
\noindent File \leanpath{VOA/\allowbreak AffineGeneral/\allowbreak CartanLatticeSimple.lean}.
\begin{leancode}
variable {n : ℕ} (G : Matrix (Fin n) (Fin n) ℤ) (hG : G.IsSymm)

variable (hEven : (ofCartan G hG).Even) (hPos : (ofCartan G hG).PosDef)

theorem cartan_isSimple (hdiag : ∀ i, G i i = 2) (i0 : Fin n) (d : Fin n → ℕ)
    (hd : ∀ i, i = i0 ∨ ∃ j, d j < d i ∧ G j i ≠ 0) :
    LieAlgebra.IsSimple ℂ (V1 (ofCartan G hG) hEven hPos) := by

theorem cartan_finiteDimensional (hdiag : ∀ i, G i i = 2) :
    FiniteDimensional ℂ (V1 (ofCartan G hG) hEven hPos) :=

abbrev e8Lie : Type := V1 e8 e8_even e8_posDef

instance e8Lie_isSimple : LieAlgebra.IsSimple ℂ e8Lie :=
  cartan_isSimple _ _ _ _ (by decide) 3 ![2, 1, 1, 0, 1, 2, 3, 4] (by decide)

instance e8Lie_finiteDimensional : FiniteDimensional ℂ e8Lie :=
  cartan_finiteDimensional _ _ _ _ (by decide)

def e8AffineVOA (k : ℂ) (hk : k + 30 ≠ 0) :
    VertexOperatorAlgebra (Vacuum (standardData e8Lie 30 (by norm_num)) k) :=
  standardAffineVOA e8Lie 30 (by norm_num) k hk
\end{leancode}

\begin{dictionary}
The hypotheses of \lean{cartan_isSimple}, in full: the section variables \lean{G} (an $n\times n$ integer matrix),
\lean{hG : G.IsSymm} (symmetry; here, over $\Z$, plain transpose-symmetry), \lean{hEven : (ofCartan G hG).Even} and
\lean{hPos : (ofCartan G hG).PosDef} (the lattice is even and positive definite; both are needed to form the lattice VOA
\lean{latticeVOAGen}, Chapter~\ref{ch:L05}), and the explicit \lean{hdiag} (diagonal $2$, so the standard basis vectors are roots),
\lean{i0}, \lean{d}, \lean{hd} (connectivity). \lean{hEven} is in fact implied by \lean{hdiag} and symmetry (proof on paper of Definition~\ref{l23:def:cartanlattice}), but
\lean{hPos} is an independent input: it is what makes the statement about finite type. \lean{cartan_finiteDimensional} needs only
\lean{hdiag} besides the section variables. In the instances these are discharged by \lean{e8_even}, \lean{e8_posDef}
(Definition~\ref{l23:def:cartanlattice}) and \lean{by decide}.
Connectivity of the Dynkin diagram is certified by a ``distance function'' $d$ to a base node $i_0$: every node other than $i_0$ has a
neighbour ($G_{ji}\ne0$) strictly closer to $i_0$. For $E_8$, $i_0=3$ (the trivalent node in Mathlib's numbering) and
$d=(2,1,1,0,1,2,3,4)$; \lean{by decide} checks both conditions. \lean{![…]} is a vector literal. \lean{CartanSimple.e8Lie} is the same type
as \lean{DeligneLatticeLie.e8Lie}, and the file re-exports the instances for the latter. The numbers $30,18,12,6,3$ are the dual Coxeter
numbers; that $\lxxiiiKil/(2h)$ is then the standard normalization is Theorem~\ref{l23:thm:normalization}. Same as the textbook,
once ``Cartan matrix'' is read as ``of finite type'' (positive definite).
\end{dictionary}

\textbf{Proof on paper.} Theorem~\ref{l23:thm:latticesimple}(4) with $s_i$ the standard basis vectors (roots, since $G_{ii}=2$; they
generate $\Z^n$; connected by $d$). Finite dimension from spanning by finitely many vectors. \qed

\begin{physics}
The library's $\mathfrak e_8$ \emph{is} the algebra of conserved currents of the $E_8$ lattice CFT; likewise for the other four types.
The theorem says that these current algebras do not split into commuting pieces (the gauge group is simple).
\end{physics}

\begin{status}
Proved; \texttt{\#print axioms} standard for \lean{CartanSimple.e8Lie_isSimple}, \lean{cartan_isSimple},
\lean{cartan_finiteDimensional} (re-run). \lean{cartan_isSimple} is conditional on its listed hypotheses, in particular \lean{hPos};
the instances for $A_2,\dots,E_8$ discharge them and are unconditional.
\end{status}

\subsection{Dimensions: counting roots}

\begin{theorem}[$\dim\mathfrak e_8=248$, from $\theta_{E_8}=E_4$]\label{l23:thm:e8dim}
(i) For every modular form $f=\sum_nc_nq^n$ of weight $4$ for $SL_2(\Z)$, $c_1=240\,c_0$. (ii) The $E_8$ lattice has $240$ vectors of
norm $2$, and $\dim\mathfrak e_8=8+240=248$.
\end{theorem}
\noindent The input for (ii) is that $E_8$ is even and unimodular (here: an explicit inverse of the Cartan matrix), so that its theta
series is a modular form of weight $4$; that theorem is from the vendored lattice library (Chapter~\ref{ch:02}), and the modular forms are
those of Chapter~\ref{ch:L14}. Lean proves the coefficient relation (i) for an arbitrary weight-four form and applies it to
$\theta_{E_8}$; it does not state $\theta_{E_8}=E_4$ as an identity of series (true on paper, since $\dim M_4=1$).
\noindent File \leanpath{VOA/\allowbreak AffineGeneral/\allowbreak E8Dimension.lean}.
\begin{leancode}
theorem coeff_one_of_weight_four (f : ModularForm 𝒮ℒ ((4 : ℕ) : ℤ)) :
    (qExpansion 1 (f : ℍ → ℂ)).coeff 1 = 240 * (qExpansion 1 (f : ℍ → ℂ)).coeff 0 := by

theorem e8Inv_mul : e8Inv * CartanMatrix.E₈ = 1 := by decide

theorem e8E_repCount_two : e8E.repCount 2 = 240 := by

theorem card_roots_e8 : (normEq e8 e8_posDef 2).card = 240 := by

theorem finrank_e8Lie : Module.finrank ℂ e8Lie = 248 := by
\end{leancode}

\begin{dictionary}
\lean{ModularForm 𝒮ℒ 4} is Mathlib's space of level-one modular forms of weight $4$; \lean{qExpansion 1 f} its $q$-expansion
$\sum c_nq^n$ (Chapter~\ref{ch:L14}). \lean{e8Inv} is an explicit integer matrix (printed in the file) and \lean{e8Inv_mul} checks
by \lean{decide} that it inverts the Cartan matrix, so $\det=\pm1$. \lean{e8E} is the same lattice as \lean{e8}, re-expressed in the
vendored \texttt{LatticesAndCodes} library (Chapter~\ref{ch:02}), where \lean{thetaModularForm} (the theta series of an even
unimodular positive-definite lattice as a modular form of weight $\mathrm{rank}/2$) and \lean{repCount m} (the number of vectors of
norm $m$) live. Same as the textbook.
\end{dictionary}

\textbf{Proof on paper.} $M_4(SL_2(\Z))$ is spanned by $E_4=1+240q+\dots$ (Mathlib's basis of monomials $E_4^aE_6^b$; only $(a,b)=(1,0)$
has weight $4$), so $c_1=240c_0$ for every $f\in M_4$. For $\theta_{E_8}=\sum_v q^{\langle v,v\rangle/2}$, $c_0=1$ (the zero vector) and
$c_1=\#\{\langle v,v\rangle=2\}$. Then Theorem~\ref{l23:thm:latticesimple}(3). \qed

\begin{physics}
$\theta_{E_8}/\eta^8$ is the partition function of the $E_8$ lattice CFT; modular invariance of the torus partition function forces
$\theta_{E_8}=E_4$ (on paper; Lean uses only the coefficient relation), so the number of massless currents ($8+240$) is fixed by
modularity alone.
\end{physics}

\begin{status}
Proved; axioms standard (re-run for \lean{E8Dim.finrank_e8Lie}).
\end{status}

For $D_4,E_6,E_7$ no such modular shortcut is available (the lattices are not unimodular), so the roots are counted by a verified
enumeration.

\begin{definition}[Pruned short-vector enumeration]\label{l23:def:enum}
A \emph{row} is a triple $(d,\ell,R)$: a pivot $d\in\Q$, coefficients $\ell$ on the later coordinates, and a radius $R\in\N$. For rows
$r_1,\dots,r_n$ the tail sum of an integer vector $v$ is $\sum_k d_k\big(v_k+\sum_{i>k}\ell_{ki}v_i\big)^2$. \lean{enum N rows} lists
integer vectors built from the last coordinate backwards, choosing each coordinate from the window of $2R+2$ consecutive integers starting at $\lfloor c\rfloor-R$, where
$c=-\sum_{i>k}\ell_{ki}v_i$ (so the window contains $[c-R,c+R]$, though it is not exactly centred), and keeping only partial vectors whose tail sum is $\le N$.
\end{definition}
\noindent File \leanpath{VOA/\allowbreak AffineGeneral/\allowbreak RootEnum.lean}.
\begin{leancode}
structure Row where
  d : ℚ
  l : List ℚ
  R : ℕ

def dotQ : List ℚ → List ℤ → ℚ
  | a :: as, x :: xs => a * x + dotQ as xs
  | _, _ => 0

def term (r : Row) : List ℤ → ℚ
  | [] => 0
  | x :: rest => r.d * ((x : ℚ) + dotQ r.l rest) ^ 2

def tailSum : List Row → List ℤ → ℚ
  | r :: rs, x :: xs => term r (x :: xs) + tailSum rs xs
  | _, _ => 0

def cands (r : Row) (rest : List ℤ) : List ℤ :=
  (List.range (2 * r.R + 2)).map fun i : ℕ => ⌊-dotQ r.l rest⌋ - (r.R : ℤ) + (i : ℤ)

def enum (N : ℚ) : List Row → List (List ℤ)
  | [] => [[]]
  | r :: rs => (enum N rs).flatMap fun rest =>
      ((cands r rest).map (· :: rest)).filter fun s => decide (tailSum (r :: rs) s ≤ N)

def Admissible (N : ℚ) (rows : List Row) : Prop := ∀ r ∈ rows, 0 < r.d ∧ N ≤ r.d * (r.R : ℚ) ^ 2
\end{leancode}

\begin{dictionary}
\lean{dotQ a x} is $\sum_ia_ix_i$ over the common length; the list \lean{r.l} of a row is ordered like the coordinates that
\emph{follow} it, so \lean{term r (x :: rest)} is $d\,(x+\sum_i\ell_i\,\mathrm{rest}_i)^2$ and \lean{cands r rest} is the window above.
\lean{List.flatMap} concatenates; \lean{filter} keeps elements satisfying a Boolean test; \lean{decide} turns a decidable proposition
into a Boolean. The definitions are \emph{computable}, which is the point: the kernel can evaluate \lean{enum}. \lean{Admissible}:
pivots positive and $dR^2\ge N$, so the window is wide enough. Same as the Fincke--Pohst algorithm.
\end{dictionary}

\textbf{Why this definition.} A naive enumeration of all integer vectors in a coordinate box containing the norm-$2$ vectors is too
slow for the kernel (Section~\ref{l23:sec:verification}); building vectors from the last coordinate backwards and discarding a partial
vector as soon as its partial sum of squares exceeds $N$ visits only branches that can still succeed. Completeness of this pruning is
Theorem~\ref{l23:thm:rootcounts}.

\begin{physics}
The output of \lean{enum 2 rows} is the list of charge vectors $\beta$ with $\langle\beta,\beta\rangle=2$, i.e.\ of the vertex operators
$e^{i\beta\cdot X}$ of conformal weight one: the number of extra massless currents beyond the $n$ Cartan currents.
\end{physics}

\begin{status}
Definitions (no proof obligation); computable, so the kernel can evaluate them.
\end{status}

\begin{theorem}[Completeness of the enumeration; root counts and dimensions]\label{l23:thm:rootcounts}
If the rows are admissible for $N$, every integer vector with tail sum $\le N$ appears in \lean{enum N rows}. Consequently, for the
$LDL^{\mathsf T}$ rows of $D_4,E_6,E_7,E_8$, the roots are counted as $24,72,126,240$, and $\dim\mathfrak{so}_8=28$, $\dim\mathfrak e_6=78$,
$\dim\mathfrak e_7=133$, $\dim\mathfrak e_8=248$.
\end{theorem}
\noindent Files \leanpath{VOA/\allowbreak AffineGeneral/\allowbreak RootEnum.lean},
\leanpath{VOA/\allowbreak AffineGeneral/\allowbreak RootCounts.lean}.
\begin{leancode}
theorem mem_enum (N : ℚ) (rows : List Row) (hadm : Admissible N rows) (s : List ℤ)
    (hlen : s.length = rows.length) (hs : tailSum rows s ≤ N) : s ∈ enum N rows := by

theorem card_normEq_eq {n : ℕ} (G : Matrix (Fin n) (Fin n) ℤ) (hG : G.IsSymm)
    (hPos : (ofCartan G hG).PosDef) (rows : List Row) (hlen : rows.length = n)
    (hadm : Admissible 2 rows)
    (hid : ∀ v : Fin n → ℤ, ((Matrix.toBilin' G v v : ℤ) : ℚ) = tailSum rows (List.ofFn v))
    (hnd : (rootList rows).Nodup) :
    (normEq (ofCartan G hG) hPos 2).card = (rootList rows).length := by

theorem e7_length : (rootList e7Rows).length = 126 := by decide +kernel

theorem card_roots_d4 : (normEq d4 d4_posDef 2).card = 24 := by

theorem card_roots_e6 : (normEq e6 e6_posDef 2).card = 72 := by

theorem card_roots_e7 : (normEq e7 e7_posDef 2).card = 126 := by

theorem card_roots_e8' : (normEq e8 e8_posDef 2).card = 240 := by

theorem finrank_d4Lie : Module.finrank ℂ d4Lie = 28 := by

theorem finrank_e6Lie : Module.finrank ℂ e6Lie = 78 := by

theorem finrank_e7Lie : Module.finrank ℂ e7Lie = 133 := by

theorem finrank_e8Lie' : Module.finrank ℂ e8Lie = 248 := by
\end{leancode}

\begin{dictionary}
\lean{rootList rows} filters \lean{enum 2 rows} to tail sum exactly $2$. Each \lean{card_roots_*} is \lean{card_normEq_eq} applied to
that type's rows, followed by the kernel-checked list length (as \lean{e7_length} for $E_7$). The names here are in the namespace
\lean{VOA.AffineGeneral.RootCounts}; the primes in \lean{card_roots_e8'} and \lean{finrank_e8Lie'} only avoid a clash with
\lean{E8Dim.card_roots_e8} and \lean{E8Dim.finrank_e8Lie} of Theorem~\ref{l23:thm:e8dim}, which state the same facts proved via
$\theta_{E_8}=E_4$; \lean{RootCounts.finrank_e6Lie} is the lattice algebra, not \lean{SCLie.finrank_e6Lie} of Theorem~\ref{l23:thm:g2f4e6}. The hypothesis \lean{hid} says the rows are an exact
$LDL^{\mathsf T}$ decomposition of the Cartan form; it is proved per type by rewriting with the \lean{sos_*} identities of
\leanpath{CartanLattices.lean} and \lean{ring}; \lean{hadm} and \lean{hnd} (no duplicates) and the lengths are checked by
\lean{decide +kernel}, i.e.\ by evaluation in Lean's kernel, which is part of the trusted base. The counts of the enumerated lists are
therefore not trusted computation outside Lean. The dimension theorems are $\operatorname{rank}+\#\text{roots}$
(Theorem~\ref{l23:thm:latticesimple}(3)): $4+24$, $6+72$, $7+126$, $8+240$. For $A_2$ the list length $6$ is checked
(\lean{a2_length}) but no theorem \lean{finrank_a2Lie = 8} is stated. Same as the textbook.
\end{dictionary}

\textbf{Proof on paper.} If the tail sum of $v$ is $\le N$, so is each term, hence
$d_k(v_k+c)^2\le N\le d_kR^2$, i.e.\ $|v_k+c|\le R$, and $v_k$ lies in the window; induction on the number of rows. The counts are then
finite evaluations. \qed

\begin{physics}
Counting the massless charged currents of the lattice CFT; the dimensions are the numbers of conserved currents, i.e.\ the dimensions of
the gauge groups $\mathrm{Spin}(8)$, $E_6$, $E_7$, $E_8$.
\end{physics}

\begin{status}
Proved; axioms standard for \lean{mem_enum}, \lean{finrank_e7Lie}, \lean{finrank_e8Lie'} and the four \lean{card_roots_*}
(re-run).
\end{status}

\subsection{The standard normalization}

\begin{theorem}[$\lxxiiiKil/(2\lxxiiihv)$ is the standard normalization]\label{l23:thm:normalization}
Let $L$ be generated by $n=\operatorname{rank}L$ simple roots $s$. Then the root vectors together with the simple coroots form a basis of
$(V_L)_1$, and $\lxxiiiKil(h_\alpha,h_\alpha)=\sum_{\beta}\langle\alpha,\beta\rangle^2$ (sum over roots). For $A_2,D_4,E_6,E_7,E_8$ this sum is
$12,24,48,72,120=4\lxxiiihv$ for every root $\alpha$, so $\big(\lxxiiiKil/(2\lxxiiihv)\big)(h_\alpha,h_\alpha)=2$.
\end{theorem}
\noindent Files \leanpath{VOA/\allowbreak AffineGeneral/\allowbreak LatticeKilling.lean},
\leanpath{VOA/\allowbreak AffineGeneral/\allowbreak LatticeNormalization.lean}.
\begin{leancode}
def v1Basis (hrank : rank L = n) : Module.Basis (Root L hPos ⊕ Fin n) ℂ (V1 L hEven hPos) :=
  basisOfTopLeSpanOfCardEqFinrank (Sum.elim (eV L hEven hPos) fun i => hV L hEven hPos (s i))
    (span_eV_hV_simple hEven s hsgen).ge
    (by rw [finrank_V1, Fintype.card_sum, Fintype.card_fin, hrank, Fintype.card_coe, add_comm])

theorem killing_hV_hV (hrank : rank L = n) (α : Root L hPos) :
    killingForm ℂ (V1 L hEven hPos) (hV L hEven hPos α) (hV L hEven hPos α)
      = ∑ β : Root L hPos, ((L.form α.1 β.1 : ℤ) : ℂ) ^ 2 := by

theorem e8_killing (β : Root e8 e8_posDef) :
    killingForm ℂ e8Lie (hV e8 e8_even e8_posDef β) (hV e8 e8_even e8_posDef β) = 120 := by

theorem e8_standard (β : Root e8 e8_posDef) :
    ((2 * 30 : ℂ)⁻¹ • killingForm ℂ e8Lie) (hV e8 e8_even e8_posDef β)
      (hV e8 e8_even e8_posDef β) = 2 := by
\end{leancode}

\begin{dictionary}
\lean{⊕} is the disjoint union of index sets; \lean{Sum.elim} combines two families. \lean{basisOfTopLeSpanOfCardEqFinrank} is Mathlib's
``a spanning family with as many elements as the dimension is a basis''. The variables \lean{s} and \lean{hsgen} (simple roots generating
$L$) are hypotheses of both statements. The value $\sum_\beta\langle\alpha,\beta\rangle^2$ is computed per type by enumerating roots
(\lean{sum_sq_form}, \lean{decide +kernel}). Why $\kappa(h_\alpha,h_\alpha)=2$ is the standard normalization: $h_\alpha$ is the coroot
($\alpha(h_\alpha)=2$), and for an invariant form $\kappa(h_\alpha,h_\alpha)=4/(\alpha,\alpha)_\kappa$, so the condition says all roots
(here all of the same length) have squared length $2$. Same as the textbook.
\end{dictionary}

\textbf{Proof on paper.} In the basis $\{e_\beta\}\cup\{h_{s_i}\}$, $\lxxiiiad h_\alpha$ is diagonal with eigenvalue $\langle\alpha,\beta\rangle$ on
$e_\beta$ and $0$ on the Cartan part, so $\tr(\lxxiiiad h_\alpha)^2=\sum_\beta\langle\alpha,\beta\rangle^2$. The identity
$\sum_{\beta\in\Phi}\langle\alpha,\beta\rangle^2=4\lxxiiihv$ for a root of length $2$ in a simply-laced system is classical (it is equivalent
to $\lxxiiiKil=2\lxxiiihv\cdot\kappa_{\mathrm{std}}$); here it is simply evaluated. \qed

\begin{physics}
In the normalization where long roots have length $2$, integrable levels are integers and the level-one lattice theory is
$L_1(\lxxiiig)$; the Deligne levels $-\lxxiiihv/6-1$ are meaningful only in this normalization. The theorem certifies that the
$k=-6$ of $L_{-6}(\mathfrak e_8)$ is the physicists' $k=-6$.
\end{physics}

\begin{status}
Proved; axioms standard for \lean{killing_hV_hV}, \lean{e8_standard} (re-run); for the other four \lean{*_standard} theorems, project record.
These theorems are the \emph{only} place where the Casimir number $h$ of the library is tied to the dual Coxeter number $\lxxiiihv$; the
identification is per type, never for general $\lxxiiig$.
\end{status}

\begin{remark}[Why the names are justified: the classification is not formalised]\label{l23:rem:type}
Lean proves that \lean{e8Lie} is a simple Lie algebra of dimension $248$; it does not prove that it is ``the'' $\mathfrak e_8$, because
the Cartan--Killing classification is not in Mathlib. On paper the identification is immediate in two independent ways. (a) By
Frenkel--Kac, $(V_L)_1$ for a simply-laced root lattice $L$ is the simple Lie algebra with root system the roots of $L$. (b) By
dimension alone, using the classification: the dimensions of simple Lie algebras are $n(n+2)$ ($A_n$), $n(2n+1)$ ($B_n,C_n$),
$n(2n-1)$ ($D_n$, $n\ge4$), $14,52,78,133,248$. The classical values near the exceptional ones are: $A_n$: $\dots,15,24,35,48,63,80,\dots,
120,143,\dots,224,255$; $B_n,C_n$: $10,21,36,55,78,105,136,\dots,210,253$; $D_n$: $28,45,66,91,120,153,\dots,231,276$. So $14$, $52$,
$133$ and $248$ are not classical, and $28$ occurs only for $D_4$; hence these dimensions determine $G_2$, $F_4$, $E_7$, $E_8$, $D_4$.
\textbf{But $78$ is ambiguous}: $\dim B_6=\dim C_6=\dim E_6=78$, and all three have rank $6$ and $72$ roots. For the lattice $\mathfrak e_6$,
(a) settles it. For the structure-constant $\mathfrak e_6$ of Section~\ref{l23:sec:sc} it is settled by the Killing value
$\lxxiiiKil(h,h)=48$ on the chosen coroot together with simplicity: for $B_6$ ($\lxxiiihv=11$) the possible values are $44$ and $88$, for
$C_6$ ($\lxxiiihv=7$) $28$ and $56$. This argument is on paper only. \emph{In Lean no Cartan type is ever identified}: the machine-checked
content is ``concrete simple Lie algebras of dimensions $14,52,78,133,248$, built as the exceptional ones'', and the three Lean models of
$\mathfrak e_6$ are not proved isomorphic (release audit 2026-10-27, finding M3).
\end{remark}

%======================================================================
\section{Concrete simple Lie algebras III: \texorpdfstring{$\mathfrak g_2$, $\mathfrak f_4$, $\mathfrak e_6$}{g2, f4, e6} from checked structure constants}\label{l23:sec:sc}
%======================================================================

The non-simply-laced algebras are not weight-one spaces of lattice VOAs of their own root lattices. The library builds them, and a
second $\mathfrak e_6$, directly from integer tables of structure constants, generated by untrusted Python scripts and re-checked in
Lean's kernel.

\subsection{Lie algebras from sparse tables}

\begin{definition}[Structure-constant algebra]\label{l23:def:scalg}
A \emph{table} $T$ assigns to each pair $(i,j)$ of indices $<N$ a finite list of pairs $(k,c)$, meaning $[b_i,b_j]=\sum c\,b_k$ with
$c\in\Z$. The \emph{checks} are: all indices $k<N$; antisymmetry $[b_i,b_j]+[b_j,b_i]=0$; and the cyclic Jacobi identity on basis triples,
coefficientwise. Given the checks, $\C^N$ with the bilinear extension of $T$ is a complex Lie algebra \lean{Alg N T hc}.
\end{definition}
\noindent File \leanpath{VOA/\allowbreak AffineGeneral/\allowbreak SCLie.lean}.
\begin{leancode}
abbrev Table : Type := List (List (List (ℕ × ℤ)))

def cycK (T : Table) (N : ℕ) (i j l n : Fin N) : ℂ :=
  (∑ m : Fin N, sc T j l m * sc T i m n) + (∑ m : Fin N, sc T l i m * sc T j m n) +
    ∑ m : Fin N, sc T i j m * sc T l m n

structure Checks (N : ℕ) (T : Table) : Prop where
  bound : boundOK T N = true
  anti : antiOK T N = true
  jac : ∀ i j l n : Fin N, cycK T N i j l n = 0

def Alg (N : ℕ) (T : Table) (_hc : Checks N T) : Type := Fin N → ℂ

theorem leibniz (x y z : Alg N T hc) :
    B hc x (B hc y z) = B hc (B hc x y) z + B hc y (B hc x z) := by

theorem jac_of_sorted (hs : ∀ i j l : Fin N, i < j → j < l → cycK T N i j l n = 0) (i j l : Fin N) :
    cycK T N i j l n = 0 := by
\end{leancode}

\begin{dictionary}
\lean{sc T i j k} is the integer $c^{ij}{}_k$ cast to $\C$; \lean{cycK T N i j l n} is the $n$-th coordinate of
$[b_i,[b_j,b_l]]+[b_j,[b_l,b_i]]+[b_l,[b_i,b_j]]$. \lean{boundOK}, \lean{antiOK} are Boolean functions evaluated by the kernel;
\lean{jac} is a \emph{proposition}, established by \lean{jac_of_sorted} from Boolean checks of sorted triples $i<j<l$ only (by
antisymmetry the cyclic sum is alternating, so sorted triples suffice, and triples with a repeated index vanish). The antisymmetry is
a hypothesis of \lean{jac_of_sorted} that the printed signature does not show: it is the section variable
\lean{ha : antiOK T N = true} (\lean{include ha}); the output coordinate \lean{n} is fixed throughout. \lean{Alg} is a type
synonym for $\C^N$, so that it does not inherit the pointwise bracket of functions; the instance \lean{LieAlgebra ℂ (Alg N T hc)} uses the
bilinear map \lean{B hc}, and \lean{leibniz} (from the basis case by trilinearity) is its Jacobi identity. The table is stored in a
sign-magnitude flat encoding (\lean{tableOf}, \lean{decZ}) for elaboration speed. Same as the textbook.
\end{dictionary}

\textbf{Proof on paper.} A bilinear antisymmetric bracket satisfying Jacobi on a basis satisfies it everywhere by trilinearity. \qed

\begin{physics}
A symmetry algebra specified by the simple-pole coefficients $f^{ab}{}_c$ of the current OPE \eqref{l23:eq:JJ} alone. The checks are the
consistency conditions of that OPE: antisymmetry is the symmetry of $J^a(z)J^b(w)$ under exchange, and Jacobi is associativity of the
triple OPE $J^aJ^bJ^c$. A table that passes them is a consistent set of current-current OPEs.
\end{physics}

\begin{status}
Proved; axioms standard.
\end{status}

\begin{theorem}[Simplicity from root tables]\label{l23:thm:sctables}
Let \lean{Alg N T hc} be a structure-constant algebra and $D$ root data: a list of root indices $\beta\mapsto b_{\mathrm{idx}(\beta)}$, a
map $-$ on root indices (no involution axiom is needed), an integer pairing table $p$, $r$ simple roots, an integer matrix $M$ expressing every coroot $h_\alpha=[e_\alpha,e_{-\alpha}]$ in
the simple coroots, a rational right inverse $Q$ of the simple pairing matrix, and a distance function on roots. If the finite checks
\lean{RootChecks} hold, the algebra is simple.
\end{theorem}

Before the Lean statement, here is what each ingredient is, on the smallest example.

\begin{example}[$\mathfrak{sl}_2$ as a structure-constant algebra, with its root data]\label{l23:ex:sl2table}
Take $N=3$, $b_0=e$, $b_1=H$, $b_2=f$. The table lists, for each ordered pair $(i,j)$, the nonzero coordinates $(k,c)$ of $[b_i,b_j]$:
\begin{center}\small
\begin{tabular}{c|ccc}
$[b_i,b_j]$ & $j=0$ ($e$) & $j=1$ ($H$) & $j=2$ ($f$)\\
\midrule
$i=0$ ($e$) & $[\,]$ & $[(0,-2)]$ ($=-2e$) & $[(1,1)]$ ($=H$)\\
$i=1$ ($H$) & $[(0,2)]$ ($=2e$) & $[\,]$ & $[(2,-2)]$ ($=-2f$)\\
$i=2$ ($f$) & $[(1,-1)]$ ($=-H$) & $[(2,2)]$ ($=2f$) & $[\,]$
\end{tabular}
\end{center}
In Lean this is the list of rows \lean{[[[], [(0, -2)], [(1, 1)]], [[(0, 2)], [], [(2, -2)]], [[(1, -1)], [(2, 2)], []]]}; a sparse vector
such as \lean{[(1, 1)]} means $1\cdot b_1$. The checks of Definition~\ref{l23:def:scalg} hold: indices are $<3$; the table is antisymmetric;
and Jacobi holds on the one sorted triple $(0,1,2)$: $[e,[H,f]]+[H,[f,e]]+[f,[e,H]]=[e,-2f]+[H,-H]+[f,-2e]=-2H+0+2H=0$.

\emph{Root data.} The roots are $\alpha$ and $-\alpha$, with root vectors $e=b_0$ and $f=b_2$: \lean{nr = 2}, \lean{rootIdx} $=(0\mapsto0,\
1\mapsto2)$, and \lean{neg} swaps them. The coroots are \emph{computed}, not supplied: $h_{\alpha}=[e,f]=H$ (sparse \lean{[(1, 1)]}) and
$h_{-\alpha}=[f,e]=-H$ (sparse \lean{[(1, -1)]}). The pairing table is read off from $[h_\gamma,e_\beta]=p(\gamma,\beta)e_\beta$:
$[H,e]=2e$, $[H,f]=-2f$, $[-H,e]=-2e$, $[-H,f]=2f$, so \lean{pairT} $=\begin{pmatrix}2&-2\\-2&2\end{pmatrix}$. There is $r=1$ simple root,
\lean{simple 0} $=\alpha$; \lean{M} expresses each coroot in the simple coroot: $h_\alpha=1\cdot h_\alpha$, $h_{-\alpha}=-1\cdot h_\alpha$; the
simple pairing matrix is $P=(2)$ and \lean{Q} $=(\tfrac12)$. Take \lean{base} $=\alpha$, \lean{dist} $=(0,1)$.

\emph{The checks, one by one.} \lean{inRange_h}: $h_{\pm\alpha}$ have index $1<3$. \lean{hbe}: e.g.\ $[h_\alpha,e_\alpha]-2e_\alpha=2e-2e=0$.
\lean{hhb}: $[H,\pm H]=0$. \lean{diag}: $p(\alpha,\alpha)=p(-\alpha,-\alpha)=2\ne0$. \lean{sep}: for $\alpha\ne-\alpha$, $\gamma=\alpha$ gives
$p(\alpha,\alpha)=2\ne-2=p(\alpha,-\alpha)$. \lean{symm}: \lean{pairT} is symmetric. \lean{M}: $h_{-\alpha}-(-1)h_\alpha=-H+H=0$. \lean{Q}: $2\cdot\frac12=1$.
\lean{conn}: $-\alpha$ has the neighbour $\alpha$ with smaller distance and $p(\alpha,-\alpha)\ne0$. \lean{span}: $b_0,b_2$ are root vectors,
and $b_1=H=h_{\alpha}$ (the disjunct ``$h_{s_i}-b_k=0$''). So Theorem~\ref{l23:thm:sctables} proves $\mathfrak{sl}_2$ simple. For
$\mathfrak e_6$ the same data have $78$ basis vectors, $72$ roots and $6$ simple roots, and the checks are kernel evaluations.
\end{example}

\noindent File \leanpath{VOA/\allowbreak AffineGeneral/\allowbreak SCSimple.lean}.
\begin{leancode}
structure RootData (N : ℕ) where
  nr : ℕ
  rootIdx : Fin nr → Fin N
  neg : Fin nr → Fin nr
  pairT : Fin nr → Fin nr → ℤ
  r : ℕ
  simple : Fin r → Fin nr
  M : Fin nr → Fin r → ℤ
  Q : Fin r → Fin r → ℚ
  base : Fin nr
  dist : Fin nr → ℕ

def eL (D : RootData N) (β : Fin D.nr) : List (ℕ × ℤ) := [((D.rootIdx β : ℕ), 1)]

def hL (D : RootData N) (α : Fin D.nr) : List (ℕ × ℤ) := brL T (eL D α) (eL D (D.neg α))

structure RootChecks (T : Table) (D : RootData N) : Prop where
  inRange_h : ∀ α, inRange N (hL T D α) = true
  hbe : ∀ α β, zeroL (brL T (hL T D α) (eL D β) ++ scale (-(D.pairT α β)) (eL D β)) = true
  hhb : ∀ α β, zeroL (brL T (hL T D α) (hL T D β)) = true
  diag : ∀ α, D.pairT α α ≠ 0
  sep : ∀ α β, α ≠ β → ∃ γ, D.pairT γ α ≠ D.pairT γ β
  symm : ∀ α β, D.pairT α β ≠ 0 → D.pairT β α ≠ 0
  M : ∀ α, zeroL (hL T D α ++ scale (-1)
    ((List.finRange D.r).flatMap fun i => scale (D.M α i) (hL T D (D.simple i)))) = true
  Q : ∀ i j, ∑ k, (D.pairT (D.simple i) (D.simple k) : ℚ) * D.Q k j = if i = j then 1 else 0
  conn : ∀ β, β = D.base ∨ ∃ γ, D.dist γ < D.dist β ∧ D.pairT γ β ≠ 0
  span : ∀ k : Fin N, (∃ β, D.rootIdx β = k) ∨
    ∃ i, zeroL (hL T D (D.simple i) ++ [((k : ℕ), 1)]) = true ∨
      zeroL (hL T D (D.simple i) ++ [((k : ℕ), -1)]) = true

theorem isSimple_of_tables (hD : RootChecks T D) [Nonempty (Fin D.nr)] :
    LieAlgebra.IsSimple ℂ (Alg N T hc) := by
\end{leancode}

\begin{dictionary}
Sparse vectors are lists of (index, integer) pairs; \lean{brL} is the sparse bracket, \lean{zeroL l} tests that all coefficients of $l$
vanish, \lean{scale} multiplies, \lean{++} concatenates (i.e.\ adds). Thus \lean{hbe} is $[h_\alpha,e_\beta]-p(\alpha,\beta)e_\beta=0$,
\lean{hhb} is $[h_\alpha,h_\beta]=0$, \lean{M} is $h_\alpha=\sum_iM_{\alpha i}h_{s_i}$, \lean{Q} is $PQ=1$ for the simple pairing matrix $P$,
\lean{conn} is connectivity by a distance function, and \lean{span}: every basis vector is a root vector or $\pm$ a simple coroot.
\lean{eL D β} is the root vector $e_\beta=b_{\mathrm{idx}(\beta)}$ and \lean{hL T D α} is $h_\alpha:=[e_\alpha,e_{\mathrm{neg}\,\alpha}]$, so the
hypothesis \lean{hneg} of Theorem~\ref{l23:thm:rootspace} holds by construction with $c_\alpha=1$; that is why \lean{neg} needs no axiom.
How the fields feed the criterion: \lean{hbe}, \lean{hhb}, \lean{diag}, \lean{sep} are its hypotheses of the same names;
\lean{inRange_h} (the sparse vector $h_\alpha$ has all indices $<N$) is a well-formedness condition that lets the sparse computations be
read as statements in $\C^N$; \lean{span} gives the spanning hypothesis; \lean{M} and \lean{Q} give the centre condition
(\lean{RootSpace.center_of_simple}). Connectivity takes two fields: \lean{conn} only gives a path \emph{from} the base root
\lean{D.base} to every root, and \lean{symm} (a nonzero pairing has nonzero transpose) reverses such paths, so that any two roots are joined
through the base, which is the two-way \lean{Relation.ReflTransGen} hypothesis \lean{conn} of Theorem~\ref{l23:thm:rootspace}. Root vectors
are basis vectors, hence nonzero. Same as the criterion.
\end{dictionary}

\textbf{Proof on paper.} Theorem~\ref{l23:thm:rootspace}. For the centre: an element $y=\sum_id_ih_{s_i}$ (by \lean{M}) commuting with all
$e_{s_k}$ satisfies $\sum_id_iP_{ik}=0$, so $d=dPQ=0$. \qed

\begin{physics}
\lean{pairT} is the table of charges of the charged currents $e_\beta$ under the Cartan currents $h_\alpha$ (the eigenvalues of
$h_{\alpha,(0)}$), i.e.\ the weights read off from the simple poles $h_\alpha(z)e_\beta(w)\sim p(\alpha,\beta)e_\beta(w)/(z-w)$. The checks say
that these charges are those of an irreducible root system: distinct currents carry distinct charges, and the charges link all
currents, so the symmetry algebra does not factorize.
\end{physics}

\begin{status}
Proved; axioms standard (re-run for \lean{SCLie.isSimple_of_tables}).
\end{status}

\subsection{The three algebras}

\begin{theorem}[$\mathfrak g_2$, $\mathfrak f_4$, $\mathfrak e_6$]\label{l23:thm:g2f4e6}
The tables \lean{g2T} ($N=14$), \lean{f4T} ($N=52$), \lean{e6T} ($N=78$) pass all checks and root checks; the algebras are simple, of
dimensions $14,52,78$. On a long coroot $h$ the Killing form is $16,36,48=4\lxxiiihv$ with $\lxxiiihv=4,9,12$, so $\lxxiiiKil/(2\lxxiiihv)$ gives
$\kappa(h,h)=2$.
\end{theorem}
\noindent Files \leanpath{VOA/\allowbreak AffineGeneral/\allowbreak G2Lie.lean},
\leanpath{VOA/\allowbreak AffineGeneral/\allowbreak F4Lie.lean}, \leanpath{VOA/\allowbreak AffineGeneral/\allowbreak E6Lie.lean}.
\begin{leancode}
abbrev g2Lie : Type := Alg 14 g2T g2Checks

instance g2Lie_isSimple : LieAlgebra.IsSimple ℂ g2Lie :=
  haveI : Nonempty (Fin g2Data.nr) := ⟨⟨0, by decide⟩⟩
  isSimple_of_tables g2Checks g2Data g2RootChecks

theorem finrank_g2Lie : Module.finrank ℂ g2Lie = 14 := finrank_alg g2Checks

theorem g2_killing : killingForm ℂ g2Lie (hbV g2Checks g2Data g2Long)
    (hbV g2Checks g2Data g2Long) = 16 := by

theorem g2_standard : ((2 * 4 : ℂ)⁻¹ • killingForm ℂ g2Lie) (hbV g2Checks g2Data g2Long)
    (hbV g2Checks g2Data g2Long) = 2 := by

theorem g2_coroot : ⁅hbV g2Checks g2Data g2Long, eV g2Checks g2Data g2Long⁆
    = (-2 : ℂ) • eV g2Checks g2Data g2Long := by

abbrev f4Lie : Type := Alg 52 f4T f4Checks

instance f4Lie_isSimple : LieAlgebra.IsSimple ℂ f4Lie :=
  haveI : Nonempty (Fin f4Data.nr) := ⟨⟨0, by decide⟩⟩
  isSimple_of_tables f4Checks f4Data f4RootChecks

theorem finrank_f4Lie : Module.finrank ℂ f4Lie = 52 := finrank_alg f4Checks

theorem f4_killing : killingForm ℂ f4Lie (hbV f4Checks f4Data f4Long)
    (hbV f4Checks f4Data f4Long) = 36 := by

theorem f4_standard : ((2 * 9 : ℂ)⁻¹ • killingForm ℂ f4Lie) (hbV f4Checks f4Data f4Long)
    (hbV f4Checks f4Data f4Long) = 2 := by

abbrev e6Lie : Type := Alg 78 e6T e6Checks

instance e6Lie_isSimple : LieAlgebra.IsSimple ℂ e6Lie :=
  haveI : Nonempty (Fin e6Data.nr) := ⟨⟨0, by decide⟩⟩
  isSimple_of_tables e6Checks e6Data e6RootChecks

theorem finrank_e6Lie : Module.finrank ℂ e6Lie = 78 := finrank_alg e6Checks

theorem e6_killing : killingForm ℂ e6Lie (hbV e6Checks e6Data e6Long)
    (hbV e6Checks e6Data e6Long) = 48 := by

theorem e6_standard : ((2 * 12 : ℂ)⁻¹ • killingForm ℂ e6Lie) (hbV e6Checks e6Data e6Long)
    (hbV e6Checks e6Data e6Long) = 2 := by
\end{leancode}

\begin{dictionary}
All names in this block live in the namespace \lean{VOA.AffineGeneral.SCLie}; in particular this \lean{finrank_e6Lie} is
\lean{SCLie.finrank_e6Lie} (the $78$-dimensional table algebra), not \lean{RootCounts.finrank_e6Lie} of Theorem~\ref{l23:thm:rootcounts}
(the lattice algebra). \lean{g2Data}, \lean{f4Data}, \lean{e6Data} here are \lean{RootData N} (the root tables), \emph{not} Lie data (see the
warning box in Section~\ref{l23:warn:models}). \lean{hbV … g2Long} is the bracket $[e_\alpha,e_{-\alpha}]$ for the root with index $1$, which
the file names \lean{g2Long}; likewise \lean{f4Long}, \lean{e6Long}. \emph{Sign:} \lean{g2_coroot} (and \lean{f4_coroot}, \lean{e6_coroot},
same form) says $[\,\mathtt{hbV},e_\alpha]=-2\,e_\alpha$, so \lean{hbV} is the \emph{negative} of the standard coroot $h_\alpha$ (for which
$\alpha(h_\alpha)=2$). The Killing values are quadratic, hence insensitive to this sign. Nothing in Lean certifies that \lean{g2Long} is long
(Lean checks $[\,\mathtt{hbV},e]=-2e$ and the Killing value, but does not compare the chosen root with the other roots), so the
``standard normalization'' conclusion rests partly on the following paper argument (release audit 2026-10-27, finding m2): the computed Killing value $16=4\cdot4$ is the long-root value for $G_2$ (the short-root value would be $48$); similarly $36=4\cdot9$
and $48=4\cdot12$ are the long-root values for $F_4$ and $E_6$. Thus \lean{f4_standard} and \lean{e6_standard} certify that
$\lxxiiiKil/18$ and $\lxxiiiKil/24$ give these coroots squared length $2$. Where the
tables come from: $\mathfrak g_2$ is the fixed subalgebra of $\mathfrak{so}_8$ under triality and $\mathfrak f_4$ the fixed subalgebra of $\mathfrak e_6$
under the diagram involution, written in orbit-sum bases (generator \texttt{aff\_fold.py}); $\mathfrak e_6$ is the Chevalley basis of
\texttt{casimir\_cert/cert\_E6.json} with Frenkel--Kac signs. The generators are untrusted; only the tables enter Lean, and every property used
is re-checked there. The identification of type is Remark~\ref{l23:rem:type}.
\end{dictionary}

\textbf{Proof on paper.} Theorem~\ref{l23:thm:sctables} with the checked tables; $\dim\C^N=N$; the Killing form of a sparse vector is the
finite sum $\sum_k\big[h,[h,b_k]\big]_k$ (\lean{killing_vec}), evaluated by \lean{decide +kernel}. \qed

\begin{physics}
$G_2$ and $F_4$ are the non-simply-laced members of the Deligne series. $\mathfrak g_2$ is the subalgebra of $\mathfrak{so}_8$ fixed by
the triality outer automorphism (the order-three symmetry of the $D_4$ Dynkin diagram), and $\mathfrak f_4$ the subalgebra of $\mathfrak e_6$
fixed by the diagram involution, which in physics is charge conjugation of the $E_6$ currents ($\mathbf{27}\leftrightarrow\overline{\mathbf{27}}$).
\end{physics}

\begin{status}
Proved; axioms standard for \lean{g2Lie_isSimple}, \lean{f4Lie_isSimple}, \lean{e6Lie_isSimple}, \lean{g2_standard},
\lean{g2_coroot}, \lean{f4_killing}, \lean{f4_standard}, \lean{e6_killing}, \lean{e6_standard} (re-run).
\end{status}

\begin{warning}[title={Trap: \texttt{hbV} is minus the standard coroot}]\label{l23:warn:hbvsign}
In \leanpath{SCSimple.lean} the docstring of \lean{hL} calls $[e_\alpha,e_{\mathrm{neg}\,\alpha}]$ the ``sparse coroot'', and the module
docstrings of \leanpath{F4Lie.lean} and \leanpath{E6Lie.lean} speak of the Killing value ``on a long coroot''. But \lean{g2_coroot},
\lean{f4_coroot}, \lean{e6_coroot} prove $[\mathtt{hbV},e_\alpha]=-2e_\alpha$: with the tables' sign conventions this vector is $-h_\alpha$.
Nothing proved is affected (the Killing values and the \lean{*_standard} theorems are quadratic in it), but a user who feeds \lean{hbV}
into a formula linear in the coroot gets the wrong sign. Listed in Section~\ref{l23:sec:defects}.
\end{warning}

\paragraph{How the tables are checked.} Each table check (index bounds, antisymmetry, the cyclic Jacobi identity on all sorted basis
triples --- $\binom{14}{3}=364$, $\binom{52}{3}=22100$ and $\binom{78}{3}=76076$ triples, each against all $N$ output coordinates ---
and the fields of \lean{RootChecks}) is a statement ``this Boolean function of the table evaluates to \lean{true}'', proved by
\lean{decide +kernel}. A single such declaration over a whole table exceeds the kernel's memory, so the checks are split into many small
declarations in generated files (one row or block of rows per declaration), and a separate file assembles the propositional
conclusions. This is the only reason for the large number of files; what is trusted is the kernel evaluation, not the scripts that
generated the tables. The file map, the costs and the dead ends that forced this design are recorded in
Section~\ref{l23:sec:verification}.

\subsection{Lie data with an explicit invariant form}

For the explicit presentations of Chapter~\ref{ch:L24} one wants Lie data whose basis \emph{is} the table basis and whose form has integer entries.

\begin{definition}[Lie data from sparse tables]\label{l23:def:lieDataOf}
Given an integer form table $G_m$ (a list of entries $(i,j,c)$), dual vectors $d_j=\frac1{D_j}\mathrm{vec}(\ell_j)$, an integer $L\ne0$ and
$c\in\Z$, and Boolean checks that $G_m$ is symmetric, that every basis vector is $\kappa$-skew ($\kappa([a,y],z)+\kappa(y,[a,z])=0$), that
$\kappa(b_i,d_j)=\delta_{ij}$, and that $\sum_i\frac{L}{D_i}[b_i,[\mathrm{vec}\,\ell_i,b_x]]=cL\,b_x$ for every $x$, one obtains Lie data with
$h=c/2$.
\end{definition}
\noindent Files \leanpath{VOA/\allowbreak AffineGeneral/\allowbreak SCLieData.lean},
\leanpath{VOA/\allowbreak AffineGeneral/\allowbreak E6LieData.lean}.
\begin{leancode}
def Skew (a : Alg N T hc) : Prop :=
  ∀ y z : Alg N T hc, formG hc Gm ⁅a, y⁆ z + formG hc Gm y ⁅a, z⁆ = 0

def lieDataOf (hs : symmPerm Gm = true) (hsk : ∀ c : Fin N, Skew hc Gm (basis hc c))
    (dl : ℕ → List (ℕ × ℤ)) (D : ℕ → ℤ) (hp : pairOK N Gm dl D = true) (L c : ℤ) (hL : L ≠ 0)
    (hcas : ∀ x : Fin N, casOK T N dl D L c x = true) : LieData (formG hc Gm) (Fin N) where

def e6LieData : LieData (formG e6Checks e6Gm) (Fin 78) :=
  lieDataOf e6Checks e6Gm e6_gInRange e6_symm e6Skew_all e6dl e6D e6_pair 3 24 (by norm_num) e6_cas

theorem e6LieData_h : e6LieData.h = 12 := by
\end{leancode}

\begin{dictionary}
\lean{formG hc Gm} is the bilinear form with Gram entries from \lean{Gm}. Skewness of all basis vectors is derived from skewness of the 12
Chevalley generators (\lean{E6LieDataSkew*}) propagated along single-term brackets (\lean{skew_of_entry}), since the generators generate.
The Casimir check is \lean{casOK … 3 24}: $L=3$ clears denominators, $c=24=2h$. The docstring says $e6Gm$ is the form with long roots of length
$2$, i.e.\ $\lxxiiiKil/24$; no Lean theorem states $\mathrm{formG}\,e6Gm=\lxxiiiKil/24$. On paper it follows: on a simple algebra every invariant form
is $\lambda\lxxiiiKil$, the Casimir of $\lambda\lxxiiiKil$ is $\lambda^{-1}\cdot1$ (by \lean{killing_h}), and the check gives $\lambda^{-1}=24$. The same
construction gives \lean{G2Lie.g2CertData} (form $\lxxiiiKil/8$, $h=4$) and \lean{F4Lie.f4CertData} ($\lxxiiiKil/18$, $h=9$) in
\leanpath{VOA/\allowbreak DeligneAudit/\allowbreak G2LieData.lean} and \leanpath{F4LieData.lean}, with checks in \leanpath{F4LieTab.lean} and
\leanpath{F4LieChecks*.lean}.
\end{dictionary}

\textbf{Proof on paper.} Invariance of a symmetric form follows from skewness of $\lxxiiiad(a)$ for all $a$; the pairing and Casimir checks are
the fields of Definition~\ref{l23:def:liedata} in coordinates. \qed

\begin{physics}
This is the current algebra with a \emph{chosen} integral normalization of the two-point function, the form in which the Deligne singular
vectors and Joseph relations of Chapter~\ref{ch:L24} are written.
\end{physics}

\begin{status}
Proved; axioms standard for \lean{lieDataOf}, \lean{e6LieData}, \lean{G2Lie.g2CertData} (re-run).
\end{status}

%======================================================================
\section{The Deligne instances}\label{l23:sec:deligne}
%======================================================================

The higher-rank capstone of Chapter~\ref{ch:L24} (\leanpath{VOA/\allowbreak Examples/\allowbreak DeligneHigherRankConditionalCapstone.lean})
is stated for a Deligne case \lean{case : HigherRankDeligneCase} (one of \lean{A2, G2, D4, F4, E6, E7, E8}), a VOA $W$, and a
\emph{current $C_2$ presentation}: weight-one generators indexed by \lean{Fin case.currentDimension} ($8,14,28,52,78,133,248$) and an ideal
of polynomial relations that vanish in Li degree zero. This section is the interface between this chapter and
Chapter~\ref{ch:L24}, so we first explain, informally and in one place, the notions it uses; the precise definitions are in
Chapter~\ref{ch:L11} (Section~\ref{l11:sec:jets}) and Chapter~\ref{ch:L24} (Definitions~\ref{l24:def:lpcf} and~\ref{l24:def:packages}).

\paragraph{The notions, informally.}
\begin{itemize}
\item \emph{Li degree zero; the $C_2$-algebra.} For a VOA $W$, Li's filtration $F^0W\supseteq F^1W\supseteq\cdots$ measures ``how many
  derivatives'' a state contains; its associated graded $\operatorname{gr}^FW$ is a commutative algebra, and its degree-zero part is Zhu's
  $C_2$-algebra $R_W=W/C_2(W)$, $C_2(W)=\operatorname{span}\{u_{(-2)}v\}$ (Chapter~\ref{ch:L11}, Theorem~\ref{l11:thm:f1}). In $R_W$ only the
  ``leading order'' of a state survives: derivatives of fields are set to zero. Its spectrum (as a variety) is the \emph{associated
  variety} of $W$; for $L_k(\lxxiiig)$ the currents generate $R_W$, so $R_W=\C[X_1,\dots,X_n]/\ker$ with $X_i\mapsto[J^i_{-1}\lxxiiivac]$.
\item \emph{A presentation} lists generators of $R_W$ (here: the $n$ currents) and an ideal $I$ of relations known to hold among them.
\item \emph{Jets and the arc space.} For a commutative algebra $R=\C[X]/I$, the \emph{jet algebra} $J_\infty R$ is the polynomial ring in
  infinitely many variables $X_i^{(j)}$ ($j\ge0$, ``the $j$-th derivative of $X_i$'') modulo all derivatives $\partial^mf$ of the relations $f\in I$,
  where $\partial X_i^{(j)}=X_i^{(j+1)}$. Geometrically, it is the coordinate ring of the \emph{arc space} of $\operatorname{Spec}R$ (formal
  curves in the variety). There is a canonical map, the \emph{jet map}, $J_\infty(\C[X]/I)\to\operatorname{gr}^FW$ sending $X_i^{(j)}$ to the
  class of the divided derivative $\frac1{j!}L_{-1}^jJ^i_{-1}\lxxiiivac=J^i_{-1-j}\lxxiiivac$ of the current (Chapter~\ref{ch:L11}, Definition~\ref{l11:def:jetmap}).
\item \emph{Current-jet surjectivity} (\lean{CurrentJetSurjectivity}) says the jet map is onto: every state is, to leading order, a
  normally ordered polynomial in the currents and their derivatives (the currents ``strongly generate'' $W$).
\item \emph{Classical freeness} (\lean{CaseClassicallyFree}) says the jet map is moreover injective, hence an isomorphism
  $J_\infty(\C[X]/I)\cong\operatorname{gr}^FW$: the only relations in the whole VOA are those forced, by differentiation, from the relations $I$
  of the $C_2$-algebra.
\item \emph{The two remaining inputs} are a way to prove injectivity by counting. \lean{GeometricCarrierEmbedding} (\lean{hgeom}) gives, for
  every weight $N$ and Li degree $p$, an injective map from the dual of the $(p,N)$ piece of $J_\infty(\C[X]/I)$ into a finite-dimensional
  ``carrier'' space $\mathrm{Carrier}_{p,N}$, so that $\dim J_\infty(\C[X]/I)_N\le\sum_p\dim\mathrm{Carrier}_{p,N}$.
  \lean{OrdinaryCarrierCharacter} (\lean{hchar}) says $\sum_{p\le N}\dim\mathrm{Carrier}_{p,N}=\dim W_N$. Together with surjectivity, which gives
  $\dim W_N=\dim(\operatorname{gr}^FW)_N\le\dim J_\infty(\C[X]/I)_N$, the inequalities are squeezed into equalities, and a surjection between spaces
  of equal finite dimension is an isomorphism (Chapter~\ref{ch:L24}, Theorem~\ref{l24:thm:squeeze}).
\end{itemize}
\emph{A toy case.} For the universal algebra $V^k(\lxxiiig)$ itself, take $I=0$: by PBW, $\operatorname{gr}^FV^k(\lxxiiig)$ is the polynomial ring
in the symbols of all $J^a_{-1-j}\lxxiiivac$, which is exactly $J_\infty\C[X_1,\dots,X_n]$; so $V^k(\lxxiiig)$ is classically free with no
relations at all. For $L_k(\lxxiiig)$ at the Deligne level, relations appear already in weight two (the null vectors of
Theorem~\ref{l23:thm:singular}), and the question answered in Chapters~\ref{ch:L24}--\ref{ch:L27} is whether the quadratic relations
$I_2^{\mathrm{ns}}$ coming from them (the Joseph relations) generate everything.

\begin{definition}[Current $C_2$ presentation]\label{l23:def:c2pres}
Let $W$ be a VOA on $V$ and $n=\dim\lxxiiig$ for the case. A \emph{current $C_2$ presentation} of $W$ consists of vectors
$g_1,\dots,g_n\in W_1$ and an ideal $I\subseteq\C[X_1,\dots,X_n]$ contained in the kernel of the algebra map
$X_i\mapsto[g_i]\in\operatorname{gr}^F_0W=R_W$ (the class of $g_i$ in Li degree zero, i.e.\ in Zhu's $C_2$-algebra).
\end{definition}
\noindent File \leanpath{VOA/\allowbreak Examples/\allowbreak DeligneHigherRankConditionalCapstone.lean}.
\begin{leancode}
variable {V : Type uV} [AddCommGroup V] [Module ℂ V]
variable (case : HigherRankDeligneCase) (W : VertexOperatorAlgebra V)

structure CurrentC2Presentation where
  gen : Fin case.currentDimension → V
  gen_weight_one : ∀ i, gen i ∈ W.grading (1 : ℤ)
  ideal : Ideal (MvPolynomial (Fin case.currentDimension) ℂ)
  relations : ideal ≤ RingHom.ker
      (MvPolynomial.aeval fun i =>
        liDegreeZeroClass W 1 ⟨gen i, gen_weight_one i⟩).toRingHom

theorem CurrentC2Presentation.exact_kernel_of_caseClassicallyFree
    {P : CurrentC2Presentation case W} (hfree : CaseClassicallyFree case W P) :
    RingHom.ker
        (MvPolynomial.aeval fun i =>
          liDegreeZeroClass W 1 ⟨P.gen i, P.gen_weight_one i⟩).toRingHom =
      P.ideal :=
\end{leancode}

\begin{dictionary}
\lean{case}, \lean{W} (and the carrier \lean{V}) are the section variables printed first. \lean{case.currentDimension} is
$\dim\lxxiiig$ for the case ($8,14,28,52,78,133,248$). \lean{W.grading (1 : ℤ)} is $W_1$. \lean{MvPolynomial (Fin n) ℂ} is
$\C[X_1,\dots,X_n]$, \lean{MvPolynomial.aeval f} the algebra map sending $X_i$ to $f(i)$, \lean{.toRingHom} forgets $\C$-linearity, and
\lean{RingHom.ker} is its kernel. \lean{liDegreeZeroClass W 1 ⟨v, hv⟩} is the class of a weight-one vector $v$ in Li degree zero of Li's
associated graded algebra (Chapter~\ref{ch:L11}). The field \lean{relations} is the containment $I\le\ker$: a presentation is only required to
list \emph{some} relations.
The predicates used below are Lean structures of Chapter~\ref{ch:L24} (Definition~\ref{l24:def:packages}), explained informally at the
start of this section: \lean{CurrentJetSurjectivity} (the jet map $J_\infty(\C[X]/I)\to\operatorname{gr}^FW$ is surjective),
\lean{GeometricCarrierEmbedding} (an injective map from the dual of each bidegree piece of the presented jets into a carrier space),
\lean{OrdinaryCarrierCharacter} (the carrier dimensions add up to $\dim W_N$) and \lean{CaseClassicallyFree} (the presented
jet map is an isomorphism). \lean{exact_kernel_of_caseClassicallyFree} (Chapter~\ref{ch:L24}, Theorem~\ref{l24:thm:kereq}) says: if the case is classically free for the presentation $P$, then $I$ is the \emph{whole} kernel. Same as the textbook
notion of a presentation of $R_W$ by generators, with exactness not assumed.
\end{dictionary}

\textbf{Proof on paper.} Definition; for the theorem: classical freeness makes the presented jet map injective, and its degree-zero
part is $\C[X]/I\to R_W$, so $\ker\subseteq I$; the field \lean{relations} gives $I\subseteq\ker$. \qed

\begin{physics}
The currents $g_i$ generate the Higgs-branch chiral ring $R_W$ (in four-dimensional language, the moment maps of the flavour symmetry), and
$I$ is a list of chiral-ring relations among them; an exact presentation is the statement ``the Higgs branch is cut out by these
equations''.
\end{physics}

\begin{status}
Definition; \lean{exact_kernel_of_caseClassicallyFree} proved, \texttt{\#print axioms} standard (re-run).
\end{status}

Here we record what \leanpath{AffineGeneral/} supplies.

\begin{theorem}[Two capstone inputs for every case]\label{l23:thm:deligneinstance}
For every Deligne case, Lie data $D$ indexed by \lean{Fin case.currentDimension}, and $k+h\ne0$, on $W=\lxxiiiL$:
(1) the current states $[J^a_{-1}\lxxiiivac]$, together with the \emph{full} kernel of $X_a\mapsto[J^a_{-1}\lxxiiivac]\in R_{L_k}$ as ideal,
form a current $C_2$ presentation; (2) every step of Li's filtration of $\lxxiiiV$ and of $\lxxiiiL$ is spanned by words in current modes,
hence the jet map is surjective (\lean{CurrentJetSurjectivity}); (3) \emph{if} one is given a family of finite-dimensional carrier spaces
with an injective carrier embedding (\lean{hgeom}) and the carrier character identity (\lean{hchar}), then the jet map is an isomorphism,
i.e.\ $L_k(\lxxiiig)$ is classically free for this presentation.
\end{theorem}
\noindent File \leanpath{VOA/\allowbreak AffineGeneral/\allowbreak DeligneInstance.lean}.
\begin{leancode}
variable {𝔤 : Type*} [LieRing 𝔤] [LieAlgebra ℂ 𝔤] [FiniteDimensional ℂ 𝔤]
variable {κ : LinearMap.BilinForm ℂ 𝔤}
variable (case : HigherRankDeligneCase) (D : LieData κ (Fin case.currentDimension))
variable (k : ℂ) (hk : k + D.h ≠ 0)

theorem universal_genWordsSpan :
    GeneratorWordsSpanLiFiltration (affineVOA D k hk) (ugen case D k) (fun _ => 1) := by

def presentation : CurrentC2Presentation case (simpleVOA D k hk) where
  gen := gen case D k hk
  gen_weight_one := gen_weight_one case D k hk
  ideal := RingHom.ker
    (MvPolynomial.aeval fun i =>
      liDegreeZeroClass (simpleVOA D k hk) 1 ⟨gen case D k hk i, gen_weight_one case D k hk i⟩).toRingHom
  relations := le_rfl

theorem genWordsSpan :
    GeneratorWordsSpanLiFiltration (simpleVOA D k hk) (presentation case D k hk).gen (fun _ => 1) :=
  generatorWordsSpanLiFiltration_quotient (QuotVOA.maxGradedIdeal (affineVOA D k hk))
    (QuotVOA.maxGradedIdeal_isIdeal _) (QuotVOA.maxGradedIdeal_isGradedIdeal _)
    (universal_genWordsSpan case D k hk)

theorem currentJetSurjectivity : CurrentJetSurjectivity case (simpleVOA D k hk) (presentation case D k hk) :=
  DeligneAudit.currentJetSurjectivity_of_genWordSpan (presentation case D k hk)
    (genWordsSpan case D k hk)

theorem classicallyFree_of_two_inputs
    (Carrier : ℕ → ℕ → Type)
    [∀ p N, AddCommGroup (Carrier p N)]
    [∀ p N, Module ℂ (Carrier p N)]
    [∀ p N, FiniteDimensional ℂ (Carrier p N)]
    (hgeom : GeometricCarrierEmbedding case (simpleVOA D k hk) (presentation case D k hk) Carrier)
    (hchar : OrdinaryCarrierCharacter (simpleVOA D k hk) Carrier) :
    CaseClassicallyFree case (simpleVOA D k hk) (presentation case D k hk) :=
\end{leancode}

\begin{dictionary}
The namespace is \lean{VOA.AffineGeneral.Deligne}; \lean{case}, \lean{D}, \lean{k}, \lean{hk} are the section variables printed at
the top (with the instance \lean{[FiniteDimensional ℂ 𝔤]}). \lean{GeneratorWordsSpanLiFiltration W g w} (Chapter~\ref{ch:L11}) says that each
step of Li's filtration of $W$ is spanned by words in the modes of the generators \lean{g} of weights \lean{w}; \lean{ugen case D k a} is the
current state $J^a_{-1}\lxxiiivac$ in $\lxxiiiV$, so \lean{universal_genWordsSpan} is claim (2) for $\lxxiiiV$ and \lean{genWordsSpan}, obtained from it
by passing to the quotient by the maximal graded ideal, is claim (2) for $\lxxiiiL$. \lean{gen case D k hk a}
is $[J^a_{-1}\lxxiiivac]\in\lxxiiiL$. \lean{RingHom.ker} is the kernel and \lean{le_rfl} proves $I\le I$. \lean{Carrier p N} is a family of
finite-dimensional spaces (the ``geometric carrier'' of Chapter~\ref{ch:L24}); \lean{hgeom}, \lean{hchar} are the two named inputs. Nothing ties
$D$ to the case except its index type: the theorem applies to \emph{any} Lie algebra of dimension \lean{case.currentDimension}; the
per-type files choose the right one.
\end{dictionary}

\begin{warning}[title={Trap: The general presentation uses the full kernel}]\label{l23:warn:fullkernel}
\lean{Deligne.presentation} takes as ideal the \emph{entire} kernel of the current-coordinate map, so its \lean{relations} field is
\lean{le_rfl}. This presentation therefore does not \emph{display} which relations hold; it only fixes generators. The statement
\lean{CaseClassicallyFree case W (presentation …)} is not weak: by \lean{exact_kernel_of_caseClassicallyFree}
(Definition~\ref{l23:def:c2pres}) classical freeness for \emph{any} presentation forces its ideal to be the full kernel, so for this
presentation it is exactly the full capstone statement. What the full-kernel choice prevents is \emph{discharging} the two inputs
\lean{hgeom}, \lean{hchar}: they cannot be checked without knowing the kernel explicitly. The mathematical content of
classical freeness for an explicit ideal (the Joseph-type ideal $I_2^{\mathrm{ns}}$ generated by weight-two singular vectors) is in the explicit
presentations of Chapter~\ref{ch:L24} (\lean{E6.explicitPresentation} etc.), which use Theorem~\ref{l23:thm:singular} and
Proposition~\ref{l23:prop:clsG}. The theorem \lean{currentJetSurjectivity} is genuine content (jet surjectivity does not depend on the ideal
beyond the generators).
\end{warning}

\textbf{Proof on paper.} (2) Let $C_q$ be the span of current words of Li level at least $q$ (the level of a word counts its
derivatives, i.e.\ $\sum(n_i-1)$ for $J^{a_1}_{-n_1}\cdots J^{a_r}_{-n_r}\lxxiiivac$). The \emph{tracked estimate}
(\lean{GeneratorModeTracked}, \leanpath{VOA/\allowbreak Filtered/\allowbreak GeneratorCurrentLiSpanning.lean}) says that the $m$-th mode of a
generator maps $C_q$ into $C_{\max(q-m-1,0)}$; for currents it follows from the commutator \eqref{l23:eq:affine}
(\lean{generatorModeTracked_of_current_comm}), since commuting $J^a_m$ past a word produces only words with the same total mode. By the
Borcherds iterate formula the estimate extends from generators to all states; as current words span $\lxxiiiV$, this shows that the
$p$-th step of Li's filtration is contained in $C_p$ (the reverse inclusion is clear), i.e.\ it is spanned by current words; spanning of Li's filtration by generator words descends to quotients (\lean{generatorWordsSpanLiFiltration_quotient}); a spanning
statement implies jet surjectivity (Chapter~\ref{ch:L11}). (1) is \lean{le_rfl}: the kernel contains itself. (3) is the squeeze
explained at the start of this section --- the capstone of Chapter~\ref{ch:L24} (Theorem~\ref{l24:thm:capstone}) with (2) plugged in as its
surjectivity input. \qed

\begin{physics}
Classical freeness says that the chiral algebra is, as a filtered object, the quantization of the arc space of its Higgs branch: no extra
relations appear in the chiral algebra beyond those forced by the Higgs-branch chiral ring. In general this is \emph{not} true (there are
known counterexamples among chiral algebras), and for the Deligne series it is a conjecture-level statement in the physics literature;
this chapter proves nothing towards it beyond the reduction to \lean{hgeom}, \lean{hchar}, and the status of those inputs is the subject of
Chapters~\ref{ch:L24}--\ref{ch:L27}. Jet surjectivity --- every state is a normally
ordered product of currents and their derivatives --- is the ``easy half''; it says the currents strongly generate the chiral algebra.
\end{physics}

\begin{status}
Proved; \texttt{\#print axioms} standard for \lean{Deligne.currentJetSurjectivity}, \lean{Deligne.classicallyFree_of_two_inputs},
\lean{Deligne.universal_genWordsSpan}, \lean{Deligne.genWordsSpan} (re-run).
\lean{classicallyFree_of_two_inputs} is \textbf{conditional} on \lean{hgeom}, \lean{hchar}.
\end{status}

\paragraph{The per-type instances.} Each Deligne type has a file choosing Lie data and the Deligne level, and restating the two theorems.
The ``model'' column says how the Lie algebra is realized: as matrices, as the weight-one space of a lattice VOA (Section~\ref{l23:sec:lattice}),
or as a structure-constant table (Section~\ref{l23:sec:sc}). A primed name such as \lean{g2Data'} is the Lie data of a table algebra built by
\lean{standardData}; the unprimed \lean{SCLie.g2Data} is the root table of Theorem~\ref{l23:thm:sctables} (see the warning box in
Section~\ref{l23:warn:models}).

\begin{center}\small
\begin{tabular}{lllllll}
\toprule
type & model & Lie algebra & Lie data & $\lxxiiihv$ & $k$ & $c$\\
\midrule
$A_2$ & matrices & \lean{AffineSl3} & \lean{sl3Data} & 3 & $-3/2$ & $-8$\\
$G_2$ & table & \lean{SCLie.g2Lie} & \lean{G2Deligne.g2Data'} & 4 & $-5/3$ & $-10$\\
$D_4$ & lattice & \lean{CartanSimple.d4Lie} & \lean{LatticeDeligne.d4Data} & 6 & $-2$ & $-14$\\
$F_4$ & table & \lean{SCLie.f4Lie} & \lean{F4Deligne.f4Data'} & 9 & $-5/2$ & $-20$\\
$E_6$ & lattice & \lean{CartanSimple.e6Lie} & \lean{LatticeDeligne.e6Data} & 12 & $-3$ & $-26$\\
$E_6$ & table & \lean{SCLie.e6Lie} & \lean{E6Deligne.e6Data'} & 12 & $-3$ & $-26$\\
$E_7$ & lattice & \lean{CartanSimple.e7Lie} & \lean{LatticeDeligne.e7Data} & 18 & $-4$ & $-38$\\
$E_8$ & lattice & \lean{CartanSimple.e8Lie} & \lean{E8Deligne.e8Data} & 30 & $-6$ & $-62$\\
\bottomrule
\end{tabular}
\end{center}
(The $A_2$ instance lives in the file \leanpath{VOA/\allowbreak AffineGeneral/\allowbreak A2General.lean}.)
The levels are $k=-\lxxiiihv/6-1$ and the central charges $c=\dim\lxxiiig\cdot k/(k+\lxxiiihv)=-2\lxxiiihv-2$ (computed here on paper from
\lean{standardAffineVOA_centralCharge}; the simple quotient has the same conformal vector). The $E_8$ file is representative:

\begin{theorem}[The Deligne case $E_8$ on $L_{-6}(\mathfrak e_8)$]\label{l23:thm:e8deligne}
With $\kappa=\lxxiiiKil/60$ on $\mathfrak e_8=(V_{E_8})_1$, indexed by $\{0,\dots,247\}$ via $\dim\mathfrak e_8=248$, $h=30$, and on
$L_{-6}(\mathfrak e_8)$ the conclusions of Theorem~\ref{l23:thm:deligneinstance} hold.
\end{theorem}
\noindent File \leanpath{VOA/\allowbreak AffineGeneral/\allowbreak E8Deligne.lean}.
\begin{leancode}
def e8Data : LieData ((2 * 30 : ℂ)⁻¹ • killingForm ℂ e8Lie) (Fin HigherRankDeligneCase.E8.currentDimension) :=
  (standardData e8Lie 30 (by norm_num)).reindex (finCongr finrank_e8Lie)

theorem e8Data_h : e8Data.h = 30 := by

abbrev e8SimpleVOA := simpleVOA e8Data (-6) e8_level_ne

theorem e8_currentJetSurjectivity :
    CurrentJetSurjectivity .E8 e8SimpleVOA (Deligne.presentation .E8 e8Data (-6) e8_level_ne) :=
  Deligne.currentJetSurjectivity .E8 e8Data (-6) e8_level_ne

theorem e8_classicallyFree_of_two_inputs
    (Carrier : ℕ → ℕ → Type)
    [∀ p N, AddCommGroup (Carrier p N)]
    [∀ p N, Module ℂ (Carrier p N)]
    [∀ p N, FiniteDimensional ℂ (Carrier p N)]
    (hgeom : GeometricCarrierEmbedding .E8 e8SimpleVOA
      (Deligne.presentation .E8 e8Data (-6) e8_level_ne) Carrier)
    (hchar : OrdinaryCarrierCharacter e8SimpleVOA Carrier) :
    CaseClassicallyFree .E8 e8SimpleVOA (Deligne.presentation .E8 e8Data (-6) e8_level_ne) :=
  Deligne.classicallyFree_of_two_inputs .E8 e8Data (-6) e8_level_ne Carrier hgeom hchar
\end{leancode}

\begin{dictionary}
\lean{finCongr finrank_e8Lie} is the bijection $\mathrm{Fin}(\dim)\simeq\mathrm{Fin}\,248$ given by Theorem~\ref{l23:thm:e8dim}; this is
where the dimension count is used. Conclusion (1) of Theorem~\ref{l23:thm:deligneinstance} is the presentation
\lean{Deligne.presentation .E8 e8Data (-6) e8_level_ne} appearing in both statements; (2) is \lean{e8_currentJetSurjectivity}; (3) is
\lean{e8_classicallyFree_of_two_inputs}, whose proof term is the general theorem specialized, and which is conditional on \lean{hgeom}
and \lean{hchar}. \lean{e8_level_ne} is the proof that $-6+30\ne0$. That $\lxxiiiKil/60$ is the standard normalization is \lean{Normalization.e8_standard}
(Theorem~\ref{l23:thm:normalization}); without it, ``$-6$'' would be a level in an unidentified normalization. Same as the textbook.
\end{dictionary}

\textbf{Proof on paper.} Theorem~\ref{l23:thm:deligneinstance}. \qed

\begin{physics}
$L_{-6}(\mathfrak e_8)$ is the chiral algebra of the rank-one Minahan--Nemeschansky theory with $E_8$ flavour symmetry, $c_{2d}=-62$,
$c_{4d}=31/6$.
\end{physics}

\begin{status}
Proved; axioms standard for \lean{e8_classicallyFree_of_two_inputs} (re-run). Conditional on the two packages, which Chapter~\ref{ch:L24}
discharges for the \emph{explicit} presentations, not for this full-kernel one.
\end{status}

\begin{warning}[title={Trap: Coexisting models and clashing names}]\label{l23:warn:models}
Several Lie algebras exist in more than one Lean model, and short names clash across namespaces. \textbf{Always read the namespace.}
\begin{center}\footnotesize
\begin{tabular}{p{0.08\textwidth}p{0.25\textwidth}p{0.27\textwidth}p{0.27\textwidth}}
\toprule
algebra & model & Lie data & simple affine VOA\\
\midrule
$\mathfrak e_6$ (i) & lattice: \lean{CartanSimple.e6Lie} $=(V_{E_6})_1$ & \lean{LatticeDeligne.e6Data} ($\lxxiiiKil/24$) & \lean{LatticeDeligne.e6SimpleVOA}\\
$\mathfrak e_6$ (ii) & table: \lean{SCLie.e6Lie} & \lean{E6Deligne.e6Data'} ($\lxxiiiKil/24$, Mathlib basis) & \lean{E6Deligne.e6SimpleVOA}\\
$\mathfrak e_6$ (ii) & table: \lean{SCLie.e6Lie} & \lean{SCLie.e6LieData} (integer form, table basis) & \lean{W6}\\
$\mathfrak e_6$ (iii) & certificate basis \lean{E6Lat.xs} inside (i) & none (brackets only) & ---\\
$\mathfrak g_2$ & table: \lean{SCLie.g2Lie} & \lean{G2Deligne.g2Data'}; \lean{G2Lie.g2CertData} & \lean{g2W} (for the latter)\\
$\mathfrak f_4$ & table: \lean{SCLie.f4Lie} & \lean{F4Deligne.f4Data'}; \lean{F4Lie.f4CertData} & \lean{f4W} (for the latter)\\
$\mathfrak{sl}_3$ & matrices; lattice \lean{CartanSimple.a2Lie} & \lean{sl3Data} (matrices) & $\lxxiiiL$ for \lean{sl3Data} (\leanpath{A2General.lean})\\
\bottomrule
\end{tabular}
\end{center}
Here \lean{W6} is \lean{simpleVOA e6LieData (-3) _}. No Lean theorem states that the $\mathfrak e_6$ models (i) and (ii) are isomorphic;
the docstring of \leanpath{E6LatTable.lean} asserts that its table equals \lean{SCLie.e6T}, which is not checked in Lean. By the project's
decision, the $E_6$ sector data of Chapter~\ref{ch:L24} (\lean{DeligneCarrier.E6.E6sector}, the integer tables of the divisor comparison) are
built on \lean{SCLie.e6LieData}, and the
final $E_6$ theorems of Chapter~\ref{ch:L27} are about \lean{W6}.

\textbf{Name clashes.} \lean{e6Lie_isSimple} names two different instances (\lean{CartanSimple.}, \lean{SCLie.}); \lean{e6SimpleVOA} names
two different VOAs (\lean{LatticeDeligne.}, \lean{E6Deligne.}); \lean{e6Data} is Lie data in \lean{LatticeDeligne} but a \emph{root table}
(\lean{RootData 78}) in \lean{SCLie}, and likewise \lean{SCLie.g2Data}, \lean{SCLie.f4Data} are root tables while the Lie data are the primed
\lean{g2Data'}, \lean{f4Data'} and \lean{G2Lie.g2CertData}, \lean{F4Lie.f4CertData}.
\end{warning}

%======================================================================
\section{The certificate-basis models}\label{l23:sec:cert}
%======================================================================

Several steps of Chapter~\ref{ch:L24} rest on large integer certificates: linear-algebra computations (ranks, divisibilities, explicit
relations among current modes in each weight ``sector'', i.e.\ each pair of a weight and a Li degree) that were carried out by Python in a
fixed ``certificate basis'' (a Chevalley basis with Frenkel--Kac signs, files \texttt{casimir\_cert/cert\_X.json}) and are re-checked in
Lean's kernel. A certificate is a statement about structure constants \emph{in its basis}: if Lean's Lie data used a different basis, every
certificate would have to be transformed and re-verified, which for $\mathfrak e_8$ means tables of size $248^3$. So Lean needs Lie data
whose structure constants and form are \emph{literally} the certificate tables. For the simply-laced types $D_4,E_7,E_8$ (and $E_6$, brackets only) this is achieved by realizing
the certificate basis inside the lattice Lie algebra; for $G_2,F_4$ by \lean{lieDataOf} on the table algebra (above). We describe the
mechanism once, with $E_8$ quoted; $D_4$ and $E_7$ are identical in structure.

\paragraph{The basis.} With $\alpha_i$ the simple roots and, for each root index $a$, its weight $\mathrm{wt}(a)$ and an explicit sign $u_a=\pm1$,
\begin{equation*}
x_i=\alpha_i(-1)\lxxiiivac\ (i<\mathrm{rank}),\qquad x_a=u_a\,f(\mathrm{wt}\,a)\,e^{\mathrm{wt}\,a},
\end{equation*}
where $f$ is a coboundary sign that converts the library's FLM cocycle $\varepsilon$ (defined from a \emph{chosen}, noncomputable
$\Z$-basis, so its values cannot be evaluated) into the explicit cocycle $\varepsilon'(\alpha,\beta)=(-1)^{\sum_{i<j}\alpha_i\beta_jG_{ij}+\sum_i\alpha_i\beta_i}$.

\begin{theorem}[Certificate basis of $\mathfrak e_8$]\label{l23:thm:cert}
(1) $\varepsilon(\alpha,\beta)f(\alpha)f(\beta)=\varepsilon'(\alpha,\beta)f(\alpha+\beta)$. (2) For every pair $(a,b)$ the kernel check
\lean{pairOK a b} holds, hence $[x_a,x_b]=\sum_cs_{abc}x_c$ with $s$ the certificate table. (3) The $x_a$ form a basis, and
$\lxxiiiKil(x_a,x_c)=60\,G_{ac}$ with $G$ the certificate form. (4) Hence there is Lie data for $\lxxiiiKil/60$ with basis $x$, $h=30$, whose
structure constants and Gram matrix are the certificate tables.
\end{theorem}
\noindent Files \leanpath{VOA/\allowbreak DeligneAudit/\allowbreak E8LieCocycle.lean},
\leanpath{E8LieBasis.lean}, \leanpath{E8LieTable.lean}, \leanpath{E8LieModel.lean}, \leanpath{E8LieModel00.lean}, \leanpath{E8LieKilling.lean}, \leanpath{E8LieForm.lean}, \leanpath{E8LieData.lean} (all in
\leanpath{VOA/\allowbreak DeligneAudit/}); \lean{casimir_transfer} in \leanpath{VOA/\allowbreak DeligneAudit/\allowbreak D4LieCasimirTransfer.lean} (split off from
\leanpath{D4LieData.lean}, which imports it).
\begin{leancode}
theorem cocycle_coboundary (α β : Lat) :
    latticeCocycle e8 α β * f α * f β = (-1 : ℂ) ^ ep α β * f (α + β) := by

def xs (a : Fin 248) : e8Lie :=
  if a.val < 8 then hst (hvec a) else ((u a : ℂ) * f (wt a)) • evR a

def pairOK (a b : Fin 248) : Bool :=
  decide (sL a b = modelL a b) &&

theorem bracket_of_pairOK (a b : Fin 248) (hab : pairOK a b = true) :
    ⁅xs a, xs b⁆ = spL (sL a b) := by

theorem pairOK_all (a b : Fin 248) : pairOK a b = true := by

theorem row_ok0 : ∀ b, pairOK 0 b = true := by decide +kernel

theorem killing_xs (a c : Fin 248) : K (xs a) (xs c) = 60 * (Gz a c : ℂ) := by

theorem casimir_transfer (D : LieData κ ι) {ι' : Type*} [Fintype ι'] (b' d' : ι' → 𝔤)
    (hexp : ∀ y, y = ∑ j, κ y (d' j) • b' j) (x : 𝔤) :
    ∑ j, ⁅b' j, ⁅d' j, x⁆⁆ = ∑ i, ⁅D.basis i, ⁅D.dual i, x⁆⁆ := by

theorem span_xs : ⊤ ≤ Submodule.span ℂ (Set.range xs) := by

irreducible_def bB : Module.Basis (Fin 248) ℂ e8Lie :=
  basisOfTopLeSpanOfCardEqFinrank xs span_xs (by rw [Fintype.card_fin, finrank_e8Lie])

def Gi2L (a : Fin 248) : List (Fin 248 × ℤ) := Gi2List.getD a.val []

def pairRowOK (a : Fin 248) : Bool :=
  (List.finRange 248).all fun b => decide (((Gi2L b).map fun t => t.2 * Gz a t.1).sum = if a = b then 2 else 0)

def dualx (b : Fin 248) : e8Lie := ((Gi2L b).map fun t => ((t.2 : ℂ) / 2) • xs t.1).sum

irreducible_def e8CertData : LieData κ8 (Fin 248) where
  invariant := e8Data.invariant
  basis := bB
  dual := dualx
  pair i j := by rw [bB_apply]; exact pair_xs i j
  h := 30

theorem structureConstant_eq (a b c : Fin 248) :
    e8CertData.structureConstant a b c = (sZ a b c : ℂ) := by

theorem formTable_eq (a b : Fin 248) : e8CertData.formTable a b = (Gz a b : ℂ) := by
\end{leancode}

\begin{dictionary}
\lean{Lat} is $\Z^8$; \lean{ep α β} is the exponent of $\varepsilon'$. \lean{sL a b} is the certificate's sparse row, \lean{modelL a b} the bracket
of the rescaled lattice basis computed with $\varepsilon'$; \lean{pairOK} checks that they agree and some side conditions on weights (the omitted
lines: for two roots, $\langle\mathrm{wt}\,a,\mathrm{wt}\,b\rangle\ge-2$, $=-2$ only for opposite roots, $=-1$ only when the target index has weight
$\mathrm{wt}\,a+\mathrm{wt}\,b$), which let \lean{bracket_of_pairOK} evaluate $[x_a,x_b]$ in the VOA via Theorem~\ref{l23:thm:latticesimple}(1) and
the formula $e^\alpha_{(0)}e^\beta=\varepsilon(\alpha,\beta)e^{\alpha+\beta}$ for $\langle\alpha,\beta\rangle=-1$. \lean{spL} turns a sparse list into a vector,
\lean{sZ} is the dense integer table. \lean{K} is the Killing form of \lean{e8Lie}. Claim (3) is \lean{span_xs} (the $248$
vectors span; \lean{⊤ ≤ span} means the span is everything) together with \lean{bB}, Mathlib's ``spanning family of size $\dim$ is a
basis'', using \lean{finrank_e8Lie}~$=248$, and \lean{killing_xs}. Claim (4) is \lean{e8CertData} with \lean{structureConstant_eq},
\lean{formTable_eq}. \lean{irreducible_def} makes a definition opaque except through its defining equation (used for elaboration speed).
\lean{Gz} is the integer certificate form $G$ (in \leanpath{E8LieForm.lean}); \lean{Gi2L b} is the sparse $b$-th row of the integer
matrix $\mathrm{Gi2}$, and the kernel check \lean{pairRowOK} (all rows, in \leanpath{E8LieFormChecks.lean}) says $\sum_c\mathrm{Gi2}_{bc}G_{ac}=2\delta_{ab}$, i.e.\
$\mathrm{Gi2}=2G^{-1}$. So \lean{dualx b}$=\sum_c\tfrac12\mathrm{Gi2}_{bc}\,x_c=\sum_c(G^{-1})_{bc}\,x_c$, the $\kappa$-dual of $x_b$ for the form
with Gram matrix $G$. The omitted last field of \lean{e8CertData} is the Casimir proof. \lean{e8CertData : LieData κ8 (Fin 248)} has form
\lean{κ8 = (2 * 30 : ℂ)⁻¹ • killingForm ℂ e8Lie} (the same as \lean{e8Data}), basis \lean{bB} (with \lean{bB a = xs a}, by \lean{bB_apply}), dual vectors \lean{dualx}, and
$h=30$, its Casimir field supplied by \lean{casimir_transfer} from \lean{E8Deligne.e8Data}: the Casimir does not depend on the choice of dual
pair. Same as the textbook.
\end{dictionary}

\textbf{Proof on paper.} (1) Any two bimultiplicative cocycles with the same commutator $(-1)^{\langle\alpha,\beta\rangle}$ differ by a coboundary
$f(\alpha)f(\beta)/f(\alpha+\beta)$ with $f=(-1)^q$, $q$ a quadratic refinement of the difference of exponents (computed bilinearly from the
values on basis vectors). (2) Finite evaluation of each bracket of basis vectors. (3) Spanning: the $x_a$ cover the Cartan part and every root
(the root count of Theorem~\ref{l23:thm:rootcounts} shows the certificate weights are all the roots), $248$ vectors in dimension $248$; the Killing form is obtained without
computing any trace, by invariance from the coroot values $\lxxiiiKil(h_\beta,h_\beta)=120$ (Theorem~\ref{l23:thm:normalization}):
$\lxxiiiKil(E_i,F_i)=60$, $\lxxiiiKil(H_j,H_i)=60A_{ji}$, pairs of nonzero total weight are orthogonal, and opposite root pairs are determined by
$\lambda\lxxiiiKil(x_a,x_{-a})=\lxxiiiKil(H,[x_a,x_{-a}])$. (4) assembles these. \qed

\begin{physics}
The certificate basis is a choice of currents with integer OPE coefficients; this theorem identifies the integer OPE table used in the
computer algebra with the actual current algebra of the $E_8$ lattice CFT, sign by sign.
\end{physics}

\begin{status}
Proved; axioms standard for \lean{E8Lie.e8CertData}, \lean{E8Lie.structureConstant_eq}, \lean{E8Lie.pairOK_all} (re-run), and
\lean{E6Lat.bracket_xs}. \lean{span_xs} and \lean{bB} were not re-run in the revision (that needs the heavy $E_8$ imports); \lean{e8CertData},
which was re-run, uses \lean{bB} as its basis, so its clean report covers \lean{bB} and \lean{span_xs}, and \leanpath{E8LieData.lean} has
its own \texttt{\#print axioms bB}.
\end{status}

\paragraph{What is checked, collectively, and at what size.}
\begin{itemize}
\item \textbf{$D_4$} (\leanpath{D4Lie\{Basis,Cocycle,Data,Table\}}): the table comparison \lean{sZ_eq_model} ($28^3$ entries) is one
  \lean{decide +kernel} in \leanpath{D4LieTable.lean}, together with the weight facts and \lean{killZ_eq} ($\lxxiiiKil=12G$); result
  \lean{D4Lie.d4CertData : LieData ((2 * 6)⁻¹ • killingForm ℂ d4Lie) (Fin 28)}, $h=6$.
\item \textbf{$E_7$} (\leanpath{E7Lie\{Basis,CasimirChecks,Cocycle,Data,Facts,Form,FormChecks,Killing,Table\}} and 20
  \leanpath{E7LieModel*}): all $133^2=17689$ pairs checked by \lean{pairOK}; the Killing
  form by invariance from \lean{e7_killing} $=72$; \lean{E7LieCasimirChecks} checks that the certificate Casimir is the inverse-form quadric;
  result \lean{e7CertData}, form $\lxxiiiKil/36$, $h=18$.
\item \textbf{$E_8$} (\leanpath{E8Lie\{Basis,Cocycle,Data,DataChecks,Facts,Form,FormChecks,Killing,Table\}} and 32 \leanpath{E8LieModel*}):
  all $248^2=61504$ pairs checked by \lean{pairOK} (assembled in \leanpath{E8LieModel.lean}); result \lean{e8CertData}, $\lxxiiiKil/60$, $h=30$.
\item \textbf{$E_6$} (18 \leanpath{E6Lat*} files, 13 of them \leanpath{E6LatModel*}): $78^2=6084$ pairs; result only
  \lean{E6Lat.bracket_xs} (the certificate structure constants realized in the lattice algebra), used by the lattice realizations of
  Chapter~\ref{ch:L26}; no Lie data (the $E_6$ Lie data in the certificate basis is \lean{SCLie.e6LieData}).
\item \textbf{$G_2$, $F_4$} (\leanpath{G2LieData.lean}; \leanpath{F4LieData}, \leanpath{F4LieTab}, ten \leanpath{F4LieChecks*}): the
  \lean{lieDataOf} checks (symmetry, skewness in four chunks, pairing in two, Casimir in four for $\mathfrak f_4$); results
  \lean{g2CertData} ($\lxxiiiKil/8$, $h=4$) and \lean{f4CertData} ($\lxxiiiKil/18$, $h=9$), with \lean{g2W}, \lean{f4W} the simple quotients at
  $k=-5/3,-5/2$.
\end{itemize}
In every case the data were generated by untrusted scripts and every fact
used is re-established in the kernel; no data file contains a \lean{sorry} or \lean{native_decide}. How the checks are split across files,
and what they cost, is recorded in Section~\ref{l23:sec:verification}.

%======================================================================
\section{Verification record}\label{l23:sec:verification}
%======================================================================

This section collects, in one place, what the machine confirmed and where the record comes from. The rest of the chapter refers here.

\textbf{Provenance.} The chapter was written against an earlier state of the library (the Lean sources under \leanpath{VOA/} did not
change between the snapshot against which it was first written and the commit at which it was revised). All quotations are checked
against the released sources by \texttt{scripts/textbook/sync.py}. Runs made before 28 October 2026 used the former names
\lean{simpleAffineVOA} and \lean{simpleAffineVOA_centralCharge}; they are printed below under the current names
\lean{universalAffineVOA} and \lean{universalAffineVOA_centralCharge}.

\textbf{Builds (from the project record, not re-run for this chapter).} The project's workstream log
\leanpath{Notes/\allowbreak deligne/\allowbreak campaign/\allowbreak WS-AFF.md} records \texttt{campaign\_build.sh VOA.AffineGeneral.All}
$\to$ ``Build completed successfully (4329 jobs)'', and a build with \texttt{VOA.DeligneAudit.E6CarrierData} of 9212 jobs. The project's
internal build record of 27 October 2026 reports \texttt{campaign\_build.sh VOA.Deligne} green at an earlier commit, 11099 jobs, no
\lean{sorryAx}, and the last aggregate \texttt{lake build VOA} (26 October 2026) green with 13868 jobs. This chapter did \emph{not} run
\texttt{lake build}.

\textbf{What ``project record'' means in a status box.} Where a status box says ``axioms standard (project record)'', the
\texttt{\#print axioms} check is the one recorded in the workstream log for that file, not re-run for this chapter; where it says
``re-run'', the declaration is in one of the lists below.

\textbf{Axioms (re-run for this chapter).} Two scratch files importing the relevant modules were checked with
\texttt{lake env lean -DmaxSynthPendingDepth=3} against the existing build products; every declaration below reported
\lean{[propext, Classical.choice, Quot.sound]}:
\begin{plaincode}
affineVOA  universalAffineVOA  universalAffineVOA_centralCharge  standardAffineVOA_centralCharge
killing_h  casimir_scalar  simpleVOA_isSimple  simple_grading_one_eq_range
mem_maxIdeal_of_SingP  SlSimple.instIsSimple  RootSpace.isSimple_of_roots
LatticeLie.isSimple_of_simpleRoots  LatticeLie.finrank_V1  LatticeLie.killing_hV_hV
CartanSimple.e8Lie_isSimple  RootCounts.finrank_e7Lie  RootCounts.finrank_e8Lie'
E8Dim.finrank_e8Lie  RootEnum.mem_enum  Normalization.e8_standard
SCLie.isSimple_of_tables  SCLie.{g2,f4,e6}Lie_isSimple  SCLie.g2_standard
SCLie.lieDataOf  SCLie.e6LieData  Deligne.currentJetSurjectivity
Deligne.classicallyFree_of_two_inputs  E8Deligne.e8_classicallyFree_of_two_inputs
DeligneAudit.E8Lie.{e8CertData, structureConstant_eq, pairOK_all}
DeligneAudit.E6Lat.bracket_xs  DeligneAudit.G2Lie.g2CertData
DeligneLatticeLie.e8VOA_centralCharge
\end{plaincode}
In the revision after the fidelity report, a third scratch file (importing the \texttt{AffineGeneral} modules \texttt{RootCounts}, \texttt{G2Lie}, \texttt{F4Lie},
\texttt{E6Lie}, \texttt{CartanLatticeSimple}, \texttt{LatticeV1Simple}, \texttt{DeligneInstance} and the capstone file; no $E_7$/$E_8$
data modules) gave the same three axioms for:
\begin{plaincode}
CartanSimple.cartan_isSimple  CartanSimple.cartan_finiteDimensional  LatticeLie.eV_ne_zero
Deligne.universal_genWordsSpan  Deligne.genWordsSpan
CurrentC2Presentation.exact_kernel_of_caseClassicallyFree
RootCounts.card_roots_{d4,e6,e7,e8'}  SCLie.g2_coroot
SCLie.{f4,e6}_killing  SCLie.{f4,e6}_standard  SCLie.{f4,e6}Lie_isSimple
\end{plaincode}
For the named simple quotient (Section~\ref{l23:sec:simplename}, added 2026-10-28; not re-run for this edition), a fourth scratch file --- a copy of
\leanpath{VOA/\allowbreak AffineGeneral/\allowbreak SimpleAffineQuotient.lean} followed by \texttt{\#print axioms} commands --- was
elaborated with \texttt{lake env lean -DmaxSynthPendingDepth=3} against the existing build products of its imports (no errors, no
warnings) and gave the same three axioms for:
\begin{plaincode}
simpleAffineQuotientVOA  simpleAffineQuotientVOA_isSimple  simpleAffineQuotientVOA_centralCharge
simpleAffineQuotientProjection  simpleAffineQuotientProjection_smul
simpleAffineQuotientProjection_surjective  simpleAffineQuotientProjection_grading
\end{plaincode}
The targeted and umbrella builds including this file are recorded green with standard axioms in \texttt{Notes/RELEASE\_PLAN.md}
(``Caveats removed in Lean'', 2026-10-28); the aggregate build \texttt{bash scripts/campaign\_build.sh VOA} (the full \texttt{lake build VOA} under the project's memory-capped wrapper) was green in two pre-release builds on 2026-10-28 (18{,}901 jobs, and 18{,}907 jobs after merging the then-current
\texttt{origin/main}); the release build is recorded in Chapter~\ref{ch:15}, Section~\ref{ch15:sec:record}. No build was re-run
for this chapter.
The library has no project axiom (\lean{VOA.pbw} is a theorem, Chapter~\ref{ch:L02}).

\paragraph{Release audit 2026-10-27: RELEASE\_AUDIT\_L17\_L18\_L23.md.} The independent release audit
\texttt{Notes/\allowbreak Textbook/\allowbreak referee/\allowbreak RELEASE\_AUDIT\_L17\_L18\_L23.md} re-elaborated the headline
declarations of this chapter from the existing build products (\texttt{lake env lean -DmaxSynthPendingDepth=3}; \texttt{lake build}
not run) and re-ran \texttt{\#print axioms}: \lean{universalAffineVOA}, \lean{universalAffineVOA_centralCharge},
\lean{standardAffineVOA_centralCharge}, \lean{killing_h}, \lean{simpleVOA_isSimple}, \lean{simple_grading_one_eq_range},
\lean{cartan_isSimple}, the lattice and table \lean{*Lie_isSimple}, \lean{SCLie.finrank_{g2,f4,e6}Lie}, \lean{SCLie.*_standard},
\lean{g2Data'}, \lean{g2_currentJetSurjectivity} and \lean{g2_classicallyFree_of_two_inputs} all report only the three standard
axioms. The $E_7$/$E_8$ dimension and root-count theorems were not re-run by the audit (their sources were read and are
unconditional). Verdict: \emph{announceable with caveats}, the caveats being mandatory, and they are those stated in this chapter:
the declaration then called \lean{simpleAffineVOA} (now \lean{universalAffineVOA}) is the universal $\lxxiiiV$ and $\lxxiiiL$ is
\lean{simpleVOA (simpleData …)}, at audit time without a named wrapper (since 2026-10-28: \lean{simpleAffineQuotientVOA},
Section~\ref{l23:sec:simplename}) (M1, warning box in Section~\ref{l23:warn:simplename}); $h$ is the Casimir number of the chosen form, equal to $\lxxiiihv$ only per type
via the \lean{*_standard} theorems (M2); Cartan types are not identified in Lean (M3, Remark~\ref{l23:rem:type}); \lean{weightOneEquiv}
is only a linear isomorphism (m1); longness of the chosen roots of the table algebras is not checked (m2); classical freeness remains
conditional on \lean{hgeom}, \lean{hchar} (m3). The audit also notes (M7) that the aggregate build is not certified for the release
commit: the last aggregate green cited above predates later commits, and a transitive-import walk found umbrella modules without build
products. (Since then the aggregate build was green on 2026-10-28, as recorded above.) Nothing in this chapter's status boxes is
upgraded by the audit.

\subsection{File map}
The chapter covers the 268 files assigned to it in \texttt{Notes/Textbook/assignment.txt}:
\begin{itemize}
\item the 161 files of \leanpath{VOA/\allowbreak AffineGeneral/}: 43 theory files (the general affine vertex algebra, its simple quotient,
  simplicity criteria, root counts, normalizations, the per-type instances) and 118 generated files for the structure-constant algebras
  $\mathfrak g_2$, $\mathfrak f_4$, $\mathfrak e_6$;
\item \leanpath{VOA/\allowbreak Foundations/\allowbreak CartanLattices.lean} and \leanpath{VOA/\allowbreak Examples/\allowbreak DeligneLatticeLie.lean}
  (the simply-laced root lattices and their lattice vertex algebras);
\item the certificate-basis models in \leanpath{VOA/\allowbreak DeligneAudit/}: \leanpath{D4Lie\{Basis,Cocycle,Data,Table\}.lean}, the
  \leanpath{E7Lie*} and \leanpath{E8Lie*} data files (basis, cocycle, data, table, form, Killing form, and 20 resp.\ 32 \leanpath{Model*}
  check files), the 18 \leanpath{E6Lat*.lean} files, \leanpath{G2LieData.lean}, \leanpath{F4LieData.lean}, \leanpath{F4LieTab.lean} and the ten
  \leanpath{F4LieChecks*.lean} files.
\end{itemize}
\emph{The 118 generated files of \leanpath{AffineGeneral/}.} For each $X\in\{\mathfrak g_2,\mathfrak f_4,\mathfrak e_6\}$ the generator
\texttt{aff\_sc\_emit2.py} writes: \leanpath{XTable.lean} (the flat rows only; $3$\,KB, $24$\,KB, $46$\,KB of data), four \leanpath{XBound*} and
four \leanpath{XAnti*} files (index bounds and antisymmetry, one row chunk each, combined by \lean{boundOK_of_chunks}/\lean{antiOK_of_chunks}),
the Jacobi files \leanpath{XJac*} (one declaration \lean{jacPiece T N i jlo jhi} per row $i$ and column block, covering all sorted triples
$i<j<l$; $12$ declarations for $\mathfrak g_2$ in one file, $125$ for $\mathfrak f_4$ in ten files, $266$ for $\mathfrak e_6$ in 29 files), and the
root-table files \leanpath{XRootsA}, \leanpath{XHbe*}, \leanpath{XHhb*}, \leanpath{XSep*}, \leanpath{XSymm*} (the fields of \lean{RootChecks}). For
$\mathfrak e_6$ there are in addition the Lie-data check files \leanpath{E6LieDataTab}, \leanpath{E6LieDataPair\{A,B,C\}},
\leanpath{E6LieDataCas\{A..F\}}, \leanpath{E6LieDataSkew\{A,B\}}. Apart from the three \leanpath{XTable.lean} files, which hold data only, every
file is a list of statements ``this Boolean function of the table evaluates to \lean{true}'', proved by \lean{decide +kernel}. In numbers:
$23$, $32$ and $51$ files for $\mathfrak g_2$, $\mathfrak f_4$, $\mathfrak e_6$ (106 in all), plus the 12 \leanpath{E6LieData*} files.
\emph{The certificate-basis check files.} $E_7$: $133^2$ pairs, one row per declaration, combined by \lean{fin_cases}. $E_8$: $248^2$ pairs in
$31$ files of $8$ rows (one row, i.e.\ $248$ pairs, per \lean{decide +kernel}), plus the assembly \leanpath{E8LieModel.lean}.

\subsection{Machine cost, and the dead ends that shaped the files}
Measured costs from the project record: every $\mathfrak e_6$ table-check file $\le25$\,s and $\le5.4$\,GB; each $E_8$ certificate row about
3\,s and under 3\,GB; the pruned $E_7$ root enumeration (Definition~\ref{l23:def:enum}) about 2\,s in the kernel.

\begin{warning}[title={Dead ends: kernel cost}]\label{l23:warn:kernel}
These routes were tried and abandoned; they explain why the checks are split as they are.
\begin{itemize}
\item \emph{Plain box enumeration of roots.} Enumerating all integer vectors in a coordinate box and checking norms under
  \lean{decide +kernel}: for $E_6$ (7875 vectors) about 11\,s; for $E_7$ (box $(2,2,3,4,3,2,1)$, 165375 vectors) it did not finish in
  10 minutes. The pruned $LDL^{\mathsf T}$ enumeration visits only branches compatible with the bound (about 2\,s for $E_7$).
\item \emph{Table encodings.} A nested-tuple $\mathfrak e_6$ table took 112\,s to elaborate (flat \lean{List ℕ} rows: about 5\,s); \lean{![…]}
  vector tables are slow in the kernel (\lean{List} with \lean{getD} is used instead).
\item \emph{Monolithic checks.} The first $\mathfrak f_4$ file checked the Leibniz identity on all $52^3$ triples in one declaration (2\,min
  15\,s; superseded by sorted cyclic Jacobi in pieces). A single \lean{decide} comparing a $133\times133$ sector table with the lattice-model
  table used 41\,GB and timed out after 445\,s; the $E_8$ version was killed above 20\,GB; even $D_4$ ($28^2$) needs 6.1\,GB as one declaration.
\end{itemize}
The working pattern: light-import check files with one row or block per declaration, and a heavy file that only assembles the
propositional conclusions.
\end{warning}

%======================================================================
\section{Lean defects found}\label{l23:sec:defects}
%======================================================================

This section is an index; each item is explained in the box cited. Severities: S1, a statement or definition is wrong or vacuous; S2, a definition is weaker or stronger than its name or the literature, or a hypothesis is redundant; S3, an advertised result is only conditional (a gap); S4, documentation (docstrings, notes) is wrong or stale.
\begin{enumerate}
\item (S4) $\mathfrak{sl}_3$ prose (``$8k/(k+3)$'', ``$1/(2(k+3))$'', ``eight Chevalley directions'', ``universal sl3 vacuum'') in the
  general files \leanpath{Virasoro.lean}, \leanpath{SugawaraField.lean}, \leanpath{WeightOne.lean}, \leanpath{GradedFields.lean}: warning box in
  Section~\ref{l23:warn:sl3docs}.
\item (S3) False provenance of vacuum nonvanishing (``free-field evaluation, without PBW'') in \lean{AffineGeneral.vacuumA_ne_zero},
  \leanpath{Grading.lean}, \leanpath{Currents.lean}: warning box in Section~\ref{l23:warn:freefield}.
\item (S4) The docstring of \lean{AffineGeneral.currentState_injective_of_ne_zero} points to the $\mathfrak{sl}_3$ all-level injectivity as if
  it applied to general $\lxxiiig$: dictionary of Theorem~\ref{l23:thm:weightone} and warning box in Section~\ref{l23:warn:simpleweightone}.
\item (S4) \lean{SlSimple.slAffineVOA}: redundant hypothesis \lean{hm}, and ``standard normalization'' claimed without proof: warning box in
  Section~\ref{l23:warn:slnnorm}.
\item (S3) Name clashes (\lean{e6Lie_isSimple}, \lean{e6SimpleVOA}, \lean{e6Data}, \lean{g2Data}, \lean{f4Data}) and the unchecked docstring
  claim of \leanpath{DeligneAudit/E6LatTable.lean} that its table equals \lean{SCLie.e6T}: warning box in Section~\ref{l23:warn:models}.
\item (S4) \lean{SCLie.hL}/\lean{hbV} called ``coroot'' although they are minus the standard coroot: warning box in
  Section~\ref{l23:warn:hbvsign}.
\item (S3, resolved) The former name \lean{AffineGeneral.simpleAffineVOA} invited the reading $L_k(\lxxiiig)$ although the declaration
  is the universal $\lxxiiiV$, and the simple quotient of a general simple $\lxxiiig$ had no named declaration (release audit 2026-10-27,
  M1): warning box in Section~\ref{l23:warn:simplename}. Resolved 2026-10-28: the declaration is now \lean{universalAffineVOA} (the old
  name is a deprecated alias), and the quotient is named \lean{simpleAffineQuotientVOA} (Section~\ref{l23:sec:simplename}).
\item (S4) Stale docstring of \leanpath{Examples/DeligneLatticeLie.lean} (simplicity and dimensions called ``the next steps''): warning box in
  Section~\ref{l23:warn:stalelattice}.
\end{enumerate}

%======================================================================
\section{Referee's checklist}\label{l23:sec:checklist}
%======================================================================

\begingroup\small
\begin{longtable}{p{0.33\textwidth}p{0.23\textwidth}p{0.21\textwidth}p{0.11\textwidth}}
\toprule
Lean declaration & Intended meaning & Hypotheses / status & Fidelity check\\
\midrule\endhead
\lean{LieData} & basis, dual basis, Casimir $=2h$ & definition & Def.~\ref{l23:def:liedata}\\
\lean{struct_casimir_contract} etc. & finite identities & unconditional & Prop.~\ref{l23:prop:finite}\\
\lean{current_comm}, \lean{vacuumA_ne_zero} & affine algebra; $\lxxiiivac\ne0$ & unconditional (relative PBW) & Def.~\ref{l23:def:vacuum}\\
\lean{affineVOA}, \lean{hYA_virasoro} & $V^k(\lxxiiig)$, $c=\dim\lxxiiig\,k/(k+h)$ & Lie data, $k+h\ne0$ & Thm.~\ref{l23:thm:affineVOA}\\
\lean{casimir_scalar} & Casimir scalar on simple $\lxxiiig$ & simple, f.d. & Thm.~\ref{l23:thm:casimir}\\
\lean{universalAffineVOA}, \lean{killing_h}, \lean{standardAffineVOA} & universal $V^k(\lxxiiig)$ (\emph{not} $L_k$) for every simple $\lxxiiig$; $2h$ the Casimir of $\kappa$ & simple; $h=\lxxiiihv$ only per type (\lean{*_standard}) & Thm.~\ref{l23:thm:simpleaffine}, Sec.~\ref{l23:warn:simplename}\\
\lean{simpleVOA}, \lean{simpleVOA_isSimple}, \lean{simple_grading_one_eq_range} & $L_k(\lxxiiig)$ (for general $\lxxiiig$: \lean{simpleVOA (simpleData …)}, named \lean{simpleAffineQuotientVOA}) simple; weight one spanned by currents & $k+h\ne0$; injectivity \textbf{not stated} & Thm.~\ref{l23:thm:simple}\\
\lean{simpleAffineQuotientVOA}, \lean{_isSimple}, \lean{_centralCharge}, \lean{simpleAffineQuotientProjection}, \lean{_smul}, \lean{_surjective}, \lean{_grading} & named $L_k(\lxxiiig)$ for simple $\lxxiiig$: simple, $c=\dim\lxxiiig\,k/(k+h)$; $\pi$ a surjective, grading-preserving homomorphism & $k+h\ne0$; unconditional & Def.~\ref{l23:def:simpleaffinequotient}, Thm.~\ref{l23:thm:simpleaffinequotient}\\
\lean{weightOneEquiv} & $\lxxiiig\cong V^k_1$ linearly (no bracket compatibility) & $k\ne0$ & Thm.~\ref{l23:thm:weightone}\\
\lean{mem_maxIdeal_of_SingP}, \lean{quadState_mem_maxIdeal} & singular vectors vanish in $L_k$ & weight $>0$ & Thm.~\ref{l23:thm:singular}\\
\lean{clsG_quad_eq_zero}, \lean{clsG_adPoly_eq_zero} & kernel facts in Li degree 0 & unconditional & Prop.~\ref{l23:prop:clsG}\\
\lean{SlSimple.instIsSimple} & $\mathfrak{sl}_n$ simple & $n\ge2$ & Thm.~\ref{l23:thm:sln}\\
\lean{RootSpace.isSimple_of_roots} & root-space criterion & listed & Thm.~\ref{l23:thm:rootspace}\\
\lean{LatticeLie.}\allowbreak\lean{isSimple_of_simpleRoots}, \lean{finrank_V1} & $(V_L)_1$ simple; $\dim=\mathrm{rank}+\#$roots & connected generating roots & Thm.~\ref{l23:thm:latticesimple}\\
\lean{CartanSimple.*Lie_isSimple} & $(V_L)_1$ simple, $L=A_2,D_4,E_6,E_7,E_8$ & unconditional; Cartan type not identified in Lean & Thm.~\ref{l23:thm:cartansimple}, Rem.~\ref{l23:rem:type}\\
\lean{E8Dim.finrank_e8Lie}, \lean{RootCounts.finrank_*} & dims $248$; $28,78,133,248$ & unconditional & Thms.~\ref{l23:thm:e8dim}, \ref{l23:thm:rootcounts}\\
\lean{Normalization.*_standard} & $\lxxiiiKil/(2\lxxiiihv)$ standard & unconditional & Thm.~\ref{l23:thm:normalization}\\
\lean{SCLie.isSimple_of_tables}, \lean{g2/f4/e6Lie_isSimple}, \lean{*_standard} & table algebras of dims $14,52,78$ simple, normalized & kernel checks; type and root longness on paper & Thms.~\ref{l23:thm:sctables}, \ref{l23:thm:g2f4e6}\\
\lean{SCLie.lieDataOf}, \lean{e6LieData}, \lean{g2CertData}, \lean{f4CertData} & Lie data in table basis & kernel checks & Def.~\ref{l23:def:lieDataOf}\\
\lean{Deligne.presentation}, \lean{currentJetSurjectivity}, \lean{classicallyFree_of_two_inputs} & capstone inputs for every case & presentation = \textbf{full kernel}; last is conditional on \lean{hgeom}, \lean{hchar} & Thm.~\ref{l23:thm:deligneinstance}, Sec.~\ref{l23:warn:fullkernel}\\
\lean{E8Deligne.*} and per-type files & $L_{-\lxxiiihv/6-1}(\lxxiiig)$ & as above & Thm.~\ref{l23:thm:e8deligne}\\
\lean{E8Lie.e8CertData}, \lean{D4Lie.d4CertData}, \lean{e7CertData}, \lean{E6Lat.bracket_xs} & certificate tables are the lattice brackets & kernel checks & Thm.~\ref{l23:thm:cert}\\
\bottomrule
\end{longtable}
\endgroup

%======================================================================
\section{Notes and references}\label{l23:sec:notes}
%======================================================================

The universal affine vertex algebra and its Sugawara vector: I.~B.~Frenkel and Y.~Zhu, \emph{Vertex operator algebras associated to
representations of affine and Virasoro algebras}, Duke Math.\ J.\ 66 (1992); V.~G.~Kac, \emph{Vertex algebras for beginners}, 2nd ed., AMS
University Lecture Series 10 (1998), \S\S4.7, 5.6; T.~Gannon, \emph{Moonshine beyond the Monster}, Cambridge (2006), \S\S3.2, 5.2 (the
library's conventions). Lattice realizations of simply-laced algebras: I.~B.~Frenkel and V.~G.~Kac, \emph{Basic representations of affine Lie
algebras and dual resonance models}, Invent.\ Math.\ 62 (1980); G.~Segal, Commun.\ Math.\ Phys.\ 80 (1981). Simple Lie algebras, root systems,
Casimir and dual Coxeter numbers: J.~E.~Humphreys, \emph{Introduction to Lie Algebras and Representation Theory}, Springer GTM 9 (1972);
V.~G.~Kac, \emph{Infinite-dimensional Lie algebras}, 3rd ed., Cambridge (1990). Short-vector enumeration: U.~Fincke and M.~Pohst, Math.\
Comp.\ 44 (1985). Theta series of $E_8$: J.-P.~Serre, \emph{A Course in Arithmetic}, Springer GTM 7 (1973), Ch.~VII. The Deligne series:
P.~Deligne, \emph{La s\'erie exceptionnelle de groupes de Lie}, C.~R.\ Acad.\ Sci.\ Paris 322 (1996); T.~Arakawa and A.~Moreau, \emph{Joseph
ideals and lisse minimal W-algebras}, J.\ Inst.\ Math.\ Jussieu 17 (2018), arXiv:1506.00710; C.~Beem, M.~Lemos, P.~Liendo, W.~Peelaers,
L.~Rastelli and B.~C.~van Rees, \emph{Infinite chiral symmetry in four dimensions}, Commun.\ Math.\ Phys.\ 336 (2015), arXiv:1312.5344.
Project records: \leanpath{Notes/\allowbreak deligne/\allowbreak campaign/\allowbreak WS-AFF.md},
\leanpath{Notes/\allowbreak deligne/\allowbreak campaign/\allowbreak LEAN\_CAMPAIGN.md},
\leanpath{Notes/\allowbreak deligne/\allowbreak audit\_2026\_10\_25/\allowbreak WEIGHT\_ONE\_LIE\_NOTE.md}, and the project's internal
records of 26--27 October 2026.

\endgroup
